% =====================================================================
%  Bachelor's Thesis -- From the Risk-Neutral Measure to Girsanov's
%  Theorem in Stochastic Integration: Pricing of Financial Products
%  (English version)
%  Layout: A4, 11 pt, 1.5 line spacing, 2.5 cm margins.
% =====================================================================

\documentclass[11pt,a4paper,oneside]{report}
\usepackage{graphicx}  % University logo
\usepackage{xcolor}
% --------- Encoding and language ------------------------------------
\usepackage[utf8]{inputenc}
\usepackage[T1]{fontenc}
\usepackage[english]{babel}


% --------- Layout ---------------------------------------------------
\usepackage[a4paper,top=2.5cm,bottom=2.5cm,left=2.5cm,right=2.5cm]{geometry}
\usepackage{setspace}
\onehalfspacing                        % 1.5 line spacing

% --------- Mathematics ----------------------------------------------
\usepackage{amsmath,amssymb,amsthm,mathtools}
\usepackage{bm}
% Indicator functions 1_A are typeset with a custom command (see below).

% --------- Other packages -------------------------------------------
\usepackage{graphicx}
\usepackage{tikz}
\usetikzlibrary{arrows.meta,positioning,calc}
\usepackage{enumitem}
\usepackage{booktabs}
\usepackage{caption}
\usepackage[colorlinks=true,linkcolor=black,citecolor=blue,urlcolor=blue]{hyperref}
\hypersetup{
  pdftitle={From the Risk-Neutral Measure to Girsanov's Theorem in Stochastic Integration: Pricing of Financial Products},
  pdfauthor={Roberto Cepeda Rocandio},
  pdfsubject={Bachelor's thesis in Mathematics, University of Oviedo},
  pdfkeywords={no-arbitrage, martingale, Itô integral, Girsanov's theorem, martingale representation, Black-Scholes formula},
  pdflang={en-GB}
}

% --------- Theorem-like environments --------------------------------
% Environment names are shared with the Spanish version so that labels
% and cross-references coincide in both.
\theoremstyle{plain}
\newtheorem{teorema}{Theorem}[chapter]
\newtheorem{proposicion}[teorema]{Proposition}
\newtheorem{lema}[teorema]{Lemma}
\newtheorem{corolario}[teorema]{Corollary}

\theoremstyle{definition}
\newtheorem{definicion}[teorema]{Definition}
\newtheorem{ejemplo}[teorema]{Example}

\theoremstyle{remark}
\newtheorem{observacion}[teorema]{Remark}

% --------- Notation shortcuts ---------------------------------------
\newcommand{\R}{\mathbb{R}}
\newcommand{\N}{\mathbb{N}}
\newcommand{\Z}{\mathbb{Z}}
\newcommand{\E}{\mathbb{E}}
\newcommand{\PP}{\mathbb{P}}
\newcommand{\Q}{\mathbb{Q}}
\newcommand{\F}{\mathcal{F}}
\newcommand{\G}{\mathcal{G}}
\newcommand{\Hh}{\mathcal{H}}
\newcommand{\HH}{\mathcal{H}}
\newcommand{\1}{\mathbf{1}}
\newcommand{\dd}{\mathrm{d}}
\newcommand{\eqdef}{\coloneqq}

% =====================================================================
\begin{document}

% =====================================================================
%   TITLE PAGE
% =====================================================================
\begin{titlepage}
\centering

\vspace*{1.5cm}

% University of Oviedo logo
\includegraphics[width=0.73\textwidth]{logo_uniovi.png}

\vspace{3cm}

% Thesis title
{\LARGE\bfseries From the Risk-Neutral Measure\\[0.3cm]
        to Girsanov's Theorem in Stochastic Integration:\\[0.3cm]
        Pricing of Financial Products\par}

\vspace{1.5cm}

% Author
{\Large Roberto Cepeda Rocandio\par}

\vspace{2.5cm}

% Supervisor
{\large Supervised by: José Santos Domínguez Menchero\par}

\vspace{1.2cm}

% Institutional block
{\sc University of Oviedo}\\[0.2cm]
{Faculty of Science}\\[0.2cm]
{Bachelor's Degree in Mathematics}

\vspace{1.2cm}

% Date
{Academic year 2026--2027}

\vfill
\end{titlepage}

\clearpage

% Roman page numbering in the front matter
\pagenumbering{roman}
% =====================================================================
%   ABSTRACT
% =====================================================================
\chapter*{Abstract}
\addcontentsline{toc}{chapter}{Abstract}
\vspace*{-0.75cm}

\itshape

This work develops, in a self-contained and mathematically rigorous
manner, the theory of European option pricing, culminating in the
proof of the Black-Scholes formula for the European \emph{call}
option. The approach adopted is carried out in a unified way: we
first formulate the problem in the discrete setting of the
multi-period binomial model, where the central conceptual ideas
(no-arbitrage principle, replication, risk-neutral measure,
martingale representation) are accessible with elementary probability
tools; we then translate the problem and its resolution to the
continuous setting of the Black-Scholes model, developing along the
way the necessary apparatus of stochastic calculus. The chosen
approach, via martingales rather than the classical route of partial
differential equations, makes transparent the translation of the
ideas from the discrete case to the continuous case, and allows the
Black-Scholes formula to be presented as the natural limit of the
multi-period binomial formula when the time steps become infinitely
small. Thus, the development requires introducing and constructing in
detail a substantial part of the stochastic calculus needed for the
continuous setting: the Itô integral, Girsanov's theorem and the
martingale representation theorem. These tools are not only necessary
for the purposes of this work, but also play a central role in the
theory of stochastic processes. The dissertation also includes the
derivation of the classical Black-Scholes PDE from the formula
obtained via martingales, thus showing the consistency between the two
formulations.

\upshape

\vspace{0.3cm}
\noindent\textbf{Keywords:} no-arbitrage, martingale, Itô integral,
Girsanov's theorem, martingale representation, Black-Scholes formula.

% =====================================================================
%   TABLE OF CONTENTS
% =====================================================================
\setcounter{tocdepth}{2}    % Contents: chapters + sections + subsections
\tableofcontents

% =====================================================================
%   NOTATION
% =====================================================================
\chapter*{Notation}
\addcontentsline{toc}{chapter}{Notation}

Throughout this work the following notation will be used.

\vspace{0.3cm}

\noindent\textbf{Probability spaces and filtrations}
\begin{center}
\begin{tabular}{@{}p{4cm}p{10cm}@{}}
$(\Omega, \F, \PP)$ & Probability space. \\
$(\F_t)_{t \geq 0}$ & Filtration: encodes the flow of information. \\
$\F_t^W$ & Natural filtration of the Brownian motion $\{W_t\}$. \\
$\PP, \Q$ & Market measure and risk-neutral measure. \\
$\Q \sim \PP$ & The measures $\Q$ and $\PP$ are equivalent. \\
$d\Q/d\PP$ & Radon-Nikodym derivative of $\Q$ with respect to $\PP$. \\
\end{tabular}
\end{center}

\vspace{0.2cm}

\noindent\textbf{Expectations and indicators}
\begin{center}
\begin{tabular}{@{}p{4cm}p{10cm}@{}}
$\E_\PP[X]$ & Expectation of $X$ under $\PP$. \\
$\E_\Q[X]$ & Expectation of $X$ under $\Q$. \\
$\E[X \mid \G]$ & Conditional expectation of $X$ given $\G$. \\
$\1_A$ & Indicator function of the event $A$. \\
\end{tabular}
\end{center}

\vspace{0.2cm}

\noindent\textbf{Stochastic processes}
\begin{center}
\begin{tabular}{@{}p{4cm}p{10cm}@{}}
$\{W_t\}_{t \geq 0}$ & Standard Brownian motion under $\PP$. \\
$\{\widetilde W_t\}$ & Brownian motion under $\Q$ (Girsanov). \\
$\{S_t\}$ & Price of the underlying asset. \\
$\widetilde S_t = e^{-rt} S_t$ & Discounted price of the underlying. \\
$B_t = e^{rt}$ & Price of the risk-free bond. \\
$\{V_t\}$ & Portfolio value. \\
$(\phi_t, \psi_t)$ & Holdings of the underlying and of the bond in the portfolio. \\
$\{L_t\}$ & Radon-Nikodym density process. \\
\end{tabular}
\end{center}

\vspace{0.2cm}

\noindent\textbf{Stochastic calculus}
\begin{center}
\begin{tabular}{@{}p{4cm}p{10cm}@{}}
$\int_0^t \phi_s\, dW_s$ & Itô stochastic integral. \\
$\HH^2[0, T]$ & Processes with $\E[\int_0^T \phi_s^2\, ds] < \infty$. \\
$\langle X, Y\rangle_t$ & Quadratic covariation of $X$ and $Y$. \\
$dX_t = a_t dt + b_t dW_t$ & Differential notation for Itô processes. \\
\end{tabular}
\end{center}

\newpage

\noindent\textbf{Parameters of the Black-Scholes model}
\begin{center}
\begin{tabular}{@{}p{4cm}p{10cm}@{}}
$r$ & Constant risk-free interest rate. \\
$\mu, \sigma$ & Expected return and volatility of the underlying. \\
$T,\ K$ & Maturity and \emph{strike} of the option. \\
$\theta = (\mu - r)/\sigma$ & Market price of risk. \\
$N(\cdot)$ & Standard normal distribution function. \\
$d_1, d_2$ & Coefficients of the Black-Scholes formula. \\
$L^2$ & Space of square-integrable random variables. \\
$\eqdef$ & Equality by definition; a.s. stands for ``almost surely''. \\
\end{tabular}
\end{center}

% Beginning of the main text: Arabic page numbering
\pagenumbering{arabic}
\setcounter{page}{1}
% =====================================================================
%   CHAPTER 1 — INTRODUCTION
% =====================================================================
\chapter{Introduction}\label{cap:intro}

\section{Motivation and context}
\label{sec:motivacion}

A \emph{financial option} is a contract that grants its holder the
right, but not the obligation, to buy or sell an underlying asset (a
stock, an index, a currency, a commodity) at a predetermined price on
a future date. The simplest option, the \emph{European call}, grants
the right to buy at a price $K$ (the \emph{strike}) on a date $T$
(the \emph{maturity}). The holder of the option will exercise the
right if and only if, when the date $T$ arrives, the price of the
underlying is above $K$; in that case they will obtain a profit equal
to the difference. The central question, from the mathematical point
of view, is the following: what is the fair price of this contract at
the time it is written?

The question admits a precise mathematical formulation, anchored in
the \emph{no-arbitrage principle}: the price of the contract must be
such that no market participant can obtain a sure profit without
taking on any risk. Under this principle, together with a
mathematical description of the dynamics of the underlying, the price
is uniquely determined. Determining this price, and constructing the
strategy that allows the writer of the option to \emph{replicate} its
value, is the problem around which the present work is organized.

\medskip

\noindent\textbf{Historical note.} The first mathematical formulation
of the problem is due to Louis Bachelier, who in his doctoral thesis
\emph{Théorie de la spéculation} \cite{bachelier} modeled the
dynamics of stock prices by means of Brownian motion. Bachelier's
work, ignored in its time, was rediscovered in the mid-twentieth
century by Paul Samuelson and other authors, who refined the model by
replacing arithmetic Brownian motion with \emph{geometric Brownian
motion}, which ensures the positivity of prices. The problem was
definitively solved in 1973 by Fischer Black, Myron Scholes and
Robert Merton, in two seminal papers \cite{black-scholes,merton}
recognized in 1997 with the Nobel Prize in Economics awarded to
Scholes and Merton (Black had died in 1995). The Black-Scholes
formula, which constitutes the central result of the present work, is
today one of the most influential mathematical developments of the
last half century.

\medskip

\noindent\textbf{Why it is mathematically interesting.} Beyond its
practical relevance, the problem requires a rich mathematical
apparatus: advanced probability (martingales, conditional
expectation), stochastic analysis (Brownian motion, Itô integral and
formula), functional analysis (density in $L^2$) and, in its PDE
version, parabolic equations. The Black-Scholes formula is a meeting
point of several central areas of modern mathematics.

\medskip

\noindent\textbf{A more general framework than it seems.} It is worth
anticipating an idea that runs through the whole work and goes
beyond the financial problem that motivates it. The notion of fair
price will lead us, both in the discrete and in the continuous model,
to replace the market probability measure by an equivalent one (the
\emph{risk-neutral measure}) under which the discounted price of the
underlying becomes a \emph{martingale}, one of the fundamental
concepts of modern probability. The change of measure producing that
martingale rests on two results of a general nature: the
Radon-Nikodym theorem, which guarantees the existence of the density
between the two measures, and Girsanov's theorem, which describes
that change in the setting of Brownian motion. Neither of these tools
is specific to pricing: the principle of replacing a measure by an
equivalent one to transform a difficult problem into a tractable one
is in common use in probability and its applications, for instance
in the variance reduction of Monte Carlo simulations of rare events
by means of the \emph{importance sampling} technique (see, e.g.,
\cite{karatzas-shreve}). In this sense, the pricing problem makes it
possible to introduce mathematical tools with much wider
applications.

\section{Objectives and approach of the work}
\label{sec:objetivos}

The present work sets out to develop, in a self-contained and
mathematically rigorous manner, the theory of European option
pricing, culminating in the proof of the Black-Scholes formula for
the European \emph{call}. The development will be carried out in a
unified way: first presenting the problem in the \emph{discrete}
setting of the multi-period binomial model, where the conceptual
ideas are accessible with elementary probability tools, and then
transferring the problem and its solution to the \emph{continuous}
setting of the Black-Scholes model, by developing the necessary
stochastic calculus.

\medskip

\noindent\textbf{Justification of the discrete-to-continuous plan.}
The chosen order of exposition is pedagogical, not historical:
chronologically, the Black-Scholes formula (1973) predates the
multi-period binomial model (Cox, Ross and Rubinstein, 1979).
However, from the mathematical point of view the binomial model is
structurally simpler: it allows the central ideas of the theory
(change of measure, replication, no-arbitrage) to be articulated with
elementary probability tools, without the apparatus of stochastic
calculus. This fact is what justifies the plan adopted: once the
logic of the problem has been assimilated in the discrete setting,
the passage to the continuous one reduces to a technical translation
of the tools, keeping the conceptual ideas intact.

\medskip

\noindent\textbf{The martingale route.} The continuous Black-Scholes
theory admits two classical routes: martingales (change of measure
and representation theorems) and PDEs (instantaneous hedging). We opt
for the former since it admits a direct translation from the discrete
case. This choice has, moreover, an advantage that goes beyond mere
exposition: the tools it requires, the Itô integral, Girsanov's
theorem and the martingale representation theorem, are mathematical
objects of interest and usefulness in their own right, and they make
up a good part of the technical content of the work. The PDE route is
recovered at the end (Section~\ref{sec:edp_bs}) as a consequence of
the martingale formulation.

\medskip

\noindent\textbf{What the work does not address.} For reasons of
length and of consistency with the chosen plan, the work is
restricted to \emph{European} options on a single underlying, in the
Black-Scholes model with constant parameters. Consequently, several
natural extensions lie outside its scope—American options, models
with stochastic volatility or jumps, numerical methods for
\emph{payoffs} without a closed-form formula, calibration to market
data—and are briefly collected in Chapter~\ref{cap:conclusiones} as
further directions.

\section{Sources and own contribution}
\label{sec:aportacion}

The expository thread of this work essentially follows that of
A.~Etheridge~\cite{etheridge}, from whom we take both the overall
structure (the passage from the discrete binomial model to the
Black-Scholes model via the martingale route) and the order in which
the main ideas are introduced: the motivation of the stochastic
integral from no-arbitrage, the change of measure as the crucial step
of pricing, and the martingale representation theorem. The
introduction to stochastic calculus, which is restricted to Brownian
motion, likewise follows this structure.

As for the technical part, it relies on several introductory texts on
stochastic calculus and probability theory. For Brownian motion
(construction, main properties and augmented filtration) we mainly
follow S.~E.~Shreve~\cite{shreve}. The central results of stochastic
calculus that we use (the construction of the Itô integral and its
isometry, Itô's formula and Girsanov's theorem) are presented
directly for Brownian motion, following the introductory texts on
the subject~\cite{oksendal,shreve} and the aforementioned
Etheridge~\cite{etheridge}. A more general formulation, valid for
continuous square-integrable martingales and semimartingales, can be
found in I.~Karatzas and S.~E.~Shreve~\cite{karatzas-shreve}. The
results we use in the Brownian case are thus obtained as a particular
case, with $\langle W\rangle_t = t$. For measure theory and
martingale theory (conditional expectation, the Radon-Nikodym theorem
and discrete-time martingales) we follow D.~Williams~\cite{williams}.
Finally, the historical references~\cite{bachelier,black-scholes,merton}
are used only to put the problem in context.

Since the results presented in this work are classical and well
known, the contribution lies not so much in obtaining new results as
in their selection, organization and development. Some proofs that
the sources present in summary form, take for granted, or formulate
at a higher level of generality than is needed here have been
developed in more detail and adapted to the problem at hand. This is
the case of the construction of the Itô integral in
Chapter~\ref{cap:itocalculus}, with the isometry proved step by step,
the extension to $\mathcal{H}^2$ and the continuity of the indefinite
integral; of Itô's formula, detailing the analytic tools involved in
its proof; and of the constructive no-arbitrage argument motivating
the introduction of a new integral
(Section~\ref{sec:motivacion_ito}). In addition, we have sought to
keep a uniform notation across the different sources and to relate
the discrete and continuous cases clearly. Thus, Girsanov's change of
measure (Section~\ref{sec:girsanov}), for instance, is presented as
the continuous version of the adjustment of the probabilities of the
different paths of the binomial model, and not as an independent
result.

Some of the aspects that required the most work were, already in the
discrete model, the very idea of \emph{no-arbitrage}
(Chapter~\ref{cap:discreto}). Understanding why the absence of sure
profits without investment makes it possible to determine a price,
and how this principle becomes the foundation of the theory, requires
a change of perspective that is not immediate. In the continuous
model, the difference between the natural Brownian filtration and its
augmentation (Chapter~\ref{cap:brown}) also required special
attention, since the subsequent results needed a filtration
satisfying certain technical properties. Finally, the appearance of
the \emph{Itô term} made it necessary to explain why the quadratic
variation of Brownian motion cannot be neglected and where the
second-order correction $(dW_t)^2 = dt$, which distinguishes
stochastic calculus from ordinary calculus, comes from.

\section{Structure of the dissertation}
\label{sec:estructura}

The dissertation is organized in six chapters: introduction, the
discrete part (Chapter~\ref{cap:discreto}), the continuous
foundations (Chapters~\ref{cap:brown} and~\ref{cap:itocalculus}), the
application to Black-Scholes (Chapter~\ref{cap:bs}) and conclusions.

\medskip

\noindent\textbf{Chapter~\ref{cap:discreto}: the binomial model.} It
poses the pricing problem in the discrete setting of the multi-period
binomial model. It sets out the elementary probability tools
required (filtrations, conditional expectation, martingales, discrete
stochastic integral) and solves the problem by means of a three-step
recipe: change of measure, construction of the martingale associated
with the discounted \emph{payoff}, and application of the binomial
representation theorem to identify the replicating portfolio. The
chapter concludes with a brief introduction to the continuous
setting, anticipating the convergence of the binomial random walk to
Brownian motion.

\medskip

\noindent\textbf{Chapter~\ref{cap:brown}: Brownian motion.} It
introduces Brownian motion as a mathematical object, establishes its
existence as a theorem (without proof, citing the literature) and
studies the properties of its paths: continuity, nowhere
differentiability, infinite total variation and finite quadratic
variation equal to the length of the interval. This last property is
the technical piece that will motivate the development of the
stochastic calculus of the following chapter.

\medskip

\noindent\textbf{Chapter~\ref{cap:itocalculus}: stochastic calculus.}
It constitutes the technical core of the work. It constructs the Itô
stochastic integral first on simple processes and then, by density,
on the space $\mathcal{H}^2$. It proves Itô's formula for functions of
Brownian motion and states its general version. It formulates the
change of measure in continuous time by means of Girsanov's theorem,
and states the Brownian martingale representation theorem. Each of
these pieces is the translation to the continuous setting of an
object introduced in the discrete case in Chapter~\ref{cap:discreto}.

\medskip

\noindent\textbf{Chapter~\ref{cap:bs}: the Black-Scholes model.} It
applies the tools of Chapter~\ref{cap:itocalculus} to the pricing
problem in the Black-Scholes model, carrying out the three-step
recipe in its continuous version. It obtains the general formula for
the price of a European derivative and specializes it to the
\emph{payoff} of the European \emph{call}, explicitly deriving the
classical Black-Scholes formula with its coefficients $d_1$ and
$d_2$. The chapter concludes by deriving the Black-Scholes PDE from
the formula obtained, showing the consistency with the classical
route.

\medskip

\noindent\textbf{Chapter~\ref{cap:conclusiones}: conclusions.} It
summarizes the work carried out, presents the convergence of the
multi-period binomial formula to the Black-Scholes formula (as the
culmination of the relationship between the two models in
quantitative terms) and briefly discusses the natural further
directions of the work.
% =====================================================================
%   CHAPTER 2 — OPTION PRICING IN DISCRETE TIME
% =====================================================================

\chapter{Option pricing in discrete time}\label{cap:discreto}

\section{Financial concepts and the no-arbitrage principle}
\label{sec:conceptos_no_arbitraje}

\subsection{Basic financial definitions}
\label{subsec:def_basicas}

Throughout this work we place ourselves in an idealized market in which
two basic types of financial instruments are traded: a risk-free asset
and a risky asset. Derivative contracts will be built on them; their
value ultimately depends on the future behavior of the underlying
asset.

\medskip

\noindent\textbf{The risk-free bond.} We denote by $B_t$ the price at
time $t \geq 0$ of the risk-free asset, normalized so that $B_0 = 1$.
This asset represents an investment without uncertainty (in practice,
an interest-bearing deposit or a government bond) whose value grows
deterministically at a fixed interest rate. Before writing down its
expression, it is worth pausing to see how this growth is modeled.

Suppose one monetary unit is invested at an interest rate $r$ per unit
of time and that interest is reinvested (\emph{compounded}) $n$ times
per unit of time. Each subperiod, of length $1/n$, accrues the
proportional rate $r/n$; since interest is reinvested, the capital is
multiplied in each subperiod by the factor $1 + r/n$, so that after one
unit of time ($n$ subperiods) the initial monetary unit has become
\begin{equation*}
  \left(1 + \frac{r}{n}\right)^{n}
\end{equation*}
This is \emph{compound interest} with frequency $n$.

\emph{Continuous compounding} arises when interest is reinvested
instantaneously, that is, when the frequency $n$ is made arbitrarily
large. Taking the limit yields the classical expression
\begin{equation}\label{eq:limite_e}
  \lim_{n \to \infty}\left(1 + \frac{r}{n}\right)^{n} = e^{r}
\end{equation}
The growth factor per unit of time is therefore $e^{r}$, and the one
corresponding to an arbitrary horizon $t > 0$ is obtained by
composition:
\begin{equation*}
  B_t = \left(e^{r}\right)^{t} = e^{rt}.
\end{equation*}
We shall therefore adopt
\begin{equation}\label{eq:bono}
  B_t = e^{rt}, \qquad r > 0, \qquad B_0 = 1,
\end{equation}
as our model for the risk-free asset. The name \emph{risk-free} is
justified by the fact that the return is known with certainty at the
initial time: investing one monetary unit at $t = 0$ guarantees
$e^{rT}$ units at $t = T$, with no uncertainty about the outcome.

The choice of continuous compounding entails no loss of generality
with respect to the usual discrete regimes: it is merely a convenient
way of expressing the same phenomenon. Indeed, if a discrete regime
offers an effective rate $r'$ per unit of time, so that one monetary
unit becomes $1 + r'$ after one unit of time, the continuous rate $r$
producing the same growth factor is the one satisfying
$e^{r} = 1 + r'$, that is,
\begin{equation*}
  r = \log(1 + r').
\end{equation*}
For instance, an effective annual rate of $5\%$ ($r' = 0.05$) is
equivalent to a continuous rate $r = \log(1.05) \approx 0.0488$, that
is, $4.88\%$. From now on, $r$ will always denote the continuously
compounded rate.

\medskip

\noindent\textbf{The risky stock.} By contrast, the price of the risky
asset (which we shall generically call the \emph{stock}) is modeled as
a stochastic process $\{S_t\}_{t \geq 0}$ defined on a suitable
probability space. The specific dynamics of $\{S_t\}$ will be made
precise as the work progresses: in this chapter it will take values in
a finite set at each time (binomial model), whereas in
Chapters~\ref{cap:brown} and \ref{cap:bs} we shall consider
continuous-time models.

\medskip

\noindent\textbf{Derivative contracts.} A \emph{derivative contract}
is an agreement between two parties whose \emph{payoff} (the amount
they must exchange at a future date) depends on the price of the
underlying asset. Three derivatives will appear repeatedly throughout
this work; they are introduced in the following definitions, where
$T > 0$ denotes the maturity date and $K > 0$ the price agreed in the
contract.

\begin{definicion}[Forward contract]
A \emph{forward} contract on the stock $S$ with maturity $T$ and
delivery price $K$ is an agreement whereby the buyer is obliged to
purchase the asset $S$ at time $T$ at the price $K$, fixed in advance.
The \emph{payoff} to the buyer at $T$ is the random variable
\begin{equation*}
  S_T - K.
\end{equation*}
\end{definicion}

The \emph{forward} allows both parties to fix today the price of a
transaction that will be executed in the future, thereby removing the
uncertainty about the price at which the asset will be bought or sold
at $T$. It is, for this reason, a common hedging tool: whoever will
need the asset later is protected against a rise in its price, and
whoever intends to dispose of it, against a fall.

The \emph{forward} is a contract \emph{of obligation}: both parties
must honor the agreement, so the \emph{payoff} may be negative if
$S_T < K$. The natural question is which value of $K$ makes the
contract \emph{fair}, in the sense that it can be entered into at no
cost at $t = 0$.

\medskip

Before defining them, it is worth explaining what the two options that
follow are for. Unlike the \emph{forward}, which binds both parties,
an option works in practice as \emph{insurance} against movements in
the price of the underlying. A \emph{call} is of interest to someone
who expects to buy the asset in the future and wants protection
against a rise: it allows them to lock in a purchase price today,
benefiting from the rise if it occurs, but without being bound if the
price falls. A \emph{put}, on the other hand, is of interest to someone
who owns the asset or plans to sell it and fears a fall: it guarantees
a selling price, protecting their gains against declines. Besides this
hedging use, options are also employed for speculative purposes, since
they allow one to bet on the direction of the price with a small
initial outlay.

\begin{definicion}[European \emph{call} option]
A European \emph{call option} on $S$ with maturity $T$ and strike
price $K$ gives the buyer the right (but not the obligation) to
purchase the asset at $T$ at the price $K$. Its \emph{payoff} at $T$
is the random variable
\begin{equation*}
  (S_T - K)^+ \eqdef \max\{S_T - K,\, 0\}.
\end{equation*}
\end{definicion}

\begin{definicion}[European \emph{put} option]
A European \emph{put option} on $S$ with maturity $T$ and strike price
$K$ gives the buyer the right (but not the obligation) to sell the
asset at $T$ at the price $K$. Its \emph{payoff} at $T$ is the random
variable
\begin{equation*}
  (K - S_T)^+ \eqdef \max\{K - S_T,\, 0\}.
\end{equation*}
\end{definicion}

Unlike the \emph{forward}, the \emph{call} and the \emph{put} are
contracts \emph{of right}: the holder decides at maturity whether or
not to exercise, choosing whichever is more favorable. Their
\emph{payoff} is therefore nonnegative and, unlike the \emph{forward},
whose fair price at $t = 0$ turns out to be zero, an option always has
strictly positive value when it is written.

\medskip

\noindent\textbf{Vocabulary.} In standard terminology, the asset $S$ is
called the \emph{underlying}, the price $K$ the \emph{strike}, and $T$
the \emph{maturity}. The qualifier \emph{European} indicates that
exercise is only possible at $T$, in contrast with American options
(exercisable at any $t \in [0, T]$), which lie outside the scope of
this work.

\subsection{The no-arbitrage principle}
\label{subsec:noarbitraje}

Before turning to the specific model that will occupy the rest of the
chapter, it is worth pausing on the principle underlying the whole
pricing theory. This principle, known as \emph{absence of arbitrage}
or the \emph{no-arbitrage principle}, is not a mathematical result we
are going to prove, but an assumption about the market that we take
as our starting point.

\medskip

\noindent\textbf{The informal idea.} An \emph{arbitrage opportunity}
is, intuitively, an investment strategy which, starting from zero
initial capital, guarantees the investor a nonnegative profit in every
possible scenario and a strictly positive one in at least one of them.
In other words: a way of obtaining something from nothing, without
taking on any risk.

We shall assume that the market admits no arbitrage opportunities. The
justification is economic: any real opportunity would be detected and
exploited immediately, causing it to disappear; arbitrages, if they
exist, are transient phenomena. The mathematical formalization is
deferred to Section~\ref{sec:binario}, once the model has been
specified.

\medskip

\noindent\textbf{The replication principle.} The no-arbitrage
assumption can be reformulated in a more operational way. The
following statement, although informal, captures its fundamental idea.

\begin{quote}
\emph{If two portfolios of assets produce exactly the same future
payments in all possible scenarios, then they must have the same value
at any earlier time.}
\end{quote}

The justification is immediate: if at some time $t < T$ one of the
portfolios were worth more than the other, selling the expensive one
and buying the cheap one would produce a sure gain without taking on
any risk, contradicting the assumption. We shall call this consequence
the \emph{replication principle}.

This principle suggests a general strategy for pricing derivative
contracts: given an option or, more generally, any contract whose
\emph{payoff} at $T$ is a known random variable $C_T$, it suffices to
construct a portfolio (made up of assets whose prices we do know)
whose value at $T$ coincides with $C_T$ in every scenario. We then say
that the portfolio \emph{replicates} the contract. By the replication
principle, the fair price of the derivative at any earlier time is
necessarily equal to the value of this replicating portfolio. This is
the central idea of the whole theory we shall develop: both in the
discrete and in the continuous setting, the problem of finding the
price of a derivative will ultimately reduce to the problem of
constructing a replicating portfolio.

\medskip

\noindent\textbf{On the role of probabilities.} It is worth making a
remark that will gain significance over the course of the chapter:
the replication argument, as we have formulated it, never makes use of
the probabilities attached to each future scenario. Only the
\emph{possible} values of the asset come into play, not the
probability of their occurrence. As we shall see, this fact has deep
consequences: the fair price of a derivative is independent of the
market's subjective beliefs about the future. It is one of the most
surprising and beautiful ideas in the theory.

\subsection{Pricing the forward as a warm-up case}
\label{subsec:forward}

As a first application of the replication principle, we address the
simplest possible problem: the pricing of a \emph{forward}. Throughout
this subsection we assume that the two assets in the market can be
traded without restrictions at any time $t \in [0, T]$, with no
transaction costs and in arbitrary fractional amounts. These
assumptions are idealizations, but they will suffice for the arguments
that follow to hold.

\medskip

\noindent\textbf{The fair price of a \emph{forward} with \emph{strike}
$K$.} Consider a \emph{forward} on the stock $S$ with maturity $T$ and
delivery price $K$. Its \emph{payoff} to the buyer at $T$ is
$S_T - K$. We seek the fair price $V_0$ of the contract at the initial
time: the amount the buyer must pay the writer at $t = 0$ in exchange
for receiving $S_T - K$ at $t = T$.

Consider the following portfolio, set up at $t = 0$:
\begin{itemize}
  \item one unit of the stock is bought, at the price $S_0$;
  \item an amount $K e^{-rT}$ of the risk-free asset is borrowed (that
        is, issued).
\end{itemize}
The net cost of setting it up at $t = 0$ is $\Pi_0 = S_0 - K e^{-rT}$.
At $t = T$, the unit of stock will be worth $S_T$, while the debt
incurred will have grown to $K e^{-rT} \cdot e^{rT} = K$. Consequently,
the value of the portfolio at $t = T$ is
\begin{equation*}
  \Pi_T = S_T - K,
\end{equation*}
which coincides with the \emph{payoff} of the \emph{forward}. By the
replication principle,
\begin{equation}\label{eq:precio_forward}
  V_0 = S_0 - K e^{-rT}.
\end{equation}

It is worth checking that the argument is consistent, in the sense
that any deviation from the price \eqref{eq:precio_forward} would
produce an arbitrage opportunity. If the contract were sold at
$V_0' > S_0 - K e^{-rT}$, an investor could: at $t = 0$, sell the
\emph{forward} (receiving $V_0'$), set up the replicating portfolio
(at a cost $S_0 - K e^{-rT}$) and keep the remainder
$V_0' - (S_0 - K e^{-rT}) > 0$. At $t = T$, the payments associated
with the \emph{forward} and with the portfolio cancel out, so the
investor would have obtained a sure profit without taking on any risk.
The symmetric situation $V_0' < S_0 - K e^{-rT}$ leads to an analogous
arbitrage.

\medskip

\noindent\textbf{The fair strike and the \emph{forward} price.} In
practice, \emph{forwards} are usually entered into with $V_0 = 0$.
Imposing this condition in \eqref{eq:precio_forward} allows us to solve
for the delivery price that makes the contract fair:
\begin{equation}\label{eq:precio_forward_just}
  K = S_0 e^{rT}.
\end{equation}
This quantity is called the \emph{forward price} of the stock for the
maturity $T$, and it is interpreted as the value at $t = T$ that one
unit of stock bought today would have if compounded at the risk-free
rate.

\medskip

\noindent\textbf{Two general lessons.}

\begin{observacion}[Independence from the dynamics of the underlying]
The derivation of the \emph{forward} price has made no use whatsoever
of the specific model governing the evolution of the stock price: we
have assumed neither that $\{S_t\}$ follows a particular process, nor
even that it takes values in a specific set. Formula
\eqref{eq:precio_forward} is valid in any arbitrage-free market in
which the bond and the stock can be traded without restrictions. The
ultimate reason is that the \emph{payoff} of the \emph{forward}
depends \emph{linearly} on the underlying, and therefore replication
can be carried out by a static strategy. This universality is,
however, an exception: as soon as the \emph{payoff} of the derivative
is nonlinear in $S_T$ (as is the case for the \emph{call} and the
\emph{put}), the replication argument ceases to be static and comes
to require a specific model for the dynamics of the underlying.
\end{observacion}

\begin{observacion}[Irrelevance of subjective probabilities]
No probability associated with the future evolution of the stock
appears in formula \eqref{eq:precio_forward}. The fair price of the
\emph{forward} does not depend on how likely it is, according to the
market's beliefs, that the stock will go up or down, nor on its
expected mean return, nor on its volatility. This independence from
subjective probabilities will reappear, in a subtler form, in the
nontrivial models we study next.
\end{observacion}

Taking this initial example as a reference, we can now address the
general problem in its simplest, though by no means trivial, version.
\section{The one-period binary model}
\label{sec:binario}

The simplest model that retains the essential ingredients of the
problem is the \emph{one-period binary} model: the market is observed
only at $t = 0$ and $t = T$, and the price of the underlying takes only
two possible values at $T$.

\subsection{Description of the model and formalization of no-arbitrage}
\label{subsec:setup_binario}

\noindent\textbf{Description of the model.} The market consists of a
risk-free bond, with $B_t = e^{rt}$ and $r > 0$, and a stock whose
initial price $S_0 > 0$ is known and whose final price $S_T$ is a
random variable taking only two values. Specifically, we fix two
constants $u, d$ with $0 < d < u$ and assume that
\begin{equation*}
  S_T = \begin{cases}
    S_0\, u & \text{with probability } p, \\[2pt]
    S_0\, d & \text{with probability } 1 - p,
  \end{cases}
\end{equation*}
for some $p \in (0, 1)$. We call $u$ and $d$ the \emph{up} and
\emph{down factors}, and $p$ the \emph{subjective probability} of an
up move.

To fix ideas, we formulate the model in the language of probability
theory. Let $\Omega = \{\omega_u, \omega_d\}$ be the sample space,
$\F = 2^{\Omega}$, and let $\PP$ be the probability on $(\Omega, \F)$
defined by $\PP(\{\omega_u\}) = p$ and $\PP(\{\omega_d\}) = 1 - p$. On
$(\Omega, \F, \PP)$, the final price of the stock is the random
variable
\begin{equation*}
  S_T(\omega_u) = S_0\, u, \qquad S_T(\omega_d) = S_0\, d.
\end{equation*}
We call $\PP$ the \emph{market measure} or \emph{physical measure}.

Although we include $p$ in the description, its only role is to
specify which scenarios we regard as \emph{possible} (those with
strictly positive probability). The numerical value of $p$ will play
no part in any computation.

\medskip

\noindent\textbf{Formalization of the notion of arbitrage.}

\begin{definicion}[Strategy or portfolio]
A \emph{strategy} in the one-period binary model is a pair
$(\phi, \psi) \in \R^2$, where $\phi$ represents the number of units of
the stock and $\psi$ the number of units of the bond held in the
portfolio during the period $[0, T]$.
\end{definicion}

The value of the strategy at each time is
\begin{align*}
  V_0(\phi, \psi)   &= \phi\, S_0 + \psi, \\
  V_T(\phi, \psi)   &= \phi\, S_T + \psi\, e^{rT}.
\end{align*}
We allow $\phi$ and $\psi$ to take any real value, including negative
ones: $\phi < 0$ corresponds to \emph{going short} in the stock, and
$\psi < 0$ to borrowing an amount $|\psi|$ from the bank. This
possibility is essential for the theory: without it, not every option
could be replicated.

\begin{definicion}[Arbitrage opportunity]
\label{def:arbitraje}
A strategy $(\phi, \psi)$ is an \emph{arbitrage opportunity} if:
\begin{enumerate}[label=\textnormal{(\roman*)}]
  \item $V_0(\phi, \psi) = 0$;
  \item $V_T(\phi, \psi) \geq 0$ almost surely under $\PP$;
  \item $\PP\bigl(\{V_T(\phi, \psi) > 0\}\bigr) > 0$.
\end{enumerate}
We say that the model is \emph{arbitrage-free} if it admits no
strategy of this kind.
\end{definicion}

\medskip

\begin{teorema}[Characterization of no-arbitrage in the binary model]
\label{teo:noarb_binario}
The one-period binary model is arbitrage-free if and only if
\begin{equation}\label{eq:cond_noarb_binario}
  d < e^{rT} < u.
\end{equation}
\end{teorema}

\begin{proof}
$(\Rightarrow)$ We argue by contraposition: if
condition~\eqref{eq:cond_noarb_binario} fails, we exhibit an
arbitrage opportunity.

Suppose $e^{rT} \leq d$. Consider the strategy
$(\phi, \psi) = (1, -S_0)$, which consists in buying one unit of stock
financed with the bond. Its initial value is $V_0 = S_0 - S_0 = 0$,
and its final value is
\begin{equation*}
  V_T = S_T - S_0\, e^{rT} = S_0\bigl(S_T/S_0 - e^{rT}\bigr),
\end{equation*}
which equals $S_0(u - e^{rT}) > 0$ at $\omega_u$ and
$S_0(d - e^{rT}) \geq 0$ at $\omega_d$. The strategy is an arbitrage
opportunity. Similarly, if $e^{rT} \geq u$, the strategy
$(\phi, \psi) = (-1, S_0)$ yields an arbitrage.

$(\Leftarrow)$ Assume~\eqref{eq:cond_noarb_binario}. Let
$(\phi, \psi)$ satisfy $V_0 = 0$, that is, $\psi = -\phi\, S_0$. Its
final value is
\begin{align*}
  V_T(\omega_u) &= \phi\, S_0\bigl(u - e^{rT}\bigr), \\
  V_T(\omega_d) &= \phi\, S_0\bigl(d - e^{rT}\bigr).
\end{align*}
By~\eqref{eq:cond_noarb_binario}, the factors $u - e^{rT}$ and
$d - e^{rT}$ have opposite signs. The sign of $V_T$ is therefore
opposite in the two states, except in the trivial case $\phi = 0$,
which implies $V_T \equiv 0$. There is no strategy
$(\phi, \psi) \neq (0, 0)$ with $V_0 = 0$, $V_T \geq 0$ and $V_T > 0$
in at least one state.
\end{proof}

Condition~\eqref{eq:cond_noarb_binario} has a natural economic
interpretation: the return on the stock must be able to exceed that of
the risk-free bond (hence $u > e^{rT}$) but must also be able to fall
below it (hence $d < e^{rT}$). Otherwise trivial arbitrages would
arise. The assumption requires the stock to be genuinely \emph{risky}
relative to the bond. We assume it from now on.

\subsection{Construction of the replicating portfolio}
\label{subsec:replicante_binario}

Consider a derivative with \emph{payoff} $C_T$ defined on
$(\Omega, \F, \PP)$. By construction, $C_T$ takes at most two values:
\begin{equation*}
  C_u \eqdef C_T(\omega_u), \qquad C_d \eqdef C_T(\omega_d).
\end{equation*}
We seek a strategy $(\phi, \psi) \in \R^2$ such that
\begin{equation}\label{eq:replicante_binario}
  V_T(\phi, \psi)(\omega) = C_T(\omega) \quad \text{for all } \omega \in \Omega.
\end{equation}
A strategy satisfying \eqref{eq:replicante_binario} is called a
\emph{replicating portfolio} for the derivative.

Condition~\eqref{eq:replicante_binario}, written out explicitly in
each state, becomes the linear system
\begin{equation}\label{eq:sistema_replicante}
  \left\{
  \begin{aligned}
    \phi\, S_0\, u + \psi\, e^{rT} &= C_u, \\
    \phi\, S_0\, d + \psi\, e^{rT} &= C_d.
  \end{aligned}
  \right.
\end{equation}
The coefficient matrix has determinant $S_0\, e^{rT}\, (u - d) > 0$,
and therefore the system has a unique solution:
\begin{equation}\label{eq:phi_psi}
  \phi = \frac{C_u - C_d}{S_0\,(u - d)}, \qquad
  \psi = e^{-rT}\, \frac{u\, C_d - d\, C_u}{u - d}.
\end{equation}

\begin{teorema}[Replication in the one-period binary model]
\label{teo:replicacion_binario}
In the arbitrage-free one-period binary model, every derivative with
\emph{payoff} $C_T$ admits a unique replicating portfolio
$(\phi, \psi) \in \R^2$, given by~\eqref{eq:phi_psi}. The fair price
of the derivative at $t = 0$ is
\begin{equation}\label{eq:precio_binario}
  V_0 = \phi\, S_0 + \psi
       = \frac{C_u - C_d}{u - d}
       + e^{-rT}\, \frac{u\, C_d - d\, C_u}{u - d}.
\end{equation}
\end{teorema}

\begin{proof}
The existence and uniqueness of $(\phi, \psi)$ were established in the
preceding discussion. To see that $V_0 = \phi\, S_0 + \psi$ is the
fair price, suppose the derivative trades at $V_0' > V_0$. Selling the
derivative and setting up the replicating portfolio requires an
initial investment $V_0 - V_0' < 0$, leaving a positive remainder. At
$t = T$, the \emph{payoff} of the derivative and the value of the
portfolio cancel out, producing a sure profit. Similarly for
$V_0' < V_0$. Both situations contradict the no-arbitrage assumption.
\end{proof}

\medskip

\noindent\textbf{Two remarks.} Formula \eqref{eq:precio_binario} does
not contain the subjective probability $p$: the price depends only on
$S_0, u, d, r, T, C_u, C_d$. On the other hand, the quantity $\phi$
obtained in~\eqref{eq:phi_psi} is the ratio between the change in the
\emph{payoff} and the change in the price of the underlying: it
measures the sensitivity of the derivative to movements of the
underlying and is known in the financial literature as the
\emph{delta}.

\subsection{Emergence of the risk-neutral measure}
\label{subsec:medida_riesgo}

Formula~\eqref{eq:precio_binario} admits an elementary algebraic
manipulation that brings out a new (and central) object of the
theory. Regrouping the terms in $C_u$ and $C_d$,
\begin{equation*}
  V_0 = e^{-rT}\!\left[\frac{e^{rT} - d}{u - d}\, C_u
                       + \frac{u - e^{rT}}{u - d}\, C_d\right].
\end{equation*}
Defining
\begin{equation}\label{eq:def_q}
  q \eqdef \frac{e^{rT} - d}{u - d},
\end{equation}
and noting that $1 - q = (u - e^{rT})/(u - d)$, the formula takes the
compact form
\begin{equation}\label{eq:precio_binario_q}
  V_0 = e^{-rT}\,\bigl[\,q\, C_u + (1 - q)\, C_d\,\bigr].
\end{equation}

\medskip

\noindent\textbf{Properties of the ratio $q$.} The no-arbitrage
condition $d < e^{rT} < u$ implies that both $e^{rT} - d > 0$ and
$u - d > 0$, and that $u - e^{rT} > 0$. Consequently, $q \in (0, 1)$.
The converse statement is also easily verified. This equivalence
between $q \in (0,1)$ and the no-arbitrage
condition~\eqref{eq:cond_noarb_binario} is important: what the fact
$q \in (0,1)$ implies is that the coefficients appearing in
formula~\eqref{eq:precio_binario_q} are nonnegative numbers adding up
to one, that is, \emph{they are the probabilities of a measure on
$\Omega$}.

\medskip

\noindent\textbf{Definition of the risk-neutral measure.} The
preceding remark justifies introducing, on the space $(\Omega, \F)$,
a new probability measure
\begin{equation}\label{eq:def_Q}
  \Q(\{\omega_u\}) = q, \qquad \Q(\{\omega_d\}) = 1 - q.
\end{equation}

\begin{definicion}[Risk-neutral measure]
\label{def:medida_riesgo}
The probability $\Q$ defined in~\eqref{eq:def_Q} is called the
\emph{risk-neutral measure} (or \emph{equivalent martingale measure})
of the one-period binary model.
\end{definicion}

Note that $\Q$ and $\PP$ assign strictly positive probability to the
same events. In the language of measure theory, they are
\emph{equivalent}, written $\Q \sim \PP$. This equivalence, trivial
here, will be a delicate condition in the continuous models of
Chapter~\ref{cap:bs}.

\medskip

\noindent\textbf{The fundamental property of $\Q$.} Let us compute
the expectation (under $\Q$) of the discounted stock price at
$t = T$:
\begin{align*}
  \E_{\Q}\bigl[e^{-rT}\, S_T\bigr]
  &= e^{-rT}\, S_0\,\bigl[q\, u + (1 - q)\, d\bigr].
\end{align*}
Substituting $q = (e^{rT} - d)/(u - d)$, a direct computation shows
that $q\, u + (1 - q)\, d = e^{rT}$, and hence
\begin{equation}\label{eq:martingala_uniperiodo}
  \E_{\Q}\bigl[e^{-rT}\, S_T\bigr] = S_0.
\end{equation}

Identity~\eqref{eq:martingala_uniperiodo} contains a central idea. If
we introduce the notation $\widetilde S_t \eqdef e^{-rt} S_t$ for the
discounted price, it can be rewritten as
$\E_{\Q}[\widetilde S_T] = \widetilde S_0$, a property that in
Section~\ref{subsec:martingalas_discretas} we shall recognize as the
preliminary version of the \emph{martingale} property. This is why
$\Q$ is also called the equivalent martingale measure.

As for the adjective \emph{risk-neutral}, its origin is the following.
Rewriting \eqref{eq:martingala_uniperiodo} without discounting,
$\E_\Q[S_T] = S_0\, e^{rT}$: under $\Q$, the mean return on the stock
coincides with that of the risk-free bond. That is, $\Q$ is the
measure under which an investor would be indifferent between holding
the risky asset and the risk-free asset, since both would offer the
same expected return. In reality (under $\PP$) investors are
risk-averse and demand a premium for holding risky assets; the measure
$\Q$ removes that premium by construction.

\medskip

\noindent\textbf{The pricing formula as a discounted expectation.} By
the definition~\eqref{eq:def_Q} of $\Q$, the expression
$q\, C_u + (1-q)\, C_d$ is exactly the expectation of $C_T$ under
$\Q$. The pricing formula then takes its definitive form
\begin{equation}\label{eq:precio_esperanza}
  \boxed{\,V_0 = e^{-rT}\, \E_{\Q}[C_T].\,}
\end{equation}
The fair price of a derivative at $t = 0$ is the discounted value of
the expectation of its \emph{payoff} under the risk-neutral measure
$\Q$. The formula reduces the pricing problem to a computation of
expectations, once the appropriate measure has been identified.

Let us insist on a remarkable feature of
\eqref{eq:precio_esperanza}: we have arrived at it without ever using
the subjective probability $p$.

\subsection{Independence from subjective probabilities}
\label{subsec:indep_subjetiva}

\noindent\textbf{An intuition worth correcting.} For a reader with a
background in probability, the first natural answer to the question
``what is the fair price of a derivative with random \emph{payoff}
$C_T$?'' would be: \emph{the discounted expected value of its
\emph{payoff}}, that is,
\begin{equation}\label{eq:precio_ingenuo}
  V_0^{\text{naive}} \eqdef e^{-rT}\, \E_{\PP}[C_T]
                  = e^{-rT}\,\bigl[p\, C_u + (1-p)\, C_d\bigr].
\end{equation}
However, it suffices to compare~\eqref{eq:precio_ingenuo} with the
correct price~\eqref{eq:precio_esperanza} to notice that the two
coincide if and only if $p = q$. Except in this fortuitous case, the
two prices differ.

What is wrong with the intuitive reasoning?
Formula~\eqref{eq:precio_ingenuo} would only be reasonable if
investors were indifferent to risk, in which case they would value any
asset by its discounted expected value. In practice, investors are
risk-averse and demand premiums. If \eqref{eq:precio_ingenuo} were
used to set the price of the derivative, arbitrage opportunities would
arise\footnote{If $p > q$, then $V_0^{\text{naive}} > V_0$ (in the
case $C_u > C_d$), and an investor could sell the derivative at the
price $V_0^{\text{naive}}$, set up the replicating portfolio at cost
$V_0$, and pocket the difference without taking on any risk. The
situation $p < q$ produces a symmetric arbitrage.}.

\medskip

\noindent\textbf{Pricing by replication, not by expectation.} The
argument that led us to~\eqref{eq:precio_esperanza} does not rest on
any computation of expectations under the market measure, but on the
explicit construction of a replicating portfolio. What determines the
price of the derivative is the possibility of \emph{reproducing} its
\emph{payoff} by means of a combination of bonds and stocks; the
subjective probability $p$ does not come into play because the
replicating portfolio reproduces the \emph{payoff} in \emph{all
possible scenarios}, with no need to weight them by their
probabilities.

\medskip

\noindent\textbf{The role of probabilities in the theory.} The above
might suggest that probabilities disappear from the pricing problem,
but this is not so: the \emph{relevant} probabilities are not those
of the market, but those of the risk-neutral measure $\Q$. Under
$\Q$, the price of the derivative is recovered as the discounted
expectation of the \emph{payoff}. The initial probabilistic intuition
was not wrong in its form (the price is indeed a discounted
expectation) but in the choice of the measure under which it must be
taken.

This replacement of $\PP$ by $\Q$ is the idea that will structure the
subsequent development, both in the multi-period binomial model and
in Black-Scholes.

\subsection{The trinomial model and incompleteness}
\label{subsec:ternario}

The natural generalization of the binary model consists in allowing
more values for the underlying. We briefly study the trinomial case
to show that replication breaks down, a phenomenon that motivates the
notion of \emph{completeness} of a market.

\medskip

We keep the two assets of the model and assume that at $t = T$ the
price takes values in $\{S_0\, u_1, S_0\, u_2, S_0\, u_3\}$ with
$0 < u_3 < u_2 < u_1$, all with strictly positive probability. We
assume the model is arbitrage-free (which requires
$u_3 < e^{rT} < u_1$).

Given a derivative with \emph{payoff} $C_T$ taking the values
$C_i \eqdef C_T(\omega_i)$ for $i = 1, 2, 3$, the replication
condition $V_T(\phi, \psi) = C_T$ becomes a linear system of
\emph{three} equations in two unknowns:
\begin{equation*}
  \phi\, S_0\, u_i + \psi\, e^{rT} = C_i, \quad i = 1, 2, 3.
\end{equation*}
For an arbitrary \emph{payoff}, the vector $(C_1, C_2, C_3)^\top$ does
not in general belong to the two-dimensional subspace of $\R^3$
spanned by the columns of the system matrix, and therefore the system
has no solution. The conclusion is immediate: in the trinomial model,
the replication method fails, and the no-arbitrage principle, by
itself, does not provide a unique price.

\begin{definicion}[Complete market, preliminary version]
\label{def:completitud_preliminar}
A market is said to be \emph{complete} if every derivative on the
assets of the market admits a replicating portfolio.
\end{definicion}

In an incomplete market, the no-arbitrage principle still constrains
the price of a derivative, but not uniquely: it leaves an
\emph{interval} of arbitrage-free prices, in general nondegenerate. To
fix a price within the interval, additional considerations must be
invoked (investor preferences, equilibrium models, utility criteria)
that go beyond the pure no-arbitrage principle.

Throughout this work we shall only deal with complete models. Both the
multi-period binomial model (Section~\ref{sec:multiperiodo}) and the
Black-Scholes model (Chapter~\ref{cap:bs}) will be complete; in the
continuous setting, completeness will be guaranteed by the Brownian
martingale representation theorem (Chapter~\ref{cap:itocalculus}).
\section{The multi-period binomial model and the probabilistic machinery}
\label{sec:multiperiodo}

The one-period binary model is insufficient as a realistic model of a
market. Its natural generalization, the \emph{multi-period binomial}
model, allows the market to be observed at several intermediate times.
Its treatment requires introducing filtrations, conditional
expectations, martingales and predictable processes, concepts that
formalize natural financial ideas of the model.

\subsection{Description of the multi-period binomial model}
\label{subsec:setup_multiperiodo}

\noindent\textbf{Discretization of the time horizon.} Fix $T > 0$ and
an integer $N \geq 1$. We subdivide $[0, T]$ into $N$ subintervals of
length $\delta t \eqdef T/N$, and consider the times
$t_k \eqdef k\, \delta t$ for $k = 0, 1, \ldots, N$. The market is
observed only at these times. The bond evolves deterministically:
$B_k \eqdef e^{r\, t_k}$.

\medskip

\noindent\textbf{The dynamics of the underlying.} The stock has a known
initial price $S_0 > 0$. We fix constants $u, d$ with $0 < d < u$ and
assume that between two consecutive times the price is multiplied by
$u$ with probability $p$ or by $d$ with probability $1-p$,
independently of the behavior in the previous steps:
\begin{equation*}
  S_{k+1} = \begin{cases}
    S_k\, u & \text{with probability } p, \\[2pt]
    S_k\, d & \text{with probability } 1 - p,
  \end{cases}
\end{equation*}
The structure is visualized by a recombining binomial tree
(Figure~\ref{fig:arbol_binomial}): at time $t_k$, the price can take
$k+1$ distinct values, $S_0\, u^j\, d^{k-j}$ for $j = 0, \ldots, k$.

\begin{figure}[h]
\centering
\begin{tikzpicture}[
  scale=1,
  every node/.style={font=\small},
  edge/.style={->, >=stealth, thick}
]
  \node (S0) at (0,0) {$S_0$};
  \node (S1u) at (2,1) {$S_0 u$};
  \node (S1d) at (2,-1) {$S_0 d$};
  \node (S2uu) at (4,2) {$S_0 u^2$};
  \node (S2ud) at (4,0) {$S_0 ud$};
  \node (S2dd) at (4,-2) {$S_0 d^2$};
  \node (S3uuu) at (6,3) {$S_0 u^3$};
  \node (S3uud) at (6,1) {$S_0 u^2 d$};
  \node (S3udd) at (6,-1) {$S_0 u d^2$};
  \node (S3ddd) at (6,-3) {$S_0 d^3$};
  \draw[edge] (S0) -- (S1u);
  \draw[edge] (S0) -- (S1d);
  \draw[edge] (S1u) -- (S2uu);
  \draw[edge] (S1u) -- (S2ud);
  \draw[edge] (S1d) -- (S2ud);
  \draw[edge] (S1d) -- (S2dd);
  \draw[edge] (S2uu) -- (S3uuu);
  \draw[edge] (S2uu) -- (S3uud);
  \draw[edge] (S2ud) -- (S3uud);
  \draw[edge] (S2ud) -- (S3udd);
  \draw[edge] (S2dd) -- (S3udd);
  \draw[edge] (S2dd) -- (S3ddd);
  \node at (0,-3.7) {$t_0$};
  \node at (2,-3.7) {$t_1$};
  \node at (4,-3.7) {$t_2$};
  \node at (6,-3.7) {$t_3$};
\end{tikzpicture}
\caption{Multi-period binomial tree with $N = 3$ steps. The recombining
structure (the paths $ud$ and $du$ meet at the same node $S_0\,ud$)
reduces the number of nodes from $2^k$ to $k+1$.}
\label{fig:arbol_binomial}
\end{figure}

\medskip

\noindent\textbf{The probability space.} As sample space we take the
set of all possible paths:
\begin{equation*}
  \Omega \eqdef \{u, d\}^N,
\end{equation*}
where each $\omega = (\omega_1, \ldots, \omega_N)$ with
$\omega_j \in \{u, d\}$ represents a complete history of the market.
We take $\F \eqdef 2^{\Omega}$ and, on $(\Omega, \F)$, we define the
market measure $\PP$ by
\begin{equation*}
  \PP(\{\omega\}) \eqdef p^{\#\{j: \omega_j = u\}}\,
                        (1-p)^{\#\{j: \omega_j = d\}},
\end{equation*}
which encodes the independence of the steps. The price at time $t_k$
is the random variable
$S_k(\omega) \eqdef S_0\, \prod_{j=1}^k f(\omega_j)$, where $f(u) = u$
and $f(d) = d$.

\subsection{Filtered probability spaces}
\label{subsec:filtraciones}

In what follows we assume familiarity with the basic notions of
probability (measurable space, $\sigma$-algebra, random variable,
expectation) covered in any \emph{STEM} degree.

\medskip

\noindent\textbf{$\sigma$-algebras as information.} The idea is that a
$\sigma$-algebra $\G \subseteq \F$ on $(\Omega, \F)$ can be
interpreted as the set of those events whose occurrence an observer is
able to determine: if $A \in \G$, the observer can answer the question
``has $A$ occurred?''. The full $\sigma$-algebra $\F$ corresponds to
the omniscient observer; sub-$\sigma$-algebras $\G \subseteq \F$
correspond to observers with partial information.

\medskip

\noindent\textbf{The $\sigma$-algebra $\F_k$ of the binomial model.}
Let us illustrate this reading on the multi-period binomial model. For
each $k \in \{0, 1, \ldots, N\}$, we define $\F_k$ as the
$\sigma$-algebra generated by the partition of $\Omega$ that groups
together the paths agreeing in their first $k$ coordinates:
\begin{equation*}
  A_{(\omega_1, \ldots, \omega_k)} \eqdef
  \bigl\{\omega' \in \Omega : \omega'_j = \omega_j,\ j = 1, \ldots, k\bigr\}.
\end{equation*}
There are $2^k$ such blocks (the \emph{atoms} of $\F_k$). At the
extremes, $\F_0 = \{\emptyset, \Omega\}$ is the trivial
$\sigma$-algebra (having observed nothing, no event can be
distinguished from the whole world) and $\F_N = \F$ (after observing
the complete path, everything is known). On the tree of
Figure~\ref{fig:arbol_binomial}, $A \in \F_k$ if and only if we can
determine whether $\omega \in A$ from the first $k$ steps.

The sequence $(\F_k)_{k=0}^{N}$ is \emph{increasing}:
$\F_k \subseteq \F_{k+1}$ for all $k$. This motivates the following
definition.

\begin{definicion}[Filtration]
\label{def:filtracion}
Let $(\Omega, \F, \PP)$ be a probability space. A \emph{filtration} on
this space is a sequence $(\F_k)_{k=0}^{N}$ of sub-$\sigma$-algebras
of $\F$ with $\F_k \subseteq \F_{k+1} \subseteq \F$ for all $k$. The
quadruple $(\Omega, \F, (\F_k), \PP)$ is called a \emph{filtered
probability space}.
\end{definicion}

From the financial point of view, the filtration of the binomial model
is the \emph{market filtration}.

\medskip

\begin{definicion}[Adapted process]
\label{def:adaptado}
A \emph{stochastic process} is a sequence $(X_k)_{k=0}^{N}$ of random
variables on $(\Omega, \F)$. The process is \emph{adapted} to the
filtration $(\F_k)$ if, for each $k$, $X_k$ is $\F_k$-measurable.
\end{definicion}

Measurability with respect to $\F_k$ formalizes the fact that the
value of $X_k$ is known after observing what has happened up to time
$t_k$. In the binomial model, $S_k$ is by construction a function of
the first $k$ coordinates and hence $\F_k$-measurable: the price
process is adapted.

\medskip

\begin{definicion}[Natural filtration]
\label{def:filtracion_natural}
The \emph{natural filtration} of a process $(X_k)_{k=0}^{N}$ is
$(\F_k^X)_{k=0}^{N}$ with $\F_k^X \eqdef \sigma(X_0, X_1, \ldots, X_k)$,
the smallest $\sigma$-algebra making the variables $X_0, \ldots, X_k$
measurable.
\end{definicion}

It is the smallest filtration with respect to which $(X_k)$ is
adapted. In the multi-period binomial model, the market filtration
coincides with the natural filtration of the price, and both will be
identified and denoted simply by $(\F_k)$.

\subsection{Conditional expectation}
\label{subsec:esperanza_condicional}

\noindent\textbf{Starting point.} The reader will recall the
elementary notion of conditional expectation given an event $B \in \F$
with $\PP(B) > 0$: the number $\E[X \mid B] \eqdef \E[X\, \1_B]/\PP(B)$,
which represents the mean value of $X$ in those scenarios in which
$B$ has occurred. The notion we shall need is of a different nature:
we do not want to condition on an individual event, but on a whole
$\sigma$-algebra $\G \subseteq \F$, and the result will no longer be a
number but a random variable.

\medskip

\noindent\textbf{Motivation on the tree.} Let $X$ be an
$\F_3$-measurable random variable on the binomial tree with $N = 3$
steps, and let us ask for its mean value ``given the information
$\F_2$''. Intuitively, standing at time $t_2$ the observer has reached
one of the $4$ atoms of $\F_2$; in each of them there are two
compatible paths, corresponding to the two outcomes of the third
step. The conditional expectation given $\F_2$ must assign to each
atom the mean value of $X$ over the two compatible paths. It is an
$\F_2$-measurable random variable (constant on each atom of $\F_2$).
The axiomatic definition captures this idea without resorting to the
particular structure of the model.

\medskip

\noindent\textbf{Kolmogorov's definition.}

\begin{definicion}[Conditional expectation]
\label{def:esperanza_condicional}
Let $X$ be an integrable random variable on $(\Omega, \F, \PP)$ and
$\G \subseteq \F$ a sub-$\sigma$-algebra. A random variable $Y$ is a
\emph{conditional expectation of $X$ given $\G$} if:
\begin{enumerate}[label=\textnormal{(\roman*)}]
  \item $Y$ is $\G$-measurable;
  \item for all $A \in \G$, $\displaystyle \int_A Y \, d\PP = \int_A X \, d\PP$.
\end{enumerate}
\end{definicion}

\begin{teorema}[Existence and uniqueness]
\label{teo:existencia_esperanza_condicional}
Let $X$ be integrable and $\G \subseteq \F$. There exists a
conditional expectation of $X$ given $\G$, and any two conditional
expectations coincide $\PP$-almost surely.
\end{teorema}

The proof relies on the Radon-Nikodym theorem applied to the measure
$\nu(A) \eqdef \int_A X \, d\PP$ on $(\Omega, \G)$, which is
absolutely continuous with respect to $\PP\!\restriction_\G$ and
admits a density $Y$ satisfying (i) and (ii). It can be found in
\cite{williams}. Since the conditional expectation is unique up to
almost sure equality, we shall speak of ``\emph{the}'' conditional
expectation of $X$ given $\G$, denoted $\E[X \mid \G]$, with the
understanding that any identity in which it appears holds
$\PP$-almost surely.

Condition (i) formalizes that the conditional expectation can be
determined from the information $\G$. Condition (ii) is the
\emph{conservation of averages}: on any $A \in \G$, the integrals of
$\E[X \mid \G]$ and $X$ coincide. Taking $A = \Omega$,
$\E[\E[X \mid \G]] = \E[X]$.

\medskip

\noindent\textbf{Fundamental properties.} In what follows, $X, Y$ are
integrable random variables and $\G, \HH \subseteq \F$ are
sub-$\sigma$-algebras.

\begin{proposicion}[Properties of conditional expectation]
\label{prop:propiedades_esperanza_cond}
The following properties hold $\PP$-almost surely.
\begin{enumerate}[label=\textnormal{(\roman*)}]
  \item \emph{Linearity}: $\E[a X + b Y \mid \G] = a\, \E[X \mid \G]
        + b\, \E[Y \mid \G]$ for any $a, b \in \R$.
  \item \emph{Known measurable variables}: if $X$ is
        $\G$-measurable, then $\E[X \mid \G] = X$.
  \item \emph{Taking out what is known}: if $Y$ is $\G$-measurable and
        $XY$ is integrable, then $\E[XY \mid \G] = Y\, \E[X \mid \G]$.
  \item \emph{Tower property}: if $\HH \subseteq \G$, then
        $\E\!\bigl[\E[X \mid \G] \,\big|\, \HH\bigr] = \E[X \mid \HH]$.
  \item Independence: if $X$ is independent of $\G$, then
        $\E[X \mid \G] = \E[X]$.
\end{enumerate}
\end{proposicion}

By way of illustration (and because it will be the most frequently
used property) we prove the tower property.

\begin{proof}[Proof of the tower property]
Let $Z \eqdef \E\bigl[\E[X \mid \G] \,\big|\, \HH\bigr]$. We must
check that $Z$ satisfies the two conditions of
Definition~\ref{def:esperanza_condicional} applied to $X$ and $\HH$.
The first is immediate: $Z$ is $\HH$-measurable by construction. For
the second, let $A \in \HH$. Since $\HH \subseteq \G$, also $A \in \G$.
Applying (ii) first to $Z$ with $\HH$ and then to $\E[X \mid \G]$ with
$\G$,
\begin{equation*}
  \int_A Z \, d\PP = \int_A \E[X \mid \G] \, d\PP = \int_A X \, d\PP.
\end{equation*}
By uniqueness, $Z = \E[X \mid \HH]$ almost surely.
\end{proof}

The intuitive reading of the tower property is the following:
conditioning first on richer information $\G$ and then on coarser
information $\HH$ is equivalent to conditioning directly on $\HH$.

\medskip

\noindent\textbf{Shorthand notation.} If $Y_1, \ldots, Y_k$ are random
variables, we write $\E[X \mid Y_1, \ldots, Y_k] \eqdef
\E[X \mid \sigma(Y_1, \ldots, Y_k)]$.

\subsection{Discrete-time martingales}
\label{subsec:martingalas_discretas}

The property $\E_\Q[\widetilde S_T] = \widetilde S_0$
(eq.~\eqref{eq:martingala_uniperiodo}) that characterized the
risk-neutral measure of the one-period model admits a natural
generalization: the concept of \emph{martingale}.

\medskip

\noindent\textbf{Definition and first properties.} We work on a
filtered probability space $(\Omega, \F, (\F_k), \PP)$.

\begin{definicion}[Martingale]
\label{def:martingala}
A process $(M_k)_{k=0}^{N}$ adapted to $(\F_k)$ is a
\emph{martingale} with respect to $\PP$ if:
\begin{enumerate}[label=\textnormal{(\roman*)}]
  \item $M_k$ is integrable for all $k$;
  \item $\E[M_{k+1} \mid \F_k] = M_k$ almost surely, for $k = 0,
        \ldots, N-1$.
\end{enumerate}
\end{definicion}

Condition (ii) can be read as follows: the best prediction of the
value of the process at time $k+1$, given the information up to $k$,
is the current value $M_k$. A martingale is an ``unbiased'' process.
Iterating (ii) and applying the tower property, for
$0 \leq j \leq k \leq N$,
\begin{equation}\label{eq:martingala_iterada}
  \E[M_k \mid \F_j] = M_j \quad \text{a.s.}
\end{equation}
In particular, $\E[M_k] = \E[M_0]$ for all $k$: the expectation of a
martingale is constant in time.

\medskip

The martingale property depends on the measure: replacing $\PP$ by a
suitable measure $\Q$ will turn the discounted price of the underlying
into a martingale.

\medskip

\noindent\textbf{The risk-neutral measure in the multi-period model.}
The no-arbitrage condition becomes, taking $\delta t = T/N$,
\begin{equation}\label{eq:cond_noarb_multi}
  d < e^{r \delta t} < u,
\end{equation}
which we assume from now on. We define
\begin{equation}\label{eq:def_q_multi}
  q \eqdef \frac{e^{r \delta t} - d}{u - d}\, \in (0, 1),
\end{equation}
in analogy with the one-period case. On $(\Omega, \F)$, we define the
measure $\Q$ by
\begin{equation}\label{eq:def_Q_multi}
  \Q(\{\omega\}) \eqdef q^{\#\{j: \omega_j = u\}}\,
                       (1-q)^{\#\{j: \omega_j = d\}}.
\end{equation}
Under $\Q$, the steps are independent with probability $q$ of an up
move in each. Since $q \in (0, 1)$ and $p \in (0, 1)$, the measures
$\Q$ and $\PP$ assign strictly positive probability to the same paths,
that is, $\Q \sim \PP$.

\medskip

\noindent\textbf{The discounted price as a martingale under $\Q$.} Let
$\widetilde S_k \eqdef e^{-r\, t_k}\, S_k$ be the discounted price.

\begin{teorema}[Martingale property of the discounted price]
\label{teo:martingala_precio_descontado}
In the multi-period binomial model, under the measure $\Q$
of~\eqref{eq:def_Q_multi}, the process $(\widetilde S_k)_{k=0}^{N}$ is
a martingale with respect to $(\F_k)$.
\end{teorema}

\begin{proof}
Adaptedness is immediate since $\widetilde S_k$ is a function of
$S_k$, which is $\F_k$-measurable. Integrability is trivial. For the
martingale property, write $S_{k+1} = S_k\, f_{k+1}$, where $f_{k+1}$
equals $u$ with probability $q$ and $d$ with probability $1-q$ under
$\Q$, independently of $\F_k$. Then
\begin{equation*}
  \widetilde S_{k+1} = e^{-r \delta t}\, \widetilde S_k\, f_{k+1}.
\end{equation*}
Applying ``taking out what is known'' (Proposition
\ref{prop:propiedades_esperanza_cond}) and independence,
\begin{align*}
  \E_\Q[\widetilde S_{k+1} \mid \F_k]
  &= e^{-r \delta t}\, \widetilde S_k\, \E_\Q[f_{k+1} \mid \F_k]
   = e^{-r \delta t}\, \widetilde S_k\, \E_\Q[f_{k+1}] \\
  &= e^{-r \delta t}\, \widetilde S_k\,\bigl[q\, u + (1 - q)\, d\bigr]
   = e^{-r \delta t}\, \widetilde S_k\, e^{r \delta t}
   = \widetilde S_k.
\end{align*}
\end{proof}

The measure $\Q$ receives the same names as in the one-period case:
\emph{risk-neutral measure} or \emph{equivalent martingale measure}.
Its uniqueness will be proved in
Subsection~\ref{subsec:representacion_binomial}.

\subsection{Predictable processes and the discrete stochastic integral}
\label{subsec:predecibles_int_disc}

\medskip

\begin{definicion}[Predictable process]
\label{def:predecible}
A process $(H_k)_{k=1}^{N}$ is \emph{predictable} with respect to the
filtration $(\F_k)$ if, for each $k = 1, \ldots, N$, $H_k$ is
$\F_{k-1}$-measurable.
\end{definicion}

The difference from the notion of adapted process
(Definition~\ref{def:adaptado}) is a shift of a single index, but its
content is essential: an adapted process may use the information
available \emph{up to} $t_k$, whereas a predictable one may only use
that available \emph{before} $t_k$. In financial language, $H_k$ is
chosen at time $t_{k-1}$, before knowing how the market will move in
the next step. It is this restriction that rules out strategies
``with foresight''.

\medskip

\begin{definicion}[Discrete stochastic integral]
\label{def:integral_discreta}
Let $(M_k)_{k=0}^{N}$ be adapted and $(H_k)_{k=1}^{N}$ predictable.
The \emph{discrete stochastic integral} of $H$ with respect to $M$ is
the process
\begin{equation}\label{eq:integral_discreta}
  (H \cdot M)_n \eqdef \sum_{k=1}^{n} H_k\, (M_k - M_{k-1}),
  \qquad n = 0, 1, \ldots, N,
\end{equation}
with $(H \cdot M)_0 \eqdef 0$.
\end{definicion}

The formula has a suggestive financial reading: if $M_k$ is the price
of an asset and $H_k$ the quantity held during $(t_{k-1}, t_k]$, then
$H_k\, (M_k - M_{k-1})$ is the gain generated in that period, and
$(H \cdot M)_n$ is the total gain up to $t_n$.

The name ``discrete stochastic integral'' anticipates the role of this
object in the continuous setting: replacing the sum by an integral and
the increments $M_k - M_{k-1}$ by infinitesimal increments $dM_t$, the
formula becomes $\int_0^t H_s\, dM_s$, the \emph{Itô integral} of
Chapter~\ref{cap:itocalculus}.

\medskip

\begin{teorema}[Stability of martingales under stochastic integration]
\label{teo:integral_preserva_martingala}
Let $(M_k)$ be a $\PP$-martingale with respect to $(\F_k)$ and $(H_k)$
predictable and bounded. Then $\bigl((H \cdot M)_n\bigr)_{n=0}^{N}$ is
a $\PP$-martingale with respect to the same filtration.
\end{teorema}

\begin{proof}
\emph{Adaptedness and integrability.} For each $n$, $(H \cdot M)_n$ is
a sum of products $H_k\, (M_k - M_{k-1})$ with $k \leq n$, all of them
$\F_n$-measurable. Integrability follows from the boundedness of $H$
and the integrability of $M$.

\emph{Martingale property.} Since
$(H \cdot M)_{n+1} = (H \cdot M)_n + H_{n+1}(M_{n+1} - M_n)$, by
linearity
\begin{equation*}
  \E\bigl[(H \cdot M)_{n+1} \,\big|\, \F_n\bigr]
  = (H \cdot M)_n + \E\bigl[H_{n+1}(M_{n+1} - M_n) \,\big|\, \F_n\bigr],
\end{equation*}
where we have used that $(H\cdot M)_n$ is $\F_n$-measurable. Since
$H_{n+1}$ is $\F_n$-measurable (by predictability), the property
``taking out what is known'' gives
\begin{equation*}
  \E\bigl[H_{n+1}(M_{n+1} - M_n) \,\big|\, \F_n\bigr]
  = H_{n+1}\, \E[M_{n+1} - M_n \mid \F_n] = 0,
\end{equation*}
where the last equality uses the martingale property of $(M_k)$. We
conclude that $\E[(H \cdot M)_{n+1} \mid \F_n] = (H \cdot M)_n$.
\end{proof}

The fact that predictability of the integrand + martingale property of
the integrator $\Rightarrow$ martingale property of the integral is
the structural idea that will reappear in the continuous setting. As
an immediate consequence, $\E[(H \cdot M)_n] = 0$ for all $n$: the
expected gain of any predictable strategy on a martingale is zero.

\subsection{The binomial representation theorem}
\label{subsec:representacion_binomial}

We close the section with the result that ties the whole theory
together: the \emph{binomial representation theorem}. It states that
every $\Q$-martingale can be written as a discrete stochastic integral
with respect to the discounted price of the underlying. It is the
discrete and constructive version of the \emph{Brownian martingale
representation theorem} of Chapter~\ref{cap:itocalculus}: the proof
we give here provides an explicit formula for the integrand, whereas
in the continuous setting we shall only have its existence.

We work on the multi-period binomial model endowed with the
risk-neutral measure $\Q$ and the filtration $(\F_k)$. We write
$\widetilde S_k = e^{-r k\, \delta t} S_k$.

\begin{teorema}[Binomial representation]
\label{teo:representacion_binomial}
Let $(M_k)_{k=0}^{N}$ be a $\Q$-martingale with respect to $(\F_k)$.
There exists a unique predictable process $(\phi_k)_{k=1}^N$ such that
\begin{equation}\label{eq:representacion_binomial}
  M_n = M_0 + \sum_{k=1}^{n} \phi_k\, (\widetilde S_k - \widetilde S_{k-1}),
  \qquad n = 0, 1, \ldots, N.
\end{equation}
\end{teorema}

\begin{proof}
We construct $\phi$ node by node. Let $k \in \{1, \ldots, N\}$ and let
$A$ be an atom of $\F_{k-1}$. By the structure of the model, $A$
decomposes as the disjoint union of two atoms $A_u, A_d$ of $\F_k$. On
each of them the $\F_k$-measurable variables take a constant value;
denote by $M_{k-1}^A$, $M_k^{A_u}$, $M_k^{A_d}$,
$\widetilde S_{k-1}^A$, $\widetilde S_k^{A_u}$, $\widetilde S_k^{A_d}$
the corresponding values. We want $\phi_k$ constant on $A$ (denote it
$\phi_k^A$) such that $\eqref{eq:representacion_binomial}$ holds on
$A_u$ and $A_d$:
\begin{equation*}
\left\{
\begin{aligned}
  M_k^{A_u} &= M_{k-1}^A + \phi_k^A\,
              \bigl(\widetilde S_k^{A_u} - \widetilde S_{k-1}^A\bigr), \\
  M_k^{A_d} &= M_{k-1}^A + \phi_k^A\,
              \bigl(\widetilde S_k^{A_d} - \widetilde S_{k-1}^A\bigr).
\end{aligned}
\right.
\end{equation*}
Subtracting eliminates $M_{k-1}^A$:
\begin{equation*}
  M_k^{A_u} - M_k^{A_d}
  = \phi_k^A\, \bigl(\widetilde S_k^{A_u} - \widetilde S_k^{A_d}\bigr).
\end{equation*}
Since $\widetilde S_k^{A_u} \neq \widetilde S_k^{A_d}$,
\begin{equation}\label{eq:phi_explicito}
  \phi_k^A = \frac{M_k^{A_u} - M_k^{A_d}}{\widetilde S_k^{A_u}
                                          - \widetilde S_k^{A_d}}.
\end{equation}
Formula~\eqref{eq:phi_explicito} determines $\phi_k$ uniquely on each
atom of $\F_{k-1}$. The first equation is satisfied by construction;
it remains to verify the second. Multiplying the first by $q$ and the
second by $1 - q$ and adding,
\begin{equation*}
  q\, M_k^{A_u} + (1-q)\, M_k^{A_d}
  = M_{k-1}^A + \phi_k^A\,
    \bigl[q\, \widetilde S_k^{A_u} + (1-q)\, \widetilde S_k^{A_d}
          - \widetilde S_{k-1}^A\bigr].
\end{equation*}
The left-hand side is
$\E_\Q[M_k \mid \F_{k-1}]\restriction_A = M_{k-1}^A$ (martingale
property of $M$); the bracket on the right-hand side is
$\E_\Q[\widetilde S_k \mid \F_{k-1}]\restriction_A - \widetilde S_{k-1}^A
= \widetilde S_{k-1}^A - \widetilde S_{k-1}^A = 0$ by the martingale
property of $(\widetilde S_k)$
(Theorem~\ref{teo:martingala_precio_descontado}). The identity holds
trivially, and combining it with the first equation recovers the
second. Repeating this on each atom of each $\F_{k-1}$ yields a
predictable process $(\phi_k)_{k=1}^N$
satisfying~\eqref{eq:representacion_binomial}. Uniqueness is immediate
from~\eqref{eq:phi_explicito}.
\end{proof}

\medskip

\noindent\textbf{Reading of the integrand.}
Formula~\eqref{eq:phi_explicito} expresses $\phi_k$ as the ratio of
the change in $M$ to the change in the discounted price between the
two children of the atom. If we identify $M$ with the discounted value
of a derivative, $\phi_k$ is the quantity of stock that must be held
in order to absorb its possible change: it is the \emph{delta} of the
derivative at each node.

\medskip

\noindent\textbf{Two consequences.} The first is the \emph{uniqueness
of the risk-neutral measure}, anticipated in
Subsection~\ref{subsec:martingalas_discretas}. Let $\Q'$ be another measure under
which $(\widetilde S_k)$ is a martingale. For each atom $A$ of
$\F_{k-1}$, the martingale condition
\begin{equation*}
  \Q'(A_u \mid A)\, \widetilde S_k^{A_u}
  + \Q'(A_d \mid A)\, \widetilde S_k^{A_d}
  = \widetilde S_{k-1}^A,
\end{equation*}
is linear in $\Q'(A_u \mid A) \in (0, 1)$ with positive coefficients,
and has a unique solution, which coincides with $q$. Consequently
$\Q' = \Q$. The second consequence is the \emph{completeness of the
multi-period binomial model}, which will be fully justified in
Subsection~\ref{subsec:tres_pasos_discreto}: since every
$\Q$-martingale can be represented as a stochastic integral with
respect to $\widetilde S$, the $\Q$-martingale associated with the
discounted value of a derivative will give rise, via representation,
to the integrand that constitutes the stock holding of the replicating
portfolio.
\section{Solution of the problem in discrete time}
\label{sec:solucion_discreto}

We now bring together the results developed in
Section~\ref{sec:multiperiodo} to obtain the fair price formula and
the replicating strategy. The construction is carried out in three
steps.

\subsection{The three steps toward replication}
\label{subsec:tres_pasos_discreto}

\noindent\textbf{Self-financing strategies.} The idea is that a
self-financing strategy is one in which rebalancing the portfolio
requires neither injecting nor withdrawing capital from outside.

\begin{definicion}[Self-financing strategy]
\label{def:self_financing_discreto}
A \emph{strategy} in the multi-period binomial model is a pair of
processes $\bigl((\phi_k), (\psi_k)\bigr)$ predictable with respect to
$(\F_k)$. Its \emph{value} at $t_k$ is
$V_k \eqdef \phi_k\, S_k + \psi_k\, B_k$ (with
$V_0 \eqdef \phi_1\, S_0 + \psi_1\, B_0$). The strategy is
\emph{self-financing} if, for all $k = 1, \ldots, N - 1$,
\begin{equation}\label{eq:cond_self_financing_discreto}
  \phi_k\, S_k + \psi_k\, B_k = \phi_{k+1}\, S_k + \psi_{k+1}\, B_k.
\end{equation}
\end{definicion}

\begin{proposicion}[Characterization via the discounted process]
\label{prop:caracterizacion_self_financing}
Let $\bigl((\phi_k), (\psi_k)\bigr)$ be a strategy and
$\widetilde V_k \eqdef e^{-r\, t_k} V_k$. The strategy is
self-financing if and only if
\begin{equation}\label{eq:self_financing_descontado}
  \widetilde V_n = V_0 + \sum_{k=1}^{n} \phi_k\,
                   \bigl(\widetilde S_k - \widetilde S_{k-1}\bigr),
  \qquad n = 0, 1, \ldots, N.
\end{equation}
\end{proposicion}

\begin{proof}
$(\Rightarrow)$ Let us compute $\widetilde V_k - \widetilde V_{k-1}$.
Since $\widetilde V_k = \phi_k\, \widetilde S_k + \psi_k$,
\begin{align*}
  \widetilde V_k - \widetilde V_{k-1}
  &= \phi_k\, (\widetilde S_k - \widetilde S_{k-1})
   + (\phi_k - \phi_{k-1})\, \widetilde S_{k-1}
   + (\psi_k - \psi_{k-1}).
\end{align*}
The self-financing condition
\eqref{eq:cond_self_financing_discreto} at $t_{k-1}$, divided by
$B_{k-1}$, implies
$(\phi_k - \phi_{k-1}) \widetilde S_{k-1} + (\psi_k - \psi_{k-1}) = 0$.
Consequently
$\widetilde V_k - \widetilde V_{k-1} = \phi_k(\widetilde S_k - \widetilde S_{k-1})$.
Summing yields~\eqref{eq:self_financing_descontado}.

$(\Leftarrow)$ If \eqref{eq:self_financing_descontado} holds,
subtracting the cases $n$ and $n-1$ and combining with
$\widetilde V_k = \phi_k\, \widetilde S_k + \psi_k$ recovers
\eqref{eq:cond_self_financing_discreto}.
\end{proof}

The proposition links the financial notion of ``self-financing
strategy'' with the mathematical notion of ``discrete stochastic
integral''. In particular, combining it with
Theorem~\ref{teo:integral_preserva_martingala}, the discounted value
of every self-financing strategy is a $\Q$-martingale.

\medskip

\noindent\textbf{The three-step recipe.} Let $C_T$ be an
$\F_N$-measurable random variable (the \emph{payoff} of the
derivative).

\begin{quote}
\textbf{Recipe for pricing by replication.}
\begin{enumerate}[label=\textbf{Step \arabic*.}, leftmargin=*]
  \item \emph{Change of measure.} Determine the unique measure $\Q$
        under which $(\widetilde S_k)$ is a martingale
        (Section~\ref{subsec:martingalas_discretas}).
  \item \emph{Construction of the discounted value.} Define
        \begin{equation}\label{eq:def_M_replicante_discreto}
          M_k \eqdef \E_\Q\bigl[e^{-rT}\, C_T \,\big|\, \F_k\bigr].
        \end{equation}
        By construction, $(M_k)$ is a $\Q$-martingale with
        $M_N = e^{-rT} C_T$.
  \item \emph{Application of the representation.} By
        Theorem~\ref{teo:representacion_binomial}, there exists a
        unique predictable $(\phi_k)$ such that
        $M_k = M_0 + (\phi \cdot \widetilde S)_k$. Define
        $\psi_k \eqdef M_k - \phi_k \widetilde S_k$.
\end{enumerate}
\end{quote}

\begin{teorema}[Pricing by replication in the multi-period binomial model]
\label{teo:valoracion_binomial}
In the arbitrage-free multi-period binomial model, every derivative
with \emph{payoff} $C_T$ admits a self-financing replicating strategy.
The fair price at $t_n$ is
\begin{equation}\label{eq:precio_binomial_multi}
  V_n = \E_\Q\bigl[e^{-r(T - t_n)}\, C_T \,\big|\, \F_n\bigr],
\end{equation}
and at $t = 0$, $V_0 = \E_\Q[e^{-rT}\, C_T]$.
\end{teorema}

\begin{proof}
\emph{Step 1.} The existence and uniqueness of $\Q$ were established
in Section~\ref{subsec:martingalas_discretas}.

\emph{Step 2.} The process $(M_k)$ defined
in~\eqref{eq:def_M_replicante_discreto} is a $\Q$-martingale by the
tower property: for $j \leq k$,
\begin{equation*}
  \E_\Q[M_k \mid \F_j]
  = \E_\Q\bigl[\E_\Q[e^{-rT} C_T \mid \F_k] \,\big|\, \F_j\bigr]
  = \E_\Q[e^{-rT} C_T \mid \F_j] = M_j.
\end{equation*}

\emph{Step 3.} By Theorem~\ref{teo:representacion_binomial} there
exists a predictable $(\phi_k)$ such that
$M_n = M_0 + \sum_{k=1}^n \phi_k(\widetilde S_k - \widetilde S_{k-1})$.
We define $\psi_k \eqdef M_k - \phi_k\, \widetilde S_k$, which is
predictable (a combination of $\F_{k-1}$-measurable quantities), and
$V_n \eqdef e^{r t_n} M_n$. By
Proposition~\ref{prop:caracterizacion_self_financing}, the
representation of $M_n$ as a stochastic integral is equivalent to the
self-financing condition with $V_0 = M_0$ and $\widetilde V_n = M_n$.
The strategy is self-financing and replicates the derivative, since
$V_N = e^{rT}\, M_N = e^{rT}\, e^{-rT}\, C_T = C_T$.

Finally, by the no-arbitrage principle, the fair price of the
derivative at $t_n$ is the value of the portfolio at that time:
\begin{equation*}
  V_n = e^{r t_n}\, M_n = \E_\Q\bigl[e^{-r(T - t_n)}\, C_T \mid \F_n\bigr].
\end{equation*}
\end{proof}

Formula~\eqref{eq:precio_binomial_multi} generalizes
\eqref{eq:precio_esperanza} to the multi-period setting: the fair
price is the discounted value of the expectation of the \emph{payoff}
under the risk-neutral measure, conditioned on the information at the
time of pricing.

\subsection{Passage to continuous time}
\label{subsec:obertura_continuo}

We close the chapter with a look ahead. The multi-period binomial
model has an artificial feature: the market is observed only at a
finite number of equally spaced times. This subsection argues, in a
deliberately heuristic way, that the progressive refinement of the
time partition (letting $N \to \infty$ with $T$ fixed and
$\delta t = T/N \to 0$) leads to a continuous-time model whose
underlying follows a geometric Brownian motion. The rigorous
construction of this model, as well as the associated pricing formula
(the Black-Scholes formula), will be the subject of the following
chapters.

\medskip

\noindent\textbf{Calibration of the factors $u$ and $d$.} If we took
$u, d$ constant (independent of $\delta t$), the price paths would
become dramatically erratic as $N$ increases. A sensible choice
requires the log-returns at each step to be compatible with a stable
behavior in the limit. Specifically, let $\sigma > 0$ and $\nu \in \R$
be two constants that will represent the \emph{volatility} and the
\emph{trend} of the price in the limiting model. Let us calibrate
\begin{equation}\label{eq:calibracion_ud}
  u = e^{\sigma \sqrt{\delta t}}, \qquad
  d = e^{-\sigma \sqrt{\delta t}},
\end{equation}
and the subjective probability as
$p = \tfrac{1}{2}(1 + \tfrac{\nu}{\sigma}\sqrt{\delta t})$. Under this
calibration, the log-return $\log(S_{k+1}/S_k)$ equals
$+\sigma\sqrt{\delta t}$ with probability $p$ and
$-\sigma\sqrt{\delta t}$ with probability $1-p$. Its expectation is
$\nu\, \delta t$ and its variance $\sigma^2\, \delta t + O(\delta t^2)$.

\medskip

\noindent\textbf{Convergence to geometric Brownian motion.} We have
\begin{equation*}
  S_N = S_0\, \exp\!\left(\sum_{k=1}^{N} \log\frac{S_k}{S_{k-1}}\right).
\end{equation*}
The sum in the exponent is a sum of independent identically
distributed variables, each with expectation $\nu \delta t$ and
variance $\sigma^2 \delta t + O(\delta t^2)$. The total sum has
expectation $\nu T$ and variance tending to $\sigma^2 T$ as
$N \to \infty$. By the Central Limit Theorem,
\begin{equation*}
  \log\!\frac{S_T}{S_0} \xrightarrow{d} \mathcal{N}(\nu T,\, \sigma^2 T).
\end{equation*}
Equivalently,
\begin{equation}\label{eq:gbm_limite}
  S_T = S_0\, \exp\bigl(\nu T + \sigma\, W_T\bigr),
\end{equation}
with $W_T \sim \mathcal{N}(0, T)$. The argument extends to any
intermediate time $t \in [0, T]$, giving
\begin{equation*}
  S_t = S_0\, \exp\bigl(\nu t + \sigma\, W_t\bigr),
\end{equation*}
with $W_t \sim \mathcal{N}(0, t)$. The family $\{W_t\}_{t \geq 0}$
emerging in the limit is what is known as \emph{Brownian motion}: a
continuous-time stochastic process with independent Gaussian
increments and continuous paths. Its rigorous construction will be
addressed in Chapter~\ref{cap:brown}. The process $S_t$ defined
by~\eqref{eq:gbm_limite} is called \emph{geometric Brownian motion},
and it will be the canonical model for the stock price in the
continuous setting.

Under the risk-neutral measure $\Q$ of the discrete model, an
analogous computation (replacing $p$ by $q$ as defined
in~\eqref{eq:def_q_multi} and expanding in powers of $\delta t$) shows
that the drift of the log-price under $\Q$ becomes
$r - \tfrac{1}{2}\sigma^2$ instead of $\nu$. This \emph{change of
drift} induced by the change of measure is the discrete and heuristic
version of the phenomenon that in the continuous setting will give
rise to \emph{Girsanov's theorem}, one of the central pieces of
Chapter~\ref{cap:itocalculus}.

\medskip

\noindent\textbf{Summary of the chapter.} Starting from the static
pricing of a \emph{forward} by a replication argument, we have
progressed to a complete and rigorous formula, in the multi-period
binomial model, for the price of any European derivative, expressed as
the discounted expectation of the \emph{payoff} under a risk-neutral
measure, together with a self-financing strategy that replicates it.
The final three-step articulation is the scheme that will reappear in
the continuous setting of Black-Scholes. Each of the mathematical
concepts introduced (filtrations, conditional expectation,
martingales, predictable processes, discrete stochastic integral) has
a continuous-time version that we shall employ in
Chapter~\ref{cap:itocalculus}. Finally, the introduction given in this
subsection suggests that the underlying in the continuous setting will
be a process with independent Gaussian increments and continuous
paths: \emph{Brownian motion}, to which we devote
Chapter~\ref{cap:brown}.

% =====================================================================
%   CHAPTER 3 — BROWNIAN MOTION
% =====================================================================
\chapter{Brownian motion}\label{cap:brown}

This chapter introduces \emph{Brownian motion}, the stochastic
foundation of the continuous-time pricing theory. The financial
motivation now recedes into the background: the tools we shall need
require a certain amount of self-contained development before we
return to apply them.

\section{Definition and existence}
\label{sec:browniano_definicion}

\subsection{Formal definition of Brownian motion}
\label{subsec:browniano_def}

\begin{definicion}[Standard Brownian motion]
\label{def:browniano}
A \emph{standard Brownian motion} on a probability space
$(\Omega, \F, \PP)$ is a family $\{W_t\}_{t \geq 0}$ of real random
variables satisfying:
\begin{enumerate}[label=\textnormal{(\roman*)}]
  \item $W_0 = 0$ almost surely;
  \item for any $0 \leq s < t$, the increment $W_t - W_s$ has
        distribution $\mathcal{N}(0,\, t - s)$;
  \item increments over disjoint intervals are independent, that is,
        for every $n \geq 1$ and every choice of times
        $0 \leq t_0 < t_1 < \cdots < t_n$, the random variables
        \begin{equation*}
          W_{t_1} - W_{t_0},\ W_{t_2} - W_{t_1},\ \ldots,\
          W_{t_n} - W_{t_{n-1}}
        \end{equation*}
        are mutually independent;
  \item the paths $t \mapsto W_t(\omega)$ are continuous functions
        almost surely.
\end{enumerate}
\end{definicion}

Conditions (ii)–(iii) inherit, via the CLT, the Gaussian and
independent nature of the steps of the rescaled random walk; (i) is a
mere normalization (any process satisfying (ii)–(iv) can be shifted
so as to start at zero) and (iv) is the condition that qualitatively
distinguishes Brownian motion: continuous paths, without jumps.

It is worth noting that conditions (i)–(iii) are
\emph{distributional}: they specify the joint law of the process,
whereas (iv) is \emph{pathwise}: it requires each realization
$t \mapsto W_t(\omega)$ to be a continuous function.

\newpage

\subsection{Existence of Brownian motion}
\label{subsec:browniano_existencia}

The preceding definition immediately raises the question of whether
the mathematical object it defines actually exists. The answer is
affirmative, but its proof is not trivial: it requires constructing a
probability space $(\Omega, \F, \PP)$ and a family of random variables
$\{W_t\}_{t \geq 0}$ on it satisfying all the properties
simultaneously, which does not follow from elementary facts of
probability.

\begin{teorema}[Existence of Brownian motion]
\label{teo:existencia_browniano}
There exist a probability space $(\Omega, \F, \PP)$ and a family of
random variables $\{W_t\}_{t \geq 0}$ on it satisfying conditions
(i)--(iv) of Definition~\ref{def:browniano}.
\end{teorema}

The rigorous proof of this theorem was originally provided by Norbert
Wiener in 1923 and, in a more elementary version based on the Schauder
system of functions, by Paul Lévy in 1939. The essential idea of
Lévy's construction is the following: the values of $W$ on the dyadic
rationals of the interval $[0, 1]$ are defined by an infinite series
of successive adjustments, controlling that the partial paths converge
uniformly with probability one; the process so constructed is then
extended by continuity to the whole interval. The detailed proof can
be found in \cite{karatzas-shreve} or \cite{williams}. From this point
on, we take the existence of Brownian motion for granted and focus on
studying its properties.

\begin{observacion}[On the continuity condition in the definition]
\label{obs:continuidad_definicion}
Many authors define Brownian motion by requiring only~(i)--(iii),
which are distributional in nature: as already pointed out,
specifying~(i)--(iii) amounts to fixing the finite-dimensional
distributions of the process (the joint laws of
$(W_{t_1}, \dots, W_{t_n})$, which turn out to be Gaussian with zero
mean and covariance $\E[W_s W_t] = \min(s, t)$). Once these
distributions are fixed, the \emph{Kolmogorov extension theorem}
guarantees the existence of a process with those distributions.
However, this result does not guarantee that its paths are
continuous: Kolmogorov's construction yields the process from its
finite-dimensional distributions, but by itself it provides no
information on the regularity of the paths $t \mapsto W_t(\omega)$.

To obtain continuous paths one further step is needed. By
\emph{Doob's theorem}, the process admits a separable version, that
is, a process with the same finite-dimensional distributions whose
behavior is determined by the values it takes on a countable dense set
of times. It can then be shown that this version has continuous paths
almost surely. This is the usual way of establishing the existence of
a continuous Brownian motion in a course on stochastic processes, and
it constitutes an alternative to Lévy's explicit construction
presented above.

Requiring~(iv) from the outset, as we do here, is therefore not an
extra condition: continuity is obtained at no additional cost by
passing to the separable version, and indeed several authors choose
to build it into the definition. This choice is moreover convenient
for our purposes, since all the stochastic calculus we develop from
Chapter~\ref{cap:itocalculus} onward, starting with the very
construction of the Itô integral, depends on the continuity of the
Brownian paths.
\end{observacion}

\section{Properties of the paths}
\label{sec:propiedades_brown}

In this section we study three pathwise properties of Brownian
motion: nowhere differentiability, unbounded variation, and the
explicit computation of the quadratic variation.

\subsection{Nowhere differentiability of the paths}
\label{subsec:no_diferenciabilidad}

The pathwise continuity of Brownian motion might suggest a picture of
smooth functions; the reality is the opposite: the paths are so
irregular that they admit a derivative at no point.

\medskip

\noindent\textbf{Heuristic argument.} Suppose, arguing by
contradiction, that the path of a Brownian motion were
differentiable at some time $t_0$. This would mean that the
difference quotient
\begin{equation*}
  \frac{W_{t_0 + h} - W_{t_0}}{h}
\end{equation*}
has a finite limit as $h \to 0$. But by condition (ii) of
Definition~\ref{def:browniano}, $W_{t_0+h} - W_{t_0}$ has distribution
$\mathcal{N}(0, h)$, and therefore the difference quotient has
distribution $\mathcal{N}(0,\, 1/h)$. Its variance diverges as
$h \to 0$: the quotient cannot converge to anything finite.

The argument is heuristic (we gloss over the difference between
convergence in distribution and almost sure convergence) but it
captures the underlying phenomenon. The rigorous proof requires
technical refinements, in particular the Borel-Cantelli lemma applied
to estimates on the increments at dyadic scales, and can be found in
\cite{karatzas-shreve}.

\begin{teorema}[Nowhere differentiability of Brownian paths]
\label{teo:no_diferenciabilidad}
Let $\{W_t\}_{t \geq 0}$ be a standard Brownian motion. With
probability one, the path $t \mapsto W_t(\omega)$ is differentiable at
no point $t \geq 0$.
\end{teorema}

Brownian paths are therefore continuous everywhere but differentiable
nowhere, a combination that will motivate the construction of the
\emph{ad hoc} calculus of Chapter~\ref{cap:itocalculus}.

\subsection{Unbounded variation of the paths}
\label{subsec:variacion_no_acotada}

We next prove a stronger quantitative property: Brownian paths have,
with probability one, \emph{unbounded variation} on every interval.

\medskip

\noindent\textbf{Total variation of a function.} Let us first recall
the classical notion.

\begin{definicion}[Total variation]
\label{def:variacion_total}
Let $f : [0, T] \to \R$ be a real function defined on a compact
interval. The \emph{total variation} of $f$ on $[0, T]$ is
\begin{equation*}
  V_T(f) \eqdef \sup_{\pi}\, \sum_{k=0}^{n-1} |f(t_{k+1}) - f(t_k)|,
\end{equation*}
where the supremum is taken over all partitions
$\pi = \{0 = t_0 < t_1 < \cdots < t_n = T\}$ of the interval. The
function $f$ is said to be of \emph{bounded variation} on $[0, T]$ if
$V_T(f) < \infty$.
\end{definicion}

Differentiable functions with bounded derivative are of bounded
variation (the variation is the integral of $|f'|$), as are monotone
functions. For a function of bounded variation, the
Riemann-Stieltjes integral $\int_0^T g\, df$ is unambiguously defined
for every continuous integrand $g$. It is precisely this circumstance
that fails in the case of Brownian motion.

\medskip

\noindent\textbf{Heuristic argument.} Let
$\pi_n = \{0 = t_0 < t_1 < \cdots < t_n = T\}$ be the uniform
partition of $[0, T]$ into $n$ subintervals of length $T/n$. By
condition (ii) of Definition~\ref{def:browniano}, each increment
$W_{t_{k+1}} - W_{t_k}$ has distribution $\mathcal{N}(0,\, T/n)$, and
therefore its absolute value satisfies
\begin{equation*}
  \E\bigl[|W_{t_{k+1}} - W_{t_k}|\bigr]
  = \sqrt{\frac{2T}{\pi\, n}}.
\end{equation*}
Summing over the $n$ subintervals,
\begin{equation*}
  \E\!\left[\,\sum_{k=0}^{n-1} |W_{t_{k+1}} - W_{t_k}|\right]
  = n\, \sqrt{\frac{2T}{\pi\, n}}
  = \sqrt{\frac{2 T n}{\pi}}.
\end{equation*}
The expectation diverges as $n \to \infty$: refining the partition
makes the sum of the absolute values of the increments grow without
bound, instead of approaching a finite value as would happen for a
function of bounded variation. The reasoning therefore suggests that
the Brownian path cannot be of bounded variation.

The above argument is heuristic (it proves divergence in expectation,
not almost surely) but it clearly displays the phenomenon. The
rigorous proof can be obtained, among other routes, as a corollary of
the computation of the quadratic variation that we address in the
next subsection (see the remark at the end of that subsection).

\begin{teorema}[Unbounded variation of Brownian paths]
\label{teo:variacion_no_acotada}
Let $\{W_t\}_{t \geq 0}$ be a standard Brownian motion and $T > 0$.
With probability one, the path $t \mapsto W_t(\omega)$ has unbounded
variation on $[0, T]$, that is, $V_T(W) = \infty$ almost surely.
\end{teorema}

\medskip

\noindent\textbf{Consequence: the failure of Riemann-Stieltjes.} The
unbounded variation of the paths has an immediate consequence for the
possibility of defining integrals against Brownian motion. The
classical Riemann-Stieltjes construction of $\int_0^T g\, df$ requires
$f$ to be of bounded variation in order for the sums
\begin{equation*}
  \sum_{k=0}^{n-1} g(\tau_k)\, \bigl(f(t_{k+1}) - f(t_k)\bigr)
\end{equation*}
to converge to the same limit regardless of the choice of the points
$\tau_k \in [t_k, t_{k+1}]$. When $f = W$, this result fails: the
limit depends on the choice of the $\tau_k$. Constructing an integral
$\int_0^T g_s\, dW_s$ will therefore require fixing a specific
convention for the evaluation points and verifying that the resulting
construction is well defined.

\subsection{Quadratic variation of Brownian motion}
\label{subsec:variacion_cuadratica}

Despite the irregularity of the paths, the \emph{sum of the squares}
of the increments exhibits a remarkably regular behavior: it converges
to a finite limit as the partition is refined. This quantity, the
\emph{quadratic variation}, is the central result of the chapter.

\medskip

\begin{teorema}[Quadratic variation of Brownian motion]
\label{teo:variacion_cuadratica}
Let $\{W_t\}_{t \geq 0}$ be a standard Brownian motion and $T > 0$.
For each partition $\pi = \{0 = t_0 < t_1 < \cdots < t_n = T\}$ of
$[0, T]$, denote its mesh by $|\pi| \eqdef \max_{k} (t_{k+1} - t_k)$
and consider the random variable
\begin{equation*}
  Q_\pi \eqdef \sum_{k=0}^{n-1} (W_{t_{k+1}} - W_{t_k})^2.
\end{equation*}
Then, as $|\pi| \to 0$,
\begin{equation}\label{eq:convergencia_var_cuadratica}
  Q_\pi \xrightarrow{L^2} T,
\end{equation}
that is, $\E\bigl[(Q_\pi - T)^2\bigr] \to 0$.
\end{teorema}

The quantity $Q_\pi$ is, by definition, a random variable (its value
depends on the path $\omega$). The theorem states that, in the limit
of fine partitions, this random variable ceases to be random: it
converges to a deterministic constant, $T$, equal to the length of the
interval. This is at first sight a surprising result, since one would
expect a sum of squares of Brownian increments (each with nonzero
variance) to retain some degree of randomness in the limit. The reason
this does not happen is that, as the partition is refined, the
contributions of the subintervals become ever more numerous and ever
smaller; the law of large numbers, applied to the sum, produces the
averaging to the constant.

\begin{proof}
The strategy consists in computing directly the expectation and the
variance of $Q_\pi$, and showing that in the limit the former tends to
$T$ and the latter to zero.

\emph{Step 1: computation of the expectation.} By condition (ii) of
Definition~\ref{def:browniano}, each increment $W_{t_{k+1}} - W_{t_k}$
has distribution $\mathcal{N}(0,\, t_{k+1} - t_k)$, and therefore
\begin{equation*}
  \E\bigl[(W_{t_{k+1}} - W_{t_k})^2\bigr] = t_{k+1} - t_k.
\end{equation*}
By linearity of expectation,
\begin{equation}\label{eq:esperanza_Q_pi}
  \E[Q_\pi] = \sum_{k=0}^{n-1} (t_{k+1} - t_k) = T.
\end{equation}
The expectation of $Q_\pi$ is exactly $T$, with no error: this result
is \emph{exact}, not asymptotic, and it holds for any partition.

\emph{Step 2: computation of the variance.} By definition,
\begin{equation*}
  \mathrm{Var}(Q_\pi) = \E\!\left[\bigl(Q_\pi - T\bigr)^2\right]
                      = \E\!\left[\biggl(\sum_{k=0}^{n-1}
                          \bigl[(W_{t_{k+1}} - W_{t_k})^2 - (t_{k+1} -
                          t_k)\bigr]\biggr)^{\!2}\right].
\end{equation*}
Expanding the square,
\begin{align*}
  \mathrm{Var}(Q_\pi)
  &= \sum_{k=0}^{n-1} \E\!\left[\bigl((W_{t_{k+1}} - W_{t_k})^2 -
                       (t_{k+1} - t_k)\bigr)^{\!2}\right] \\
  &\quad + 2\sum_{j < k} \E\!\Bigl[\bigl((W_{t_{j+1}} - W_{t_j})^2 -
                  (t_{j+1} - t_j)\bigr) \\
  &\qquad\qquad\qquad \cdot \bigl((W_{t_{k+1}} - W_{t_k})^2 -
                  (t_{k+1} - t_k)\bigr)\Bigr].
\end{align*}
The cross terms ($j < k$) vanish: by the independence of increments
over disjoint intervals (condition (iii) of
Definition~\ref{def:browniano}), the variables
$(W_{t_{j+1}} - W_{t_j})^2 - (t_{j+1} - t_j)$ and
$(W_{t_{k+1}} - W_{t_k})^2 - (t_{k+1} - t_k)$ are independent and,
moreover, both have zero expectation (by~\eqref{eq:esperanza_Q_pi}
applied to a single subinterval). Hence their product has zero
expectation, and only the diagonal terms survive:
\begin{equation*}
  \mathrm{Var}(Q_\pi)
  = \sum_{k=0}^{n-1} \E\!\left[\bigl((W_{t_{k+1}} - W_{t_k})^2 -
                       (t_{k+1} - t_k)\bigr)^{\!2}\right].
\end{equation*}
For each summand, write $\Delta_k \eqdef t_{k+1} - t_k$ and expand:
\begin{align*}
  \E\!\left[\bigl((W_{t_{k+1}} - W_{t_k})^2 - \Delta_k\bigr)^{\!2}\right]
  &= \E\bigl[(W_{t_{k+1}} - W_{t_k})^4\bigr]
   - 2\, \Delta_k\, \E\bigl[(W_{t_{k+1}} - W_{t_k})^2\bigr]
   + \Delta_k^2 \\
  &= 3\, \Delta_k^2 - 2\, \Delta_k^2 + \Delta_k^2 = 2\, \Delta_k^2,
\end{align*}
where we have used that the fourth moment of a
$\mathcal{N}(0, \Delta_k)$ is $3\, \Delta_k^2$.

Substituting,
\begin{equation}\label{eq:varianza_Q_pi}
  \mathrm{Var}(Q_\pi) = 2 \sum_{k=0}^{n-1} \Delta_k^2.
\end{equation}

\emph{Step 3: passage to the limit.} To bound the above sum, we
observe that
\begin{equation*}
  \sum_{k=0}^{n-1} \Delta_k^2
  \leq \max_k \Delta_k\, \sum_{k=0}^{n-1} \Delta_k
  = |\pi|\, T.
\end{equation*}
Consequently, $\mathrm{Var}(Q_\pi) \leq 2\, |\pi|\, T \to 0$ as
$|\pi| \to 0$. Combining with $\E[Q_\pi] = T$ from Step 1,
\begin{equation*}
  \E\bigl[(Q_\pi - T)^2\bigr] = \mathrm{Var}(Q_\pi) \to 0,
\end{equation*}
which is the $L^2$ convergence stated in
\eqref{eq:convergencia_var_cuadratica}.
\end{proof}

\medskip

\noindent\textbf{The heuristic rule $(dW_t)^2 = dt$.}
Theorem~\ref{teo:variacion_cuadratica} admits a particularly
illuminating informal interpretation. The statement that
$\sum_k (W_{t_{k+1}} - W_{t_k})^2 \to T$ as the partition is refined
can be read, formally, as
\begin{equation*}
  \int_0^T (dW_t)^2 = T,
\end{equation*}
or, in differential form, as the rule
\begin{equation}\label{eq:regla_dW2}
  (dW_t)^2 = dt.
\end{equation}
Identity~\eqref{eq:regla_dW2} is, obviously, an informal shorthand
that by itself has no precise mathematical meaning. But it captures a
fundamental fact: although the ``infinitesimal increment'' $dW_t$ is
of order $\sqrt{dt}$ (its standard deviation), its square is of order
$dt$, not $dt^2$ as it would be for a smooth function. In classical
calculus the terms quadratic in $dt$ are negligible; for Brownian
motion, they are not.

\medskip

\noindent\textbf{Consequence: unbounded variation.} We close the
section by returning to the claim announced in
Subsection~\ref{subsec:variacion_no_acotada}: the finite and nonzero
quadratic variation of Brownian motion implies, with probability one,
its unbounded variation. The argument is as follows. On any partition
$\pi$,
\begin{equation*}
  Q_\pi = \sum_{k=0}^{n-1} (W_{t_{k+1}} - W_{t_k})^2
        \leq \max_k |W_{t_{k+1}} - W_{t_k}|
              \cdot \sum_{k=0}^{n-1} |W_{t_{k+1}} - W_{t_k}|.
\end{equation*}
If the path were of bounded variation, the second factor would be
finite; by the uniform continuity of $W$ on $[0, T]$ (guaranteed by
condition (iv) of Definition~\ref{def:browniano}), the first factor
would tend to zero as $|\pi| \to 0$. The product would therefore tend
to zero, contradicting Theorem~\ref{teo:variacion_cuadratica}, which
ensures $Q_\pi \to T > 0$. Consequently, the path cannot be of
bounded variation.
\section{The Brownian filtration and continuous martingales}
\label{sec:filtracion_browniana}

We close the chapter with the filtration associated with Brownian
motion and the canonical examples of continuous martingales. The
section is brief: its content is essentially the adaptation to the
continuous setting of the discrete notions of
Chapter~\ref{cap:discreto}.

\subsection{The Brownian filtration and associated martingales}
\label{subsec:filtracion_y_martingalas}

\noindent\textbf{The natural filtration of Brownian motion.} The
filtration (Definition~\ref{def:filtracion}) generalizes to the
continuous setting by replacing the sequence $(\F_k)_{k=0}^N$ with a
family $(\F_t)_{t \geq 0}$ having the same monotonicity property
$\F_s \subseteq \F_t$ for $s \leq t$.

\begin{definicion}[Natural filtration of Brownian motion]
\label{def:filtracion_browniana}
Let $\{W_t\}_{t \geq 0}$ be a standard Brownian motion on
$(\Omega, \F, \PP)$. The \emph{natural filtration} of $W$ is the
family $(\F_t^W)_{t \geq 0}$ defined by
\begin{equation*}
  \F_t^W \eqdef \sigma\bigl(W_s : 0 \leq s \leq t\bigr),
\end{equation*}
that is, the $\sigma$-algebra generated by the values of the process
on the interval $[0, t]$.
\end{definicion}

The natural filtration encodes precisely the information generated by
Brownian motion, but it has two technical limitations that should be
corrected before using it as an operational framework. The first is
that it does not necessarily contain the $\PP$-null sets of $\F$: a
modification of an adapted process on a set of measure zero may fail
to be adapted, which prevents us from handling \emph{almost sure}
questions with the necessary flexibility. The second is that
\emph{right-continuity} (that is, the identity
$\F_t^W = \bigcap_{s > t} \F_s^W$) is not guaranteed, a property that
is indispensable for handling stopping times and for the existence of
continuous versions of processes such as the Itô integral
(Chapter~\ref{cap:itocalculus}).

The standard solution consists in passing to a slightly enlarged
filtration that does satisfy both properties.

\begin{definicion}[Augmented filtration of Brownian motion]
\label{def:filtracion_aumentada}
Let $\mathcal{N} \eqdef \{N \in \F : \PP(N) = 0\}$ be the collection
of $\PP$-null sets. The \emph{augmented filtration} of Brownian motion
is the family $(\F_t)_{t \geq 0}$ defined by
\begin{equation*}
  \F_t \eqdef \sigma\bigl(\F_t^W \cup \mathcal{N}\bigr),
\end{equation*}
that is, the smallest $\sigma$-algebra containing $\F_t^W$ and all
the $\PP$-null sets.
\end{definicion}

This definition resolves the first limitation by construction.
Moreover, it also resolves the second one automatically: for Brownian
motion, the augmented filtration by itself satisfies the
right-continuity property. The result can be found in
\cite{karatzas-shreve}, Proposition~2.7.7, and rests on Blumenthal's
zero-one law.

\begin{proposicion}[Right-continuity of the augmented filtration]
\label{prop:filtracion_continua_derecha}
The filtration $(\F_t)_{t \geq 0}$ of
Definition~\ref{def:filtracion_aumentada} is right-continuous:
$\F_t = \bigcap_{s > t} \F_s$ for all $t \geq 0$.
\end{proposicion}

Consequently, $(\F_t)$ satisfies the so-called \emph{usual
conditions}, completeness and right-continuity, the standard framework
on which the general theory of continuous processes is developed.
\textbf{From now on, and throughout the rest of this work, $(\F_t)$
will without exception denote the augmented filtration of Brownian
motion.} The distinction from the natural filtration $\F_t^W$ will not
be made explicit again, since for informational purposes both encode
the same events: any set in $\F_t$ differs from one in $\F_t^W$ by at
most a $\PP$-null set.

\medskip

\begin{definicion}[Continuous-time martingale]
\label{def:martingala_continua}
Let $(\Omega, \F, (\F_t)_{t \geq 0}, \PP)$ be a filtered probability
space. A stochastic process $\{M_t\}_{t \geq 0}$ is a
\emph{martingale} with respect to $(\F_t)$ if:
\begin{enumerate}[label=\textnormal{(\roman*)}]
  \item $M_t$ is $\F_t$-measurable for all $t \geq 0$ (\emph{adaptedness});
  \item $\E[|M_t|] < \infty$ for all $t \geq 0$ (\emph{integrability});
  \item $\E[M_t \mid \F_s] = M_s$ almost surely, for all
        $0 \leq s \leq t$.
\end{enumerate}
\end{definicion}

This is the literal translation to the continuous setting of
Definition~\ref{def:martingala}; the elementary properties
(linearity, expectation constant in time) are preserved, with proofs
identical to the discrete case.

\medskip

\noindent\textbf{Canonical examples.} We present two continuous
martingales associated with Brownian motion. The filtration
considered is the Brownian $(\F_t)$ fixed in
Definition~\ref{def:filtracion_aumentada}.

\begin{proposicion}[Brownian martingales]
\label{prop:martingalas_brownianas}
Let $\{W_t\}_{t \geq 0}$ be a standard Brownian motion. The following
processes are martingales with respect to $(\F_t)$:
\begin{enumerate}[label=\textnormal{(\alph*)}]
  \item Brownian motion itself, $\{W_t\}_{t \geq 0}$;
  \item for each $\lambda \in \R$, the \emph{exponential martingale}
        \begin{equation}\label{eq:martingala_exponencial}
          \mathcal{E}_t^{\lambda} \eqdef
          \exp\!\Bigl(\lambda\, W_t - \tfrac{1}{2}\lambda^2\, t\Bigr).
        \end{equation}
\end{enumerate}
\end{proposicion}

\begin{proof}
The adaptedness of both processes to $(\F_t)$ is immediate, since they
are functions of $W_t$. Integrability follows from the Gaussian
properties of Brownian motion. It remains to verify property (iii) of
Definition~\ref{def:martingala_continua}.

\emph{(a)} For $0 \leq s \leq t$, we write $W_t = W_s + (W_t - W_s)$,
and by condition (iii) of Definition~\ref{def:browniano}, the
increment $W_t - W_s$ is independent of $\F_s$. Consequently,
\begin{equation*}
  \E[W_t \mid \F_s] = W_s + \E[W_t - W_s \mid \F_s]
                    = W_s + \E[W_t - W_s] = W_s,
\end{equation*}
since $\E[W_t - W_s] = 0$ by condition (ii) of
Definition~\ref{def:browniano}.

\emph{(b)} For the exponential martingale,
\begin{equation*}
  \mathcal{E}_t^{\lambda}
  = \exp\!\bigl(\lambda W_s - \tfrac{1}{2}\lambda^2 s\bigr)
    \cdot \exp\!\bigl(\lambda(W_t - W_s) - \tfrac{1}{2}\lambda^2(t - s)\bigr)
  = \mathcal{E}_s^{\lambda} \cdot
    \exp\!\bigl(\lambda(W_t - W_s) - \tfrac{1}{2}\lambda^2(t - s)\bigr).
\end{equation*}
The first factor is $\F_s$-measurable and the second is independent of
$\F_s$. By the property ``taking out what is known'' and independence,
\begin{equation*}
  \E\bigl[\mathcal{E}_t^{\lambda} \mid \F_s\bigr]
  = \mathcal{E}_s^{\lambda}\,
    \E\!\left[\exp\!\bigl(\lambda(W_t - W_s)\bigr)\right]
    \cdot e^{-\frac{1}{2}\lambda^2(t-s)}.
\end{equation*}
Since $W_t - W_s \sim \mathcal{N}(0, t-s)$, its moment generating
function gives
$\E[\exp(\lambda(W_t-W_s))] = e^{\frac{1}{2}\lambda^2(t-s)}$, and the
exponential factor cancels exactly with the last term, so that
$\E[\mathcal{E}_t^{\lambda} \mid \F_s] = \mathcal{E}_s^{\lambda}$.
\end{proof}

\medskip

\noindent\textbf{Comment on the exponential martingale.} The
process~\eqref{eq:martingala_exponencial}, also known as the
\emph{stochastic exponential} of $W$, will play an important role in
Chapter~\ref{cap:itocalculus}: a generalized version of
formula~\eqref{eq:martingala_exponencial}, in which $\lambda$ is
replaced by an adapted process $\theta_t$, is precisely the
Radon-Nikodym derivative defining the change of measure in
Girsanov's theorem. The quadratic form of the exponent,
$-\tfrac{1}{2}\lambda^2 t$, is no accident: it is exactly the
correction that makes the exponential have unit expectation and,
therefore, that preserves the martingale property.

\medskip

Having completed the study of Brownian motion and its associated
martingales, we are in a position to develop the stochastic calculus
that will occupy the next chapter. The tools to be introduced there
(Itô integral, Itô's formula, Girsanov's theorem, Brownian martingale
representation theorem) rest essentially on the properties studied
so far.

% =====================================================================
%   CHAPTER 4 — STOCHASTIC CALCULUS
% =====================================================================
\chapter{Stochastic calculus}\label{cap:itocalculus}

This chapter develops the mathematical foundations of stochastic
calculus on which the continuous-time pricing theory of
Chapter~\ref{cap:bs} will be built. Its fundamental pillars are four:
the \emph{Itô stochastic integral}, which generalizes to the Brownian
setting the discrete stochastic integral of Chapter~\ref{cap:discreto};
\emph{Itô's formula}, which replaces the classical chain rule for
Brownian processes; \emph{Girsanov's theorem}, which formalizes the
change of measure in continuous time; and the \emph{Brownian
martingale representation theorem}, which is the continuous analogue
of the binomial representation theorem of
Section~\ref{subsec:representacion_binomial}. Each of these pillars
thus has its counterpart in the discrete case of
Chapter~\ref{cap:discreto}: the machinery we develop here introduces
no essentially new concepts, but translates already familiar ideas to
the continuous setting, where the irregularity of Brownian paths
(Chapter~\ref{cap:brown}) demands a considerably more careful
technical construction.

\section{Motivation: the need for a new integral}
\label{sec:motivacion_ito}

Before undertaking the construction of the stochastic integral, it is
essential to articulate precisely why an \emph{ad hoc} construction is
necessary and why the classical integration methods do not suffice.
The motivation unfolds in three steps: first, showing that the price
process of the underlying must necessarily exhibit unbounded variation
for the model to be economically sensible; second, checking that this
condition makes the classical Riemann-Stieltjes procedure break down;
and third, identifying the choice that will lead to the correct
definition.

\medskip

\noindent\textbf{The price of the underlying cannot have smooth
paths.} The first observation, anchored in the no-arbitrage
principle, is the following: in an arbitrage-free continuous model,
the paths of the discounted price of the underlying cannot be of
bounded variation. The argument, which we reproduce following
\cite{etheridge}, is constructive: we shall explicitly exhibit the
arbitrage opportunity that would otherwise arise. We work with the
discounted price $\widetilde S_t = e^{-rt} S_t$, on which the
self-financing condition takes a particularly simple form; the
variation properties of $S$ and $\widetilde S$ are equivalent under
$B_t = e^{rt}$.

The rigorous definition of a self-financing portfolio in continuous
time will be formalized in Chapter~\ref{cap:bs}; here we only need
one of its consequences, the direct analogue
of~\eqref{eq:integral_discreta} in the discrete case: the discounted
value at $T$ of a self-financing portfolio starting from initial
capital $V_0$ with an adapted strategy $\phi$ having continuous paths
is
\begin{equation}\label{eq:cartera_descontada_continua_BV}
  \widetilde V_T = V_0 + \int_0^T \phi_t\, d\widetilde S_t,
\end{equation}
where the integral, under the assumption of bounded variation of the
integrator, is interpreted pathwise as a classical
Riemann-Stieltjes integral.

Let $F : \R \to \R$ be of class $C^1$ and consider the strategy
$\phi_t = F'(\widetilde S_t)$, starting from $V_0 = F(\widetilde S_0)$.
By the fundamental theorem of calculus applied to the continuous
integrator of bounded variation $\widetilde S$,
\begin{equation*}
  \int_0^T F'(\widetilde S_t)\, d\widetilde S_t
  = F(\widetilde S_T) - F(\widetilde S_0),
\end{equation*}
and substituting in~\eqref{eq:cartera_descontada_continua_BV} we
obtain $\widetilde V_T = F(\widetilde S_T)$. Under the bounded
variation assumption, any quantity of the form $F(\widetilde S_T)$
with $F\in C^1$ is thus replicable from initial capital
$F(\widetilde S_0)$.

Let us specialize to $F(x) = (x - \widetilde S_0)^2$: the strategy
$\phi_t = 2\,(\widetilde S_t - \widetilde S_0)$, starting from
$V_0 = 0$, generates the nominal value at $T$
\begin{equation*}
  V_T = e^{rT}\,\bigl(\widetilde S_T - \widetilde S_0\bigr)^2,
\end{equation*}
which is nonnegative almost surely and strictly positive on the event
$\{\widetilde S_T \neq \widetilde S_0\}$, of strictly positive
probability in any nontrivial model. That is, it generates an
arbitrage opportunity starting from zero capital, in contradiction
with the fundamental assumption of the model.

The conclusion, therefore, is that any realistic model of the price
in continuous time must employ processes whose paths have unbounded
variation. The Brownian motion introduced in Chapter~\ref{cap:brown},
far from being a mathematical eccentricity, is exactly the kind of
process the theory demands.

\medskip

\noindent\textbf{Why Riemann-Stieltjes does not work.} Once we accept
that the integrator in an expression of the type
$\int_0^T \phi_s\, dW_s$ will have unbounded variation, the technical
problem arises immediately: the classical Riemann-Stieltjes
construction ceases to work. Let us briefly recall why.

In the Riemann-Stieltjes integral, given a partition
$\pi = \{0 = t_0 < t_1 < \cdots < t_n = T\}$ and points
$\tau_k \in [t_k, t_{k+1}]$, one considers the sum
\begin{equation*}
  S_\pi(\phi, f; \tau) \eqdef \sum_{k=0}^{n-1} \phi(\tau_k)\,
                              \bigl(f(t_{k+1}) - f(t_k)\bigr).
\end{equation*}
If $f$ is of bounded variation and $\phi$ is continuous, this sum
converges to the same limit as $|\pi| \to 0$, regardless of the
choice of the points $\tau_k$. This common limit is, by definition,
$\int_0^T \phi\, df$.

When $f$ has unbounded variation, on the other hand, the result
collapses: the limit, if it exists, depends on the choice of the
points $\tau_k$. The most transparent example is that of Brownian
motion itself with $\phi = W$, that is, the sums that would
approximate $\int_0^T W_s\, dW_s$. Consider the uniform partition
$t_k = kT/n$ and the left-endpoint choice ($\tau_k = t_k$), and
observe the algebraic identity
\begin{equation*}
  W_{t_k}(W_{t_{k+1}} - W_{t_k})
  = \tfrac{1}{2}\bigl(W_{t_{k+1}}^2 - W_{t_k}^2\bigr)
   - \tfrac{1}{2}(W_{t_{k+1}} - W_{t_k})^2.
\end{equation*}
Summing over $k = 0, 1, \ldots, n-1$, the first part telescopes and
yields $\tfrac{1}{2}(W_T^2 - W_0^2) = \tfrac{1}{2} W_T^2$, while the
second part tends, by Theorem~\ref{teo:variacion_cuadratica}, to
$\tfrac{1}{2} T$ in the sense of $L^2$ convergence. Consequently,
\begin{equation}\label{eq:suma_izq}
  \sum_{k=0}^{n-1} W_{t_k}(W_{t_{k+1}} - W_{t_k})
  \xrightarrow{L^2} \tfrac{1}{2}\bigl(W_T^2 - T\bigr).
\end{equation}
The analogous right-endpoint choice ($\tau_k = t_{k+1}$) produces, by
an entirely parallel computation in which the sign of the quadratic
term changes, the limit $\tfrac{1}{2}(W_T^2 + T)$. The two
approximating sums therefore tend to different limits: they differ
exactly by $T$, which is the quadratic variation of Brownian motion
on $[0, T]$. The cause is now evident: the difference between the two
sums is $\sum_k (W_{t_{k+1}} - W_{t_k})^2$, which for a function of
bounded variation would tend to zero, but which for Brownian motion
converges to $T$. The pathology that breaks Riemann-Stieltjes is
exactly the property that distinguishes Brownian motion from
classical processes.

\medskip

\noindent\textbf{Itô's choice.} The conclusion of the preceding
analysis is clear: to construct a well-defined integral
$\int_0^T \phi_s\, dW_s$, we cannot simply take the limit of the
partial sums, but must fix a specific convention for the evaluation
points. The choice we shall adopt, \emph{Itô's choice}, consists in
always evaluating the integrand at the \emph{left} endpoint of each
subinterval,
\begin{equation*}
  \sum_{k=0}^{n-1} \phi_{t_k}\, (W_{t_{k+1}} - W_{t_k}).
\end{equation*}
This choice is justified for two reasons. The first, mathematical in
nature, is that the resulting integral will preserve the martingale
property (Theorem~\ref{teo:martingala_integral_indefinida}). The
second, financial in nature, is that the left endpoint formalizes the
notion of \emph{predictability}: the quantity of stock held during
the period $[t_k, t_{k+1})$ is chosen at the beginning of the period,
at $t = t_k$, before knowing how the market will move during the
subinterval. It is exactly the same idea we articulated in the
discrete case in Section~\ref{subsec:predecibles_int_disc}, now
carried over to the continuous setting.

\section{Construction of the Itô stochastic integral}
\label{sec:construccion_ito}

Having presented the motivation, we turn to the rigorous construction
of the stochastic integral $\int_0^T \phi_s\, dW_s$. The general
strategy follows a standard scheme: we define the integral on a class
of elementary integrands where the object reduces to a finite sum; we
prove on that class a fundamental identity controlling the $L^2$ norm
of the result; and we extend the integral by density to a
sufficiently broad class of integrands, exploiting the completeness
of $L^2$. Throughout the section we fix a time horizon $T > 0$ and a
complete probability space $(\Omega, \F, \PP)$ on which a standard
Brownian motion $\{W_t\}_{t \in [0, T]}$ is considered together with
its augmented filtration $(\F_t)_{0 \le t \le T}$.

\subsection{Simple processes and the elementary integral}
\label{subsec:integral_simples}

We begin by defining the class of elementary integrands on which the
stochastic integral admits a direct definition, free of technical
difficulties.

\begin{definicion}[Simple predictable process]
\label{def:proceso_simple}
A stochastic process $\phi : [0, T] \times \Omega \to \R$ is
\emph{simple predictable}—or, for short, \emph{simple}—if there
exist a partition $0 = t_0 < t_1 < \cdots < t_n = T$ and random
variables $\xi_0, \xi_1, \ldots, \xi_{n-1}$ such that:
\begin{enumerate}[label=\textnormal{(\roman*)}]
  \item for each $i = 0, 1, \ldots, n-1$, $\xi_i$ is bounded and
        $\F_{t_i}$-measurable;
  \item the process $\phi$ admits the representation
        \begin{equation}\label{eq:proceso_simple}
          \phi_t(\omega) = \sum_{i=0}^{n-1} \xi_i(\omega)\,
                           \mathbf{1}_{(t_i, t_{i+1}]}(t).
        \end{equation}
\end{enumerate}
We denote by $\mathcal{S}$ the class of all simple processes on
$[0, T]$.
\end{definicion}

Let us pause briefly on condition (i). The measurability of $\xi_i$
with respect to the $\sigma$-algebra $\F_{t_i}$ (the information
available at the beginning of the subinterval $(t_i, t_{i+1}]$, not
at its end) is the exact translation to the continuous setting of the
notion of \emph{predictability}\footnote{The rigorous notion of
predictable process in continuous time is formalized via the
$\sigma$-algebra generated by the left-continuous adapted processes
on $[0, T] \times \Omega$, and coincides for the class of simple
processes with condition~(i). It will not be necessary to work with
it explicitly, since all the subsequent theory is extended by density
from $\mathcal{S}$; see \cite{karatzas-shreve}.} introduced in
Section~\ref{subsec:predecibles_int_disc} for the discrete case. The
value of the integrand during the subinterval is determined with the
information available at its beginning, before observing the
increment of Brownian motion.

\medskip

\begin{definicion}[Stochastic integral of a simple process]
\label{def:integral_simple}
Let $\phi \in \mathcal{S}$ be a simple process with the
representation \eqref{eq:proceso_simple}. Its \emph{stochastic
integral} (or \emph{Itô integral}) over $[0, T]$ is the random
variable
\begin{equation}\label{eq:integral_simple}
  I(\phi) \eqdef \int_0^T \phi_s\, dW_s
            \eqdef \sum_{i=0}^{n-1} \xi_i\,
                   \bigl(W_{t_{i+1}} - W_{t_i}\bigr).
\end{equation}
\end{definicion}

Formula~\eqref{eq:integral_simple} is the continuous version of the
discrete stochastic integral $(H \cdot M)_n$ defined
in~\eqref{eq:integral_discreta}: each increment of Brownian motion
over a subinterval is weighted by the value the integrand takes at
the beginning of that subinterval. Since the $\xi_i$ are random
variables, $I(\phi)$ is also a random variable, and by construction
it is $\F_T$-measurable.

\medskip

\noindent\textbf{Well-definedness and linearity.} Before proceeding,
it is worth noting that the integral~\eqref{eq:integral_simple} does
not depend on the particular representation used for $\phi$: refining
the partition,
$\xi_i\, (W_{t_{i+1}} - W_{t_i}) = \xi_i\, (W_s - W_{t_i}) + \xi_i\, (W_{t_{i+1}} - W_s)$,
a trivial identity by the additivity of Brownian increments, which
implies that any two representations (for which a common refinement,
their union, always exists) produce the same integral. Likewise, the
integral is linear: for $\phi, \psi \in \mathcal{S}$ and
$a, b \in \R$, $I(a\phi + b\psi) = a\, I(\phi) + b\, I(\psi)$ almost
surely, which follows immediately from
formula~\eqref{eq:integral_simple} after taking a common refinement
of the partitions.

\subsection{The Itô isometry}
\label{subsec:isometria_ito}

We next establish the \emph{Itô isometry}: an identity that controls
the $L^2$ norm of the integral in terms of the integrand.

\medskip

\begin{definicion}[The space $\mathcal{H}^2$]
\label{def:H2}
We denote by $\mathcal{H}^2[0, T]$, or simply $\mathcal{H}^2$ when
the time horizon is clear from the context, the space of stochastic
processes $\phi : [0, T] \times \Omega \to \R$ satisfying:
\begin{enumerate}[label=\textnormal{(\roman*)}]
   \item $\phi$ is progressively measurable\footnote{That is, for each
      $t \in [0, T]$, the restriction of $\phi$ to
      $[0, t] \times \Omega$ is $\mathcal{B}([0, t]) \otimes \F_t$-measurable.
      The property implies adaptedness and holds automatically for
      adapted processes with left-continuous paths, in particular for
      the simple processes of Definition~\ref{def:proceso_simple}.}
      with respect to the filtration $(\F_t)$;
  \item $\displaystyle \E\!\left[\int_0^T \phi_s^2\, ds\right] < \infty.$
\end{enumerate}
On $\mathcal{H}^2$ we consider the norm
\begin{equation}\label{eq:norma_H2}
  \|\phi\|_{\mathcal{H}^2}
  \eqdef \E\!\left[\int_0^T \phi_s^2\, ds\right]^{\!1/2}.
\end{equation}
\end{definicion}

The progressive measurability condition guarantees joint
measurability in $(t, \omega)$ and, in particular, that the integral
$\int_0^T \phi_s^2\, ds$ is well defined. For simple processes
(Definition~\ref{def:proceso_simple}) both conditions hold trivially,
and therefore $\mathcal{S} \subseteq \mathcal{H}^2$.

\medskip

\noindent\textbf{The isometry.}

\begin{teorema}[Itô isometry for simple processes]
\label{teo:isometria_ito_simples}
Let $\phi \in \mathcal{S}$ be a simple process. Then
\begin{equation}\label{eq:isometria_ito_simples}
  \E\!\left[\bigl(I(\phi)\bigr)^2\right]
  = \E\!\left[\int_0^T \phi_s^2\, ds\right]
  = \|\phi\|_{\mathcal{H}^2}^2.
\end{equation}
As an immediate consequence, $\E[I(\phi)] = 0$.
\end{teorema}

\begin{proof}
Let $\phi \in \mathcal{S}$ with the representation
\eqref{eq:proceso_simple}, that is,
$\phi_t = \sum_{i=0}^{n-1} \xi_i\, \mathbf{1}_{(t_i, t_{i+1}]}(t)$,
with $\xi_i$ bounded and $\F_{t_i}$-measurable. By
Definition~\ref{def:integral_simple},
$I(\phi) = \sum_{i=0}^{n-1} \xi_i\, (W_{t_{i+1}} - W_{t_i})$.

To lighten the notation, write $\Delta W_i \eqdef W_{t_{i+1}} - W_{t_i}$
and $\Delta_i \eqdef t_{i+1} - t_i$.

\emph{Step 1: computation of $\E[I(\phi)^2]$.} Expanding the square
of the sum,
\begin{equation*}
  I(\phi)^2 = \sum_{i=0}^{n-1} \xi_i^2\, (\Delta W_i)^2
            + 2 \sum_{j < i} \xi_j\, \xi_i\, \Delta W_j\, \Delta W_i.
\end{equation*}
Taking expectations,
\begin{equation}\label{eq:expansion_cuadrado}
  \E[I(\phi)^2]
  = \sum_{i=0}^{n-1} \E\bigl[\xi_i^2\, (\Delta W_i)^2\bigr]
  + 2 \sum_{j < i} \E\bigl[\xi_j\, \xi_i\, \Delta W_j\, \Delta W_i\bigr].
\end{equation}

\emph{Step 2: vanishing of the cross terms.} Let $j < i$. The
variables $\xi_j$, $\xi_i$ and $\Delta W_j$ are all
$\F_{t_i}$-measurable—indeed, $\xi_j$ is $\F_{t_j}$-measurable and
$\F_{t_j} \subseteq \F_{t_i}$, $\xi_i$ is $\F_{t_i}$-measurable by
assumption, and $\Delta W_j = W_{t_{j+1}} - W_{t_j}$ is
$\F_{t_i}$-measurable since $t_{j+1} \leq t_i$. By the property
``taking out what is known'' of conditional expectation
(Proposition~\ref{prop:propiedades_esperanza_cond}, property (iii)),
\begin{equation*}
  \E\bigl[\xi_j\, \xi_i\, \Delta W_j\, \Delta W_i \,\big|\, \F_{t_i}\bigr]
  = \xi_j\, \xi_i\, \Delta W_j\, \E[\Delta W_i \mid \F_{t_i}].
\end{equation*}
But $\Delta W_i$ is independent of $\F_{t_i}$ by condition (iii) of
Definition~\ref{def:browniano}, and has zero mean, hence
$\E[\Delta W_i \mid \F_{t_i}] = 0$. The conditional expectation of the
product is zero, and taking unconditional expectations, so is the
expectation:
\begin{equation*}
  \E\bigl[\xi_j\, \xi_i\, \Delta W_j\, \Delta W_i\bigr] = 0,
  \qquad \text{for all } j < i.
\end{equation*}
The cross terms of~\eqref{eq:expansion_cuadrado} vanish.

\emph{Step 3: computation of the diagonal terms.} For each $i$, the
variable $\xi_i^2$ is $\F_{t_i}$-measurable and $(\Delta W_i)^2$ is
independent of $\F_{t_i}$. Applying again ``taking out what is
known'' and independence,
\begin{equation*}
  \E\bigl[\xi_i^2\, (\Delta W_i)^2\bigr]
  = \E\bigl[\xi_i^2\, \E[(\Delta W_i)^2 \mid \F_{t_i}]\bigr]
  = \E[\xi_i^2]\, \E[(\Delta W_i)^2]
  = \E[\xi_i^2]\, \Delta_i,
\end{equation*}
where the last equality uses that
$\Delta W_i \sim \mathcal{N}(0, \Delta_i)$.

\emph{Step 4: identification with the $\mathcal{H}^2$ norm.} By
\eqref{eq:expansion_cuadrado} and the previous steps,
\begin{equation*}
  \E[I(\phi)^2] = \sum_{i=0}^{n-1} \E[\xi_i^2]\, \Delta_i.
\end{equation*}
On the other hand, on each subinterval $(t_i, t_{i+1}]$ the process
$\phi$ equals $\xi_i$, hence $\phi_s^2 = \xi_i^2$ and
\begin{equation*}
  \int_0^T \phi_s^2\, ds = \sum_{i=0}^{n-1} \xi_i^2\, \Delta_i.
\end{equation*}
Taking expectations,
\begin{equation*}
  \E\!\left[\int_0^T \phi_s^2\, ds\right]
  = \sum_{i=0}^{n-1} \E[\xi_i^2]\, \Delta_i,
\end{equation*}
which coincides with $\E[I(\phi)^2]$. Identity
\eqref{eq:isometria_ito_simples} is proved.

For the consequence $\E[I(\phi)] = 0$, it suffices to take
expectations in $I(\phi) = \sum_i \xi_i\, \Delta W_i$ and apply the
same argument of Step 2 to each individual summand:
$\E[\xi_i\, \Delta W_i] = \E[\xi_i\, \E[\Delta W_i \mid \F_{t_i}]] = 0$.
\end{proof}

\medskip

\noindent\textbf{Comment on the role of predictability.} The
preceding proof brings out with remarkable clarity the mathematical
reason why Itô's choice, evaluating the integrand at the left
endpoint of each subinterval, is the correct one. The decisive step
is \emph{Step 2}: the vanishing of the cross terms is obtained
precisely because $\xi_i$ is $\F_{t_i}$-measurable while the
increment $\Delta W_i$ is independent of $\F_{t_i}$. Had we chosen to
evaluate the integrand at the right endpoint, that is, if the
variables were $\F_{t_{i+1}}$-measurable, the independence of the
increment would no longer hold and the cross terms would no longer
vanish. The isometry is thus a property specific to the Itô integral,
not shared by integrals constructed with other conventions.

\medskip

\noindent\textbf{Reading as an isometry.} Identity
\eqref{eq:isometria_ito_simples} admits a natural geometric reading.
Regarding $\mathcal{S}$ as a subset of the Hilbert space
$\mathcal{H}^2$ with the norm $\|\cdot\|_{\mathcal{H}^2}$ defined
in~\eqref{eq:norma_H2}, and $L^2(\Omega, \F_T, \PP)$ as a Hilbert
space with its usual norm, identity~\eqref{eq:isometria_ito_simples}
says exactly that the operator
$I : (\mathcal{S},\, \|\cdot\|_{\mathcal{H}^2}) \to (L^2(\Omega),\, \|\cdot\|_{L^2})$
is linear and norm-preserving. This is exactly the notion of
\emph{linear isometry} in the sense of functional analysis.
\subsection{Extension to $\mathcal{H}^2$ by density}
\label{subsec:extension_H2}

The isometry established in Theorem~\ref{teo:isometria_ito_simples}
allows us to extend the integral from $\mathcal{S}$ to $\mathcal{H}^2$
by the classical density argument: we approximate the integrand by
simple processes and use the completeness of $L^2(\Omega)$ to pass to
the limit.

\medskip

\noindent\textbf{Density of simple processes.}

\begin{lema}[Density of $\mathcal{S}$ in $\mathcal{H}^2$]
\label{lema:densidad_simples}
The class $\mathcal{S}$ of simple processes is dense in
$\mathcal{H}^2$ with respect to the norm $\|\cdot\|_{\mathcal{H}^2}$.
That is, for every process $\phi \in \mathcal{H}^2$ there exists a
sequence $(\phi^{(n)}) \subset \mathcal{S}$ such that
\begin{equation*}
  \|\phi^{(n)} - \phi\|_{\mathcal{H}^2} \xrightarrow{n \to \infty} 0.
\end{equation*}
\end{lema}

The proof of the lemma, which we omit, proceeds in two stages. In the
first, the problem is reduced to the case of bounded processes
$\phi$, via truncation. In the second, a given bounded process is
approximated by constructing time regularizations of the form
\begin{equation*}
  \phi^{(n)}_t \eqdef \sum_{i=0}^{n-1}
                       \biggl(\frac{n}{T} \int_{(i-1) T/n}^{i T/n} \phi_s\,
                              ds\biggr)
                       \mathbf{1}_{(i T/n,\, (i+1) T/n]}(t),
\end{equation*}
which are simple processes and converge to $\phi$ in $\mathcal{H}^2$
by the continuity of translation in $L^2$. The detailed argument can
be found in \cite{karatzas-shreve}, Section 3.2.

\medskip

\noindent\textbf{Construction of the integral on $\mathcal{H}^2$.}
Let $\phi \in \mathcal{H}^2$. By Lemma~\ref{lema:densidad_simples}
there exists a sequence $(\phi^{(n)}) \subset \mathcal{S}$ with
$\|\phi^{(n)} - \phi\|_{\mathcal{H}^2} \to 0$. In particular,
$(\phi^{(n)})$ is Cauchy in $\mathcal{H}^2$. Applying the Itô
isometry (Theorem~\ref{teo:isometria_ito_simples}) and linearity to
the difference $\phi^{(n)} - \phi^{(m)} \in \mathcal{S}$,
\begin{equation}\label{eq:cauchy_integral}
  \E\!\left[\bigl(I(\phi^{(n)}) - I(\phi^{(m)})\bigr)^{\!2}\right]
  = \E\!\left[\bigl(I(\phi^{(n)} - \phi^{(m)})\bigr)^{\!2}\right]
  = \|\phi^{(n)} - \phi^{(m)}\|_{\mathcal{H}^2}^2 \xrightarrow{n,m \to \infty} 0.
\end{equation}
The sequence $(I(\phi^{(n)}))$ is therefore Cauchy in
$L^2(\Omega, \F_T, \PP)$. Since this space is complete, there exists
a unique element $L \in L^2(\Omega)$ such that $I(\phi^{(n)}) \to L$
in $L^2$.

It remains to verify that the limit $L$ does not depend on the
particular approximating sequence. Let
$(\phi^{(n)}), (\widetilde{\phi}^{(n)}) \subset \mathcal{S}$ be two
sequences both converging to $\phi$ in $\mathcal{H}^2$, with
respective limits $L, \widetilde{L} \in L^2(\Omega)$. Considering the
interleaved sequence
$\phi^{(1)}, \widetilde{\phi}^{(1)}, \phi^{(2)}, \widetilde{\phi}^{(2)}, \ldots$,
which also converges to $\phi$ in $\mathcal{H}^2$, the images
$I(\phi^{(1)}), I(\widetilde{\phi}^{(1)}), I(\phi^{(2)}), \ldots$ form
a Cauchy sequence in $L^2(\Omega)$ by the same argument
as~\eqref{eq:cauchy_integral}, and hence converge to a unique limit.
But the two subsequences $(I(\phi^{(n)}))$ and
$(I(\widetilde{\phi}^{(n)}))$ must converge to that same limit, which
forces $L = \widetilde{L}$ almost surely.

\begin{definicion}[Stochastic integral on $\mathcal{H}^2$]
\label{def:integral_H2}
Let $\phi \in \mathcal{H}^2$ and let $(\phi^{(n)}) \subset \mathcal{S}$
be any sequence of simple processes with
$\|\phi^{(n)} - \phi\|_{\mathcal{H}^2} \to 0$. The \emph{Itô
stochastic integral} of $\phi$ over $[0, T]$ is the random variable
\begin{equation*}
  I(\phi) \eqdef \int_0^T \phi_s\, dW_s
            \eqdef L^2\text{-}\!\lim_{n \to \infty}\, I(\phi^{(n)}),
\end{equation*}
an element of $L^2(\Omega, \F_T, \PP)$ defined up to modifications on
sets of measure zero.
\end{definicion}

The definition is consistent: for $\phi \in \mathcal{S}$, taking the
constant sequence $\phi^{(n)} = \phi$ recovers
Definition~\ref{def:integral_simple}.

\begin{teorema}[Properties of the stochastic integral on $\mathcal{H}^2$]
\label{teo:propiedades_integral_H2}
The operator $I : \mathcal{H}^2 \to L^2(\Omega, \F_T, \PP)$ given by
Definition~\ref{def:integral_H2} satisfies the following properties:
\begin{enumerate}[label=\textnormal{(\roman*)}]
  \item \emph{Linearity:} for $\phi, \psi \in \mathcal{H}^2$ and
        $a, b \in \R$,
        \begin{equation*}
          I(a\phi + b\psi) = a\, I(\phi) + b\, I(\psi)
          \quad \text{a.s.}
        \end{equation*}
  \item \emph{Itô isometry:} for $\phi \in \mathcal{H}^2$,
        \begin{equation}\label{eq:isometria_H2}
          \E\!\left[\bigl(I(\phi)\bigr)^{\!2}\right]
          = \E\!\left[\int_0^T \phi_s^2\, ds\right]
          = \|\phi\|_{\mathcal{H}^2}^2.
        \end{equation}
  \item \emph{Zero expectation:} for $\phi \in \mathcal{H}^2$,
        $\E[I(\phi)] = 0$.
\end{enumerate}
\end{teorema}

\begin{proof}
The three properties follow from the corresponding properties on
$\mathcal{S}$ by passing to the limit, thanks to the fact that $I$ is
continuous with respect to the norms $\|\cdot\|_{\mathcal{H}^2}$ and
$\|\cdot\|_{L^2}$ (a continuity which is in turn a consequence of the
isometry on $\mathcal{S}$).

By way of illustration, we prove (ii). Let $\phi \in \mathcal{H}^2$
and $(\phi^{(n)}) \subset \mathcal{S}$ with $\phi^{(n)} \to \phi$ in
$\mathcal{H}^2$. Then $I(\phi^{(n)}) \to I(\phi)$ in $L^2(\Omega)$ by
Definition~\ref{def:integral_H2}. The continuity of the norm
$\|\cdot\|_{L^2}$ gives
\begin{equation*}
  \E[I(\phi)^2]
  = \|I(\phi)\|_{L^2}^2
  = \lim_{n\to\infty} \|I(\phi^{(n)})\|_{L^2}^2
  = \lim_{n\to\infty} \E[I(\phi^{(n)})^2].
\end{equation*}
By Theorem~\ref{teo:isometria_ito_simples} applied to each
$\phi^{(n)} \in \mathcal{S}$,
\begin{equation*}
  \E[I(\phi^{(n)})^2] = \|\phi^{(n)}\|_{\mathcal{H}^2}^2,
\end{equation*}
and by the continuity of the norm $\|\cdot\|_{\mathcal{H}^2}$,
$\|\phi^{(n)}\|_{\mathcal{H}^2}^2 \to \|\phi\|_{\mathcal{H}^2}^2$.
Identity~\eqref{eq:isometria_H2} is obtained by passing to the limit.
The proofs of (i) and (iii) are entirely analogous.
\end{proof}

The integral so defined is, however, a \emph{global} object on
$[0, T]$. In many contexts we shall need the integral over $[0, t]$
with $t$ variable: this motivates the study of the process $\{M_t\}$
in the next subsection.

\subsection{The indefinite integral as a process}
\label{subsec:integral_indefinida}

For each $t \in [0, T]$, the truncated process
$\phi\, \mathbf{1}_{[0, t]}$ belongs to $\mathcal{H}^2$ (the
restriction of a square-integrable process is square-integrable),
which allows us to apply the integral operator $I$ to it and thus
obtain a family $\{M_t\}_{t \in [0, T]}$ that we shall study as a
process indexed by $t$.

\begin{definicion}[Indefinite integral]
\label{def:integral_indefinida}
Let $\phi \in \mathcal{H}^2$. For each $t \in [0, T]$, we define
\begin{equation}\label{eq:def_integral_indefinida}
  M_t \eqdef \int_0^t \phi_s\, dW_s
        \eqdef I\!\left(\phi\, \mathbf{1}_{[0, t]}\right).
\end{equation}
The family $\{M_t\}_{t \in [0, T]}$ so obtained is called the
\emph{indefinite integral} of $\phi$ with respect to the Brownian
motion $W$.
\end{definicion}

Note that for $t = T$ we recover $M_T = I(\phi)$, the global operator
of the previous subsection. Likewise,
$M_0 = I(\phi\, \mathbf{1}_{\{0\}}) = 0$ almost surely, since the
process $\phi\, \mathbf{1}_{\{0\}}$ has zero $\mathcal{H}^2$ norm.

\medskip

\noindent\textbf{Continuity of the paths.} As we have constructed the
process, each $M_t$ is defined as an element of $L^2(\Omega)$, that
is, up to modifications on sets of measure zero. This form of
definition does not by itself guarantee that the paths
$t \mapsto M_t(\omega)$ are continuous functions: the family
$\{M_t\}_{t \in [0, T]}$ is indexed by a continuous parameter, and
the uncountable union of the null sets associated with each $M_t$
might, \emph{a priori}, fail to be a null set. It is therefore
necessary to verify the existence of a continuous \emph{version} of
the process, that is, a choice of representatives in each equivalence
class that yields continuous paths almost surely.

\begin{teorema}[Continuity of the indefinite integral]
\label{teo:continuidad_integral_indefinida}
Let $\phi \in \mathcal{H}^2$ and let $\{M_t\}_{t \in [0, T]}$ be the
process defined in~\eqref{eq:def_integral_indefinida}. There exists a
version of this process whose paths are continuous almost surely.
\end{teorema}

The rigorous proof, which we omit, rests essentially on Doob's
maximal inequality for $L^2$ martingales, together with the
approximation argument by simple processes already employed. For
simple processes, the paths of the process $M_t$ are continuous by
construction (they are piecewise linear in $t$); the passage to the
limit by density, controlled by means of the maximal inequality
applied to the difference of two approximations, guarantees almost
sure uniform convergence on $[0, T]$, which preserves continuity. The
detailed construction can be found in \cite{karatzas-shreve},
Section 3.2.

From here on, we shall always identify $\{M_t\}$ with its continuous
version, without explicit mention.

\medskip

\noindent\textbf{The martingale property.}

\begin{teorema}[Martingale property of the indefinite integral]
\label{teo:martingala_integral_indefinida}
Let $\phi \in \mathcal{H}^2$ and let $\{M_t\}_{t \in [0, T]}$ be the
indefinite integral of $\phi$. Then $\{M_t\}$ is a martingale with
respect to the filtration $(\F_t)$, that is:
\begin{enumerate}[label=\textnormal{(\roman*)}]
  \item $M_t$ is $\F_t$-measurable for all $t \in [0, T]$;
  \item $\E[|M_t|] < \infty$ for all $t \in [0, T]$;
  \item $\E[M_t \mid \F_s] = M_s$ almost surely, for all
        $0 \leq s \leq t \leq T$.
\end{enumerate}
\end{teorema}

\begin{proof}
\emph{Adaptedness and integrability.} For simple processes, $M_t$ is
a finite sum of products of $\F_t$-measurable variables and hence
$\F_t$-measurable. For general $\phi \in \mathcal{H}^2$, the $L^2$
convergence of the approximations $M_t^{(n)}$ to $M_t$
(Definition~\ref{def:integral_H2}) yields an almost surely convergent
subsequence, and $\F_t$-measurability is preserved under almost sure
limits. Integrability follows from Cauchy-Schwarz:
$\E[|M_t|] \leq \E[M_t^2]^{1/2} = \|\phi\, \mathbf{1}_{[0, t]}\|_{\mathcal{H}^2} \leq \|\phi\|_{\mathcal{H}^2} < \infty$,
where the middle equality uses the Itô isometry.

\emph{Martingale property.} Let $0 \leq s \leq t \leq T$. We begin
with the case of a simple process $\phi \in \mathcal{S}$ with
representation $\phi_r = \sum_i \xi_i\, \mathbf{1}_{(t_i, t_{i+1}]}(r)$.
Taking a refinement of the partition that includes the points $s$ and
$t$—which alters neither the process $\phi$ nor the indefinite
integral—we may assume without loss of generality that $s = t_p$
and $t = t_q$ for some $p \leq q$. Then
\begin{equation*}
  M_t - M_s = \int_s^t \phi_r\, dW_r
            = \sum_{i = p}^{q-1} \xi_i\,
              \bigl(W_{t_{i+1}} - W_{t_i}\bigr).
\end{equation*}
Taking conditional expectations with respect to $\F_s = \F_{t_p}$ and
using linearity,
\begin{equation*}
  \E[M_t - M_s \mid \F_s]
  = \sum_{i = p}^{q-1} \E\bigl[\xi_i\, (W_{t_{i+1}} - W_{t_i})
                               \,\big|\, \F_{t_p}\bigr].
\end{equation*}
For each summand, observe that $\F_{t_p} \subseteq \F_{t_i}$ since
$i \geq p$. Applying the tower property
(Proposition~\ref{prop:propiedades_esperanza_cond}, property (iv))
with $\F_{t_p} \subseteq \F_{t_i}$,
\begin{equation*}
  \E\bigl[\xi_i\, (W_{t_{i+1}} - W_{t_i}) \mid \F_{t_p}\bigr]
  = \E\!\left[\E\bigl[\xi_i\, (W_{t_{i+1}} - W_{t_i})
                       \mid \F_{t_i}\bigr] \,\Big|\, \F_{t_p}\right].
\end{equation*}
The inner conditional expectation is zero: $\xi_i$ is
$\F_{t_i}$-measurable and $W_{t_{i+1}} - W_{t_i}$ is independent of
$\F_{t_i}$ with zero mean, hence taking out what is known and using
independence,
$\E[\xi_i\, (W_{t_{i+1}} - W_{t_i}) \mid \F_{t_i}] = \xi_i\, \E[W_{t_{i+1}} - W_{t_i} \mid \F_{t_i}] = 0$.
Consequently, $\E[M_t - M_s \mid \F_s] = 0$, that is,
$\E[M_t \mid \F_s] = M_s$ almost surely. The martingale property is
established for simple processes.

For general $\phi \in \mathcal{H}^2$, take a sequence
$(\phi^{(n)}) \subset \mathcal{S}$ with $\phi^{(n)} \to \phi$ in
$\mathcal{H}^2$, and let $M_t^{(n)}$ be the corresponding indefinite
integrals. By what has just been proved,
$\E[M_t^{(n)} \mid \F_s] = M_s^{(n)}$ almost surely. Taking the limit
in $L^2(\Omega)$ as $n \to \infty$, $M_t^{(n)} \to M_t$ and
$M_s^{(n)} \to M_s$ by the isometry applied to the restrictions to
$[0, t]$ and $[0, s]$ respectively. Since conditional expectation
with respect to $\F_s$ is a continuous operator on $L^2$—it is the
orthogonal projection onto the subspace $L^2(\Omega, \F_s, \PP)$—the
equality passes to the limit: $\E[M_t \mid \F_s] = M_s$ almost
surely.
\end{proof}

\medskip

\noindent\textbf{The isometry in process form.} As an immediate
consequence of the construction, the Itô isometry generalizes to the
indefinite process: for all $t \in [0, T]$,
\begin{equation}\label{eq:isometria_proceso}
  \E[M_t^2] = \E\!\left[\int_0^t \phi_s^2\, ds\right].
\end{equation}
The identity follows by applying
Theorem~\ref{teo:propiedades_integral_H2} to the restricted process
$\phi\, \mathbf{1}_{[0, t]} \in \mathcal{H}^2$.
\section{Itô's formula}
\label{sec:formula_ito}

This section develops \emph{Itô's formula}, the continuous-time
substitute for the classical chain rule and, together with Girsanov's
theorem, one of the fundamental pieces of Chapter~\ref{cap:bs}.

\subsection{Itô processes and Itô's formula for $f(W_t)$}
\label{subsec:procesos_ito_y_formula}

\begin{definicion}[Itô process]
\label{def:proceso_ito}
A stochastic process $\{X_t\}_{t \in [0, T]}$ is an \emph{Itô
process} if it admits a representation of the form
\begin{equation}\label{eq:def_proceso_ito}
  X_t = X_0 + \int_0^t a_s\, ds + \int_0^t b_s\, dW_s,
  \qquad t \in [0, T],
\end{equation}
where $X_0$ is an $\F_0$-measurable random variable, and the
processes $\{a_t\}$ and $\{b_t\}$ are adapted to $(\F_t)$ and satisfy
the integrability conditions
\begin{equation*}
  \int_0^T |a_s|\, ds < \infty \text{ almost surely},
  \qquad b \in \mathcal{H}^2[0, T].
\end{equation*}
The processes $a$ and $b$ are called, respectively, the \emph{drift
coefficient} and the \emph{diffusion coefficient} of the process $X$.
\end{definicion}

The condition $b \in \mathcal{H}^2[0, T]$ guarantees that the
stochastic integral $\int_0^t b_s\, dW_s$ is well defined in the
sense of Theorem~\ref{teo:propiedades_integral_H2}, while the
condition on $a$ ensures that the integral $\int_0^t a_s\, ds$ is
well defined pathwise as a Lebesgue integral.

\medskip

\noindent\textbf{Differential notation.} The integral
equality~\eqref{eq:def_proceso_ito} is usually written, in
abbreviated form, as
\begin{equation}\label{eq:notacion_diferencial}
  dX_t = a_t\, dt + b_t\, dW_t.
\end{equation}
This notation is very convenient and will be used frequently from now
on, but it should be stressed that it has no mathematical meaning of
its own: it is nothing more than a compact way of expressing
\eqref{eq:def_proceso_ito}. Identities of the
form~\eqref{eq:notacion_diferencial} must always be understood as
shorthand for the corresponding integral identities.

\medskip

\noindent\textbf{Statement of Itô's formula.}

\begin{teorema}[Itô's formula for $f(W_t)$]
\label{teo:formula_ito_brown}
Let $f : \R \to \R$ be a function of class $C^2$ with $f$, $f'$ and
$f''$ bounded. Then, for all $t \in [0, T]$, almost surely,
\begin{equation}\label{eq:formula_ito_brown}
  f(W_t) = f(0) + \int_0^t f'(W_s)\, dW_s
        + \frac{1}{2} \int_0^t f''(W_s)\, ds.
\end{equation}
Equivalently, in differential notation,
\begin{equation*}
  d\bigl(f(W_t)\bigr) = f'(W_t)\, dW_t + \tfrac{1}{2} f''(W_t)\, dt.
\end{equation*}
\end{teorema}

\noindent\textbf{Reading of the result.} The differential of $f(W_t)$
is the sum of a term in $dW_t$, analogous to that of the classical
chain rule, and an extra term in $dt$, the \emph{Itô term}, absent in
ordinary calculus. This correction is a consequence of the heuristic
rule $(dW_t)^2 = dt$ anticipated in~\eqref{eq:regla_dW2}: when $f$ is
expanded in a Taylor series, the second-order term does not vanish
and contributes exactly $\tfrac{1}{2} f''(W_t)\, dt$. The boundedness
assumptions are the technically most convenient ones; the formula
holds under weaker conditions via localization.

\begin{proof}
For each partition $\pi = \{0 = t_0 < t_1 < \cdots < t_n = t\}$ of
$[0, t]$, we write the telescoping expansion
\begin{equation}\label{eq:telescopica}
  f(W_t) - f(0)
  = \sum_{k=0}^{n-1} \bigl[f(W_{t_{k+1}}) - f(W_{t_k})\bigr];
\end{equation}
the resulting terms will produce the two integrals on the right-hand
side of~\eqref{eq:formula_ito_brown}, plus a negligible remainder.

\emph{Step 1: Taylor expansion.} By the assumption $f \in C^2$, for
each $k$ there exists $\theta_k \in [W_{t_k}, W_{t_{k+1}}]$ (or the
corresponding interval with reversed endpoints) such that
\begin{equation*}
  f(W_{t_{k+1}}) - f(W_{t_k})
  = f'(W_{t_k})\, \Delta W_k
   + \tfrac{1}{2} f''(W_{t_k})\, (\Delta W_k)^2
   + R_k,
\end{equation*}
where $\Delta W_k \eqdef W_{t_{k+1}} - W_{t_k}$ and the remainder is
\begin{equation*}
  R_k = \tfrac{1}{2}\bigl[f''(\theta_k) - f''(W_{t_k})\bigr]\,
        (\Delta W_k)^2.
\end{equation*}
Summing over $k$ and comparing with~\eqref{eq:telescopica},
\begin{equation}\label{eq:descomposicion_telescopica}
  f(W_t) - f(0)
  = \underbrace{\sum_{k=0}^{n-1} f'(W_{t_k})\, \Delta W_k}_{S_1^\pi}
   + \tfrac{1}{2} \underbrace{\sum_{k=0}^{n-1} f''(W_{t_k})\,
                              (\Delta W_k)^2}_{S_2^\pi}
   + \underbrace{\sum_{k=0}^{n-1} R_k}_{R^\pi}.
\end{equation}

\emph{Step 2: convergence of $S_1^\pi$.} The sum $S_1^\pi$ is exactly
the stochastic integral of the simple process
$\phi^\pi_s \eqdef \sum_k f'(W_{t_k})\, \mathbf{1}_{(t_k, t_{k+1}]}(s)$
with respect to $W$ over $[0, t]$. Let us show that $\phi^\pi$
converges to $f'(W_s)$ in $\mathcal{H}^2[0, t]$ as $|\pi| \to 0$. By
the definition of the norm,
\begin{equation*}
  \|\phi^\pi - f'(W)\|_{\mathcal{H}^2}^2
  = \E\!\left[\int_0^t \bigl(\phi^\pi_s - f'(W_s)\bigr)^{\!2}\, ds\right]
  = \E\!\left[\sum_{k=0}^{n-1} \int_{t_k}^{t_{k+1}}
             \bigl(f'(W_{t_k}) - f'(W_s)\bigr)^{\!2}\, ds\right].
\end{equation*}
By the uniform continuity of $f'$ on any compact set and the almost
sure continuity of Brownian paths, the pointwise integrand tends to
zero as $|\pi| \to 0$; by the boundedness of $f'$, it is dominated by
a constant. The dominated convergence theorem guarantees that the
integral converges to zero, that is, $\phi^\pi \to f'(W)$ in
$\mathcal{H}^2$. By the Itô isometry,
\begin{equation*}
  S_1^\pi \xrightarrow{L^2} \int_0^t f'(W_s)\, dW_s.
\end{equation*}

\emph{Step 3: convergence of $S_2^\pi$.} We shall prove that
\begin{equation*}
  S_2^\pi \xrightarrow{L^2} \int_0^t f''(W_s)\, ds
  \quad \text{as } |\pi| \to 0.
\end{equation*}
To this end, we decompose
\begin{equation*}
  S_2^\pi - \int_0^t f''(W_s)\, ds
  = \underbrace{\sum_{k=0}^{n-1} f''(W_{t_k})\,
                \bigl[(\Delta W_k)^2 - \Delta_k\bigr]}_{A^\pi}
  + \underbrace{\sum_{k=0}^{n-1} f''(W_{t_k})\, \Delta_k
              - \int_0^t f''(W_s)\, ds}_{B^\pi},
\end{equation*}
where $\Delta_k \eqdef t_{k+1} - t_k$, and we shall prove that both
$A^\pi$ and $B^\pi$ tend to zero in $L^2$.

For $B^\pi$: the quantity $\sum_k f''(W_{t_k})\, \Delta_k$ is the
Riemann sum of the continuous function $s \mapsto f''(W_s)$
associated with the partition $\pi$ and the choice $\tau_k = t_k$. By
the almost sure pathwise continuity, it converges pathwise to
$\int_0^t f''(W_s)\, ds$ as $|\pi| \to 0$, and the uniform bound on
$f''$ allows us to apply dominated convergence to obtain convergence
in $L^2$.

For $A^\pi$: by an argument analogous to that of the Itô isometry
(Theorem~\ref{teo:isometria_ito_simples}, Step 2), using that
$f''(W_{t_k})$ is $\F_{t_k}$-measurable and that
$(\Delta W_k)^2 - \Delta_k$ has zero conditional expectation given
$\F_{t_k}$, the cross terms vanish when squaring and taking
expectations. Applying the uniform bound $|f''| \leq M$ and the
Gaussian fourth-moment computation
$\E[((\Delta W_k)^2 - \Delta_k)^2] = 2\, \Delta_k^2$,
\begin{equation*}
  \E[(A^\pi)^2]
  = \sum_{k=0}^{n-1} \E\bigl[(f''(W_{t_k}))^2\bigr]\,
    \E\bigl[((\Delta W_k)^2 - \Delta_k)^2\bigr]
  \leq 2 M^2 \sum_{k=0}^{n-1} \Delta_k^2
  \leq 2 M^2 \cdot |\pi| \cdot t \xrightarrow{|\pi| \to 0} 0.
\end{equation*}

\emph{Step 4: the remainder $R^\pi$ is negligible.} Recall that
\begin{equation*}
  R^\pi = \tfrac{1}{2} \sum_{k=0}^{n-1}
            \bigl[f''(\theta_k) - f''(W_{t_k})\bigr]\, (\Delta W_k)^2,
\end{equation*}
with $\theta_k$ between $W_{t_k}$ and $W_{t_{k+1}}$. Since $f''$ is
uniformly continuous on any compact set and Brownian paths are
continuous almost surely, the factor $|f''(\theta_k) - f''(W_{t_k})|$
tends uniformly to zero as $|\pi| \to 0$, that is,
\begin{equation*}
  \varepsilon^\pi \eqdef \max_{0 \leq k \leq n-1}
                          |f''(\theta_k) - f''(W_{t_k})|
                          \xrightarrow{|\pi| \to 0} 0
  \quad \text{almost surely}.
\end{equation*}
Therefore,
\begin{equation*}
  |R^\pi| \leq \tfrac{1}{2}\, \varepsilon^\pi
              \sum_{k=0}^{n-1} (\Delta W_k)^2.
\end{equation*}
The sum on the right-hand side converges to $t$ in $L^2$ by
Theorem~\ref{teo:variacion_cuadratica}, and $\varepsilon^\pi \to 0$
almost surely with uniform bound $2M$. Combining both facts (via
dominated convergence) we conclude that $R^\pi \to 0$ in $L^1$, and
hence in probability. Passing to a subsequence if necessary,
$R^\pi \to 0$ almost surely.

\emph{Step 5: conclusion.} Returning to the
decomposition~\eqref{eq:descomposicion_telescopica} and combining the
limits obtained in the previous steps (after passing to a common
subsequence ensuring that all convergences hold almost surely),
\begin{equation*}
  f(W_t) - f(0) = \int_0^t f'(W_s)\, dW_s
                + \tfrac{1}{2} \int_0^t f''(W_s)\, ds
                + 0,
\end{equation*}
almost surely, which is identity~\eqref{eq:formula_ito_brown}.
\end{proof}

\medskip

\subsection{General version and examples of application}
\label{subsec:formula_ito_general}

\medskip

\begin{teorema}[Itô's formula for Itô processes]
\label{teo:formula_ito_general}
Let $\{X_t\}$ be an Itô process with dynamics
$dX_t = a_t\, dt + b_t\, dW_t$, and let $f : [0, T] \times \R \to \R$
be a function of class $C^{1,2}$, that is, continuously
differentiable in its first variable and twice continuously
differentiable in the second, with the relevant derivatives bounded
on the range of $X$. Then the process $\{f(t, X_t)\}_{t \in [0, T]}$
is also an Itô process, and it satisfies
\begin{equation}\label{eq:formula_ito_general}
  d\bigl(f(t, X_t)\bigr)
  = \biggl[\partial_t f(t, X_t) + a_t\, \partial_x f(t, X_t)
           + \tfrac{1}{2}\, b_t^2\, \partial_{xx} f(t, X_t)\biggr]\, dt
   + b_t\, \partial_x f(t, X_t)\, dW_t.
\end{equation}
Equivalently, in integral form,
\begin{multline*}
  f(t, X_t) = f(0, X_0)
      + \int_0^t \!\Bigl[\partial_t f(s, X_s) + a_s\, \partial_x f(s, X_s)
                          + \tfrac{1}{2}\, b_s^2\, \partial_{xx} f(s, X_s)\Bigr]
                          ds \\
      + \int_0^t b_s\, \partial_x f(s, X_s)\, dW_s.
\end{multline*}
\end{teorema}

The proof of this theorem, which we omit, follows the same strategy
as that of Theorem~\ref{teo:formula_ito_brown}: a telescoping
expansion followed by a second-order Taylor expansion in the space
variable, now with an additional first-order term in the time
variable. The identification of the limits is entirely analogous; the
details can be found in \cite{karatzas-shreve}, Section 3.3.

\medskip

\noindent\textbf{Itô's mnemonic rules.} The structure
of~\eqref{eq:formula_ito_general} is easily memorized by operating
with $dX_t = a_t\, dt + b_t\, dW_t$ as with an ordinary differential
object, except for the multiplication rules
\begin{equation}\label{eq:reglas_ito}
  (dt)^2 = 0, \qquad dt\, dW_t = 0, \qquad (dW_t)^2 = dt,
\end{equation}
which encode, respectively, that $dt$ is of higher order than
$\sqrt{dt}$ and the quadratic variation of Brownian motion of
Theorem~\ref{teo:variacion_cuadratica}. Applying a formal Taylor
expansion to $f(t, X_t)$ and using these rules to discard negligible
terms, one recovers exactly \eqref{eq:formula_ito_general}.

\medskip

\noindent\textbf{Three examples.} We illustrate the use of the
formula with three examples that will reappear later on.

\medskip

\noindent\textit{Example 1: the computation of $\int_0^t W_s\, dW_s$.}
We apply Theorem~\ref{teo:formula_ito_brown} to the function
$f(x) = \tfrac{1}{2} x^2$, whose derivatives are $f'(x) = x$ and
$f''(x) = 1$. Strictly speaking this function does not satisfy the
boundedness assumptions, but the result still holds by localization
arguments. The formula gives
\begin{equation*}
  \tfrac{1}{2} W_t^2 = 0 + \int_0^t W_s\, dW_s
                       + \tfrac{1}{2} \int_0^t 1\, ds,
\end{equation*}
that is,
\begin{equation}\label{eq:integral_W_dW}
  \int_0^t W_s\, dW_s = \tfrac{1}{2}\bigl(W_t^2 - t\bigr).
\end{equation}
We thus recover, by the systematic route of Itô calculus, the result
that in the motivation (Section~\ref{sec:motivacion_ito}) we had
obtained by direct computation of the sums.
Formula~\eqref{eq:integral_W_dW} also shows, via
Theorem~\ref{teo:martingala_integral_indefinida}, that
$\{W_t^2 - t\}$ is a martingale (a result that states the role of
this quantity already anticipated in
Proposition~\ref{prop:martingalas_brownianas}).

\medskip

\noindent\textit{Example 2: geometric Brownian motion.} Consider the
process
\begin{equation*}
  S_t \eqdef S_0\, \exp\bigl(\nu t + \sigma\, W_t\bigr),
\end{equation*}
where $S_0 > 0$, $\nu \in \R$ and $\sigma > 0$. We apply
Theorem~\ref{teo:formula_ito_general} to the function
$f(t, x) = S_0\, \exp(\nu t + \sigma x)$ and the process $X_t = W_t$
(which is an Itô process with $a \equiv 0$, $b \equiv 1$). The
derivatives of $f$ are
\begin{equation*}
  \partial_t f = \nu\, f, \qquad
  \partial_x f = \sigma\, f, \qquad
  \partial_{xx} f = \sigma^2\, f.
\end{equation*}
Substituting in the formula,
\begin{equation*}
  dS_t = \bigl[\nu\, S_t + 0 + \tfrac{1}{2} \sigma^2\, S_t\bigr]\, dt
        + \sigma\, S_t\, dW_t,
\end{equation*}
that is,
\begin{equation}\label{eq:gbm_sde}
  \boxed{\; dS_t = \mu\, S_t\, dt + \sigma\, S_t\, dW_t, \qquad
  \mu \eqdef \nu + \tfrac{1}{2}\sigma^2. \;}
\end{equation}
Identity~\eqref{eq:gbm_sde} is the \emph{stochastic differential
equation of geometric Brownian motion}, the canonical model for the
price of the underlying in Black-Scholes (Chapter~\ref{cap:bs}).

\medskip

\noindent\textit{Example 3: the generalized exponential martingale.}
Let $\{\theta_t\}_{t \in [0, T]}$ be an adapted process with
$\theta \in \mathcal{H}^2$. Consider the process $\{L_t\}$ defined by
\begin{equation}\label{eq:martingala_exp_general}
  L_t \eqdef \exp\!\biggl(- \int_0^t \theta_s\, dW_s
                           - \tfrac{1}{2} \int_0^t \theta_s^2\, ds\biggr).
\end{equation}
We apply Itô's formula to identify its SDE. Define
$Y_t \eqdef -\int_0^t \theta_s\, dW_s - \tfrac{1}{2} \int_0^t \theta_s^2\, ds$,
which is an Itô process with
$dY_t = -\tfrac{1}{2} \theta_t^2\, dt - \theta_t\, dW_t$. The process
$L_t = \exp(Y_t)$ is a function of the Itô process $Y$ via
$f(y) = e^y$, with $f'(y) = f''(y) = e^y = L_t$. By
Theorem~\ref{teo:formula_ito_general} applied to $f(Y_t)$, where the
drift coefficient of $Y$ is $-\tfrac{1}{2}\theta_t^2$ and the
diffusion coefficient is $-\theta_t$,
\begin{equation*}
  dL_t = L_t\, \biggl[-\tfrac{1}{2}\theta_t^2 +
                       \tfrac{1}{2}\, (-\theta_t)^2\biggr]\, dt
        + L_t\, (-\theta_t)\, dW_t.
\end{equation*}
The two terms in the bracket cancel exactly, leaving
\begin{equation}\label{eq:sde_martingala_exp}
  dL_t = -\theta_t\, L_t\, dW_t.
\end{equation}
The cancellation that produces the zero drift term is not accidental:
the correction $-\tfrac{1}{2}\int \theta_s^2\, ds$ appearing in the
exponent of~\eqref{eq:martingala_exp_general} is calibrated precisely
so that this happens. The consequence is that $\{L_t\}$ is a
driftless process, expressible as the stochastic integral of
$-\theta_t\, L_t$ with respect to $W$. By
Theorem~\ref{teo:martingala_integral_indefinida}, if
$\theta\, L \in \mathcal{H}^2$, then $\{L_t\}$ is a martingale.

The process~\eqref{eq:martingala_exp_general} generalizes the
exponential martingale $\mathcal{E}_t^\lambda$ of
Chapter~\ref{cap:brown}, replacing the parameter $\lambda$ by an
adapted process $\theta_t$. It will reappear in
§\ref{sec:girsanov} as the Radon-Nikodym derivative of Girsanov's
change of measure.

\medskip
\section{Girsanov's theorem}
\label{sec:girsanov}

In the binomial model of Chapter~\ref{cap:discreto}, the change
$\PP \to \Q$ was constructed explicitly by modifying the
probabilities of the random walk. In the continuous setting, what is
modified is the measure on the space of paths, and the question is
under which measure $\Q$ a process with drift $\theta_t$ relative to
$W$ becomes a martingale. The answer is provided by \emph{Girsanov's
theorem}.

\subsection{Change of measure in continuous time}
\label{subsec:cambio_medida}

Before stating the theorem, it is worth introducing the language and
concepts on which it will rest: equivalence of measures and the
Radon-Nikodym derivative in the continuous setting.

\medskip

\begin{definicion}[Equivalent measures]
\label{def:medidas_equivalentes}
Two probability measures $\PP$ and $\Q$ on a measurable space
$(\Omega, \F)$ are said to be \emph{equivalent}, written
$\Q \sim \PP$, if for all $A \in \F$,
\begin{equation*}
  \PP(A) = 0 \iff \Q(A) = 0.
\end{equation*}
\end{definicion}

Equivalence of measures is the minimal degree of compatibility we
shall need: two equivalent measures assign zero probability to
exactly the same events, although they may weigh the remaining ones
differently. In particular, statements of the type ``almost surely''
are interchangeable between the two.

\medskip

\noindent\textbf{Radon-Nikodym derivative.}

\begin{teorema}[Radon-Nikodym, equivalent version]
\label{teo:radon_nikodym}
Let $\PP$ and $\Q$ be two equivalent probability measures on
$(\Omega, \F)$. There exists a unique (up to almost sure equality)
random variable $L : \Omega \to (0, \infty)$ which is
$\F$-measurable, $\PP$-integrable, and satisfies
\begin{equation}\label{eq:def_RN}
  \Q(A) = \int_A L\, d\PP, \qquad A \in \F.
\end{equation}
The variable $L$ is called the \emph{Radon-Nikodym derivative of $\Q$
with respect to $\PP$}, and is denoted $L = d\Q/d\PP$.
\end{teorema}

The proof can be found in \cite{williams}.
Identity~\eqref{eq:def_RN} can be reformulated, equivalently, in
terms of expectations: for any random variable $X$ such that $L X$ is
$\PP$-integrable,
\begin{equation}\label{eq:cambio_esperanza_RN}
  \E_\Q[X] = \E_\PP[L X],
\end{equation}
a formula that determines how expectations under $\Q$ are computed
from expectations under $\PP$.

\begin{observacion}[Change of measure as a general method]
\label{obs:cambio_medida_general}
Identity~\eqref{eq:cambio_esperanza_RN} reveals an idea that goes
beyond the pricing framework: it allows one to compute expectations
under a measure $\Q$ by operating under another measure $\PP$, simply
by reweighting through the density $L$. Replacing a measure by an
equivalent one in order to turn a difficult computation into a
tractable one is a device of general use in probability; in
simulation, for instance, it is the foundation of \emph{importance
sampling}, which reduces the variance of Monte Carlo estimates of
rare events. Girsanov's theorem, which we address next, is the
version of this method tailored to Brownian motion.
\end{observacion}

\medskip

\noindent\textbf{The filtered case: the density as a process.} In
the setting of interest to us, recall that $(\Omega, \F)$ is
moreover endowed with a filtration $(\F_t)_{t \in [0, T]}$, the
augmented Brownian filtration. If $\PP$ and $\Q$ are equivalent on
$(\Omega, \F_T)$, we can restrict them to each $\F_t$ with $t < T$,
and thus obtain a family of Radon-Nikodym derivatives
\begin{equation}\label{eq:proceso_densidad}
  L_t \eqdef \frac{d(\Q\!\restriction_{\F_t})}{d(\PP\!\restriction_{\F_t})},
  \qquad t \in [0, T].
\end{equation}
The Radon-Nikodym derivative thus becomes a \emph{process} rather
than a random variable. The fundamental structural property of this
process is the following.

\begin{proposicion}[Martingale property of the density process]
\label{prop:martingala_densidad}
Let $\PP$ and $\Q$ be equivalent measures on $(\Omega, \F_T)$, and
let $\{L_t\}_{t \in [0, T]}$ be the process defined
in~\eqref{eq:proceso_densidad}. Then $\{L_t\}$ is a $\PP$-martingale
with respect to $(\F_t)$ with $\E_\PP[L_t] = 1$ for all $t$.
\end{proposicion}

\begin{proof}
Adaptedness and integrability are immediate: $L_t$ is
$\F_t$-measurable by construction, and $\E_\PP[L_t] = \Q(\Omega) = 1$
by \eqref{eq:def_RN} applied with $A = \Omega$.

For the martingale property, let $0 \leq s < t \leq T$ and
$A \in \F_s$. Since $\F_s \subseteq \F_t$, also $A \in \F_t$.
Applying \eqref{eq:def_RN} twice,
\begin{equation*}
  \int_A L_t\, d\PP = \Q(A) = \int_A L_s\, d\PP.
\end{equation*}
By the uniqueness of conditional expectation
(Definition~\ref{def:esperanza_condicional}),
$\E_\PP[L_t \mid \F_s] = L_s$ almost surely.
\end{proof}

\medskip

\noindent\textbf{Construction of equivalent measures via exponential
martingales.} We shall follow the reverse path: given a positive
$\PP$-martingale $\{L_t\}$ with $\E_\PP[L_T] = 1$, we can
\emph{define} a new measure $\Q$ on $(\Omega, \F_T)$ by
\begin{equation}\label{eq:def_Q_continuo}
  \Q(A) \eqdef \int_A L_T\, d\PP, \qquad A \in \F_T.
\end{equation}
Integrability and the condition $\E_\PP[L_T] = 1$ guarantee that $\Q$
is a probability measure, and the strict positivity of $L_T$ ensures
$\Q \sim \PP$. The Radon-Nikodym derivative is then
$d\Q/d\PP = L_T$, and the density process recovers, by
Proposition~\ref{prop:martingala_densidad}, the original martingale
$\{L_t\}$.

The question, therefore, is: what kind of positive martingales with
value one at $t = 0$ can we construct? Example~3 of
Section~\ref{subsec:formula_ito_general} offers the natural answer:
the exponential processes
\begin{equation}\label{eq:L_t_recordada}
  L_t = \exp\!\biggl(-\int_0^t \theta_s\, dW_s
                      - \tfrac{1}{2} \int_0^t \theta_s^2\, ds\biggr),
\end{equation}
associated with an adapted process $\theta$, are canonical
candidates. By construction they are positive and satisfy $L_0 = 1$.
We have also verified that they satisfy the SDE
$dL_t = -\theta_t\, L_t\, dW_t$, that is, they have no drift. By
Theorem~\ref{teo:martingala_integral_indefinida}, if
$\theta\, L \in \mathcal{H}^2$, then $\{L_t\}$ is a $\PP$-martingale.
And in that case $\E_\PP[L_T] = \E_\PP[L_0] = 1$, as required
by~\eqref{eq:def_Q_continuo}.

\medskip

\noindent\textbf{Novikov's condition.} The condition
$\theta\, L \in \mathcal{H}^2$ is delicate to verify in practice,
since it involves the very process $L$ we are defining. There is a
more manageable standard criterion, known as \emph{Novikov's
condition}, which guarantees that the
process~\eqref{eq:L_t_recordada} is indeed a martingale.

\begin{teorema}[Novikov's condition]
\label{teo:novikov}
Let $\{\theta_t\}_{t \in [0, T]}$ be an adapted process such that
\begin{equation}\label{eq:novikov}
  \E_\PP\!\left[\exp\!\biggl(\tfrac{1}{2} \int_0^T \theta_s^2\, ds\biggr)\right]
  < \infty.
\end{equation}
Then the process $\{L_t\}$ defined in~\eqref{eq:L_t_recordada} is a
$\PP$-martingale.
\end{teorema}

The proof of this result lies beyond the scope of the present work
and can be found in \cite{karatzas-shreve}, Section 3.5. We shall use
it as a tool to guarantee, in practice, the martingale condition. In
the applications we address in Chapter~\ref{cap:bs}, the process
$\theta$ will turn out to be deterministic (in fact constant), and
Novikov's condition will be trivially satisfied.

\subsection{Girsanov's theorem}
\label{subsec:girsanov_teorema}

\noindent\textbf{Lévy's characterization.}

\begin{teorema}[Lévy's characterization of Brownian motion]
\label{teo:caracterizacion_levy}
Let $\{X_t\}_{t \in [0, T]}$ be a continuous stochastic process
adapted to a filtration $(\F_t)$ on a probability space
$(\Omega, \F, \PP)$, with $X_0 = 0$ almost surely. If:
\begin{enumerate}[label=\textnormal{(\roman*)}]
  \item $\{X_t\}$ is a martingale with respect to $(\F_t)$ under $\PP$;
  \item the process $\{X_t^2 - t\}$ is also a martingale with respect
        to $(\F_t)$ under $\PP$;
\end{enumerate}
then $\{X_t\}$ is a standard Brownian motion under $\PP$.
\end{teorema}

The proof of this theorem, as formulated here, can be found in
\cite{etheridge}, Section 4.3.

Condition (ii) is equivalent, for continuous martingales, to
$[X]_t = t$. The characterization therefore says: a continuous
martingale with $X_0 = 0$ and quadratic variation $t$ is necessarily
a Brownian motion.

\medskip

\noindent\textbf{Girsanov's theorem.} We turn to the statement of the
central result.

\begin{teorema}[Girsanov]
\label{teo:girsanov}
Let $\{W_t\}_{t \in [0, T]}$ be a standard Brownian motion on
$(\Omega, \F, (\F_t), \PP)$, and let $\{\theta_t\}_{t \in [0, T]}$ be
an adapted process satisfying Novikov's condition
\eqref{eq:novikov}. Define $\{L_t\}$ by~\eqref{eq:L_t_recordada} and
the measure $\Q$ on $(\Omega, \F_T)$ by
$\Q(A) = \int_A L_T\, d\PP$. Then:
\begin{enumerate}[label=\textnormal{(\roman*)}]
  \item $\Q$ is a probability measure equivalent to $\PP$, with
        $d\Q/d\PP = L_T$;
  \item the process defined by
        \begin{equation}\label{eq:def_W_tilde}
          \widetilde{W}_t \eqdef W_t + \int_0^t \theta_s\, ds,
          \qquad t \in [0, T],
        \end{equation}
        is a standard Brownian motion on
        $(\Omega, \F, (\F_t), \Q)$.
\end{enumerate}
\end{teorema}

\begin{proof}
Claim (i) is essentially proved in the preceding subsection. By
Novikov's condition and Theorem~\ref{teo:novikov}, $\{L_t\}$ is a
strictly positive $\PP$-martingale with $L_0 = 1$, and therefore
$\E_\PP[L_T] = 1$. The map $A \mapsto \int_A L_T\, d\PP$ then defines
a probability measure on $(\Omega, \F_T)$, equivalent to $\PP$ by the
strict positivity of $L_T$, and with Radon-Nikodym derivative $L_T$
by construction.

For claim (ii) we apply Lévy's characterization
(Theorem~\ref{teo:caracterizacion_levy}) under the measure $\Q$. We
must verify that $\widetilde{W}$ satisfies the hypotheses of the
theorem: pathwise continuity, $\widetilde{W}_0 = 0$, the martingale
property under $\Q$, and the property that
$\{\widetilde{W}_t^2 - t\}$ is a $\Q$-martingale.

\emph{Step 1: continuity and initial value.} The continuity of
$\{\widetilde{W}_t\}$ is immediate: $W_t$ is continuous almost surely
(condition (iv) of Definition~\ref{def:browniano}) and the integral
$\int_0^t \theta_s\, ds$ is continuous in $t$ by construction. The
sum is therefore continuous. Likewise,
$\widetilde{W}_0 = W_0 + 0 = 0$.

\emph{Step 2: $\{\widetilde{W}_t\}$ is a $\Q$-martingale.} The
strategy is the following: by a classical lemma that we now detail,
a process $\{Y_t\}$ is a $\Q$-martingale if and only if
$\{Y_t\, L_t\}$ is a $\PP$-martingale. We shall apply this
characterization to $Y_t = \widetilde{W}_t$.

\noindent Let us first prove the lemma. If $\{Y_t L_t\}$ is a
$\PP$-martingale, then for $0 \leq s \leq t$ and $A \in \F_s$,
\begin{equation*}
  \int_A Y_t\, L_t\, d\PP = \int_A Y_s\, L_s\, d\PP.
\end{equation*}
The left-hand side, by~\eqref{eq:def_Q_continuo} applied with
restriction to $\F_t$ (recall that $L_t$ is the Radon-Nikodym
derivative of $\Q\!\restriction_{\F_t}$ with respect to
$\PP\!\restriction_{\F_t}$), equals $\int_A Y_t\, d\Q$. Similarly, the
right-hand side equals $\int_A Y_s\, d\Q$. By the uniqueness of
conditional expectation, $\E_\Q[Y_t \mid \F_s] = Y_s$, that is,
$\{Y_t\}$ is a $\Q$-martingale. The converse implication is obtained
by reversing the argument.

Having established the lemma, it suffices to prove that
$\{\widetilde{W}_t\, L_t\}$ is a $\PP$-martingale. Let us compute its
SDE under $\PP$ by applying the product rule, which is a particular
case of the general Itô formula
(Theorem~\ref{teo:formula_ito_general}) applied to $f(x, y) = xy$
with $X_t = \widetilde{W}_t$ and $Y_t = L_t$ as arguments. The
differentials of the processes are
\begin{equation*}
  d\widetilde{W}_t = \theta_t\, dt + dW_t, \qquad
  dL_t = -\theta_t\, L_t\, dW_t,
\end{equation*}
where the first identity follows from the
definition~\eqref{eq:def_W_tilde}, and the second was established
in~\eqref{eq:sde_martingala_exp}. The product rule yields
\begin{equation*}
  d(\widetilde{W}_t\, L_t)
  = L_t\, d\widetilde{W}_t + \widetilde{W}_t\, dL_t
   + d\widetilde{W}_t\, dL_t,
\end{equation*}
where the last term is the formal product of the two differentials,
evaluated by means of Itô's rules~\eqref{eq:reglas_ito}. Let us
compute it:
\begin{equation*}
  d\widetilde{W}_t\, dL_t
  = (\theta_t\, dt + dW_t) \cdot (-\theta_t\, L_t\, dW_t)
  = -\theta_t^2\, L_t\, dt\, dW_t - \theta_t\, L_t\, (dW_t)^2.
\end{equation*}
By Itô's rules, $dt\, dW_t = 0$ and $(dW_t)^2 = dt$, so that
$d\widetilde{W}_t\, dL_t = -\theta_t\, L_t\, dt$. Substituting,
\begin{align*}
  d(\widetilde{W}_t\, L_t)
  &= L_t\,(\theta_t\, dt + dW_t)
   + \widetilde{W}_t\, (-\theta_t\, L_t\, dW_t)
   - \theta_t\, L_t\, dt \\
  &= \underbrace{(L_t\, \theta_t - \theta_t\, L_t)}_{= 0}\, dt
   + L_t\, dW_t - \widetilde{W}_t\, \theta_t\, L_t\, dW_t \\
  &= L_t\, (1 - \theta_t\, \widetilde{W}_t)\, dW_t.
\end{align*}
The cancellation of the $dt$ term is the crucial computation of the
proof: it shows that $\{\widetilde{W}_t\, L_t\}$ is a driftless
process, expressible as the stochastic integral of the process
$L_t (1 - \theta_t\, \widetilde{W}_t)$ with respect to $W$. By
Theorem~\ref{teo:martingala_integral_indefinida}, assuming sufficient
integrability conditions on $\theta$, ultimately guaranteed by
Novikov's condition, $\{\widetilde{W}_t\, L_t\}$ is a
$\PP$-martingale. By the preceding lemma, $\{\widetilde{W}_t\}$ is a
$\Q$-martingale.

\emph{Step 3: $\{\widetilde{W}_t^2 - t\}$ is a $\Q$-martingale.} By
the characterization of Step 2, it suffices to prove that
$\{(\widetilde{W}_t^2 - t)\, L_t\}$ is a $\PP$-martingale. Applying
Itô's formula to $f(x) = x^2$ with $X_t = \widetilde{W}_t$ and using
$(d\widetilde{W}_t)^2 = dt$ by Itô's rules, we obtain
$d(\widetilde{W}_t^2 - t) = 2\, \widetilde{W}_t\, \theta_t\, dt + 2\, \widetilde{W}_t\, dW_t$.
Multiplying by $L_t$ and applying the product rule exactly as in
Step 2, the $dt$ terms cancel for the same underlying reason (the
correction $-\tfrac{1}{2} \int \theta_s^2\, ds$ in the exponent of
$L_t$ is calibrated to produce this cancellation), leaving
\begin{equation*}
  d\bigl((\widetilde{W}_t^2 - t)\, L_t\bigr)
  = (\text{terms in } dW_t),
\end{equation*}
with no $dt$ term. The process $(\widetilde{W}_t^2 - t)\, L_t$ is
therefore a $\PP$-martingale under the appropriate integrability
conditions, and by the equivalence lemma $\{\widetilde{W}_t^2 - t\}$
is a $\Q$-martingale.

\emph{Conclusion.} The process $\{\widetilde{W}_t\}$ is continuous,
satisfies $\widetilde{W}_0 = 0$, is a $\Q$-martingale, and
$\{\widetilde{W}_t^2 - t\}$ is also a $\Q$-martingale. By Lévy's
characterization (Theorem~\ref{teo:caracterizacion_levy}),
$\widetilde{W}$ is a standard Brownian motion under $\Q$.
\end{proof}

\medskip

\noindent\textbf{Reading of the theorem.} Girsanov's theorem is a
result of great interest for the purposes of our work: the change of
measure encoded by the exponential martingale $\{L_t\}$
\emph{absorbs the additional drift $\theta_t$ of the process
$W_t + \int_0^t \theta_s\, ds$ until it disappears}. Under the new
measure $\Q$, what was a process with drift $\theta_t$ becomes a pure
Brownian motion. The operation is reversible: applying the theorem
with $-\theta$ recovers $W$ from $\widetilde{W}$.

\medskip

\noindent\textbf{Reformulation in terms of Itô processes.} The
practical use of the theorem in Chapter~\ref{cap:bs} will go through
the following reformulation. Let $\{X_t\}$ be an Itô process with
dynamics
\begin{equation*}
  dX_t = a_t\, dt + b_t\, dW_t,
\end{equation*}
and suppose $b_t > 0$ almost surely. Define
$\theta_t \eqdef a_t / b_t$, and suppose it satisfies Novikov's
condition. Applying Theorem~\ref{teo:girsanov} with this $\theta$,
under the resulting measure $\Q$ the process $\widetilde{W}$ defined
in~\eqref{eq:def_W_tilde} is a Brownian motion. Substituting
$dW_t = d\widetilde{W}_t - \theta_t\, dt$ into the SDE of $X$,
\begin{equation*}
  dX_t = a_t\, dt + b_t\, (d\widetilde{W}_t - \theta_t\, dt)
       = (a_t - b_t\, \theta_t)\, dt + b_t\, d\widetilde{W}_t.
\end{equation*}
Substituting $\theta_t = a_t/b_t$, the coefficient of $dt$ vanishes,
and
\begin{equation}\label{eq:girsanov_operativo}
  \boxed{\; dX_t = b_t\, d\widetilde{W}_t. \;}
\end{equation}
Under the measure $\Q$, the process $X$ loses its drift and becomes a
pure stochastic integral with respect to the new Brownian motion
$\widetilde{W}$. In particular, by
Theorem~\ref{teo:martingala_integral_indefinida}, $\{X_t\}$ is a
$\Q$-martingale.
\section{The Brownian martingale representation theorem}
\label{sec:representacion_martingalas}

This section presents the generalization to the continuous setting of
the binomial representation theorem
(Subsection~\ref{subsec:representacion_binomial}). Unlike the
discrete case, where the integrand was constructed explicitly by
solving a linear system node by node, the continuous version relies
on functional-analytic arguments and \emph{does not provide an
explicit formula}: it guarantees the existence of the integrand, not
its construction.

\medskip

Throughout the section we work on the filtered probability space
$(\Omega, \F, (\F_t)_{t \in [0, T]}, \PP)$, where $(\F_t)$ is the
filtration of the Brownian motion $\{W_t\}$ fixed in
Definition~\ref{def:filtracion_aumentada}. It is essential to stress
that the theorem requires working with exactly this filtration
(generated by the Brownian motion together with the null sets): if
the underlying filtration contained strictly additional information,
extraneous to the Brownian motion, the result would cease to hold.
This restriction will not be a problem in the applications of
Chapter~\ref{cap:bs}, where the only source of randomness will be
precisely the Brownian motion modeling the underlying.

\begin{teorema}[Representation of $L^2$ random variables]
\label{teo:representacion_l2}
Let $X$ be an $\F_T$-measurable random variable with
$\E[X^2] < \infty$. There exists a unique process
$\phi \in \mathcal{H}^2[0, T]$ such that
\begin{equation}\label{eq:representacion_l2}
  X = \E[X] + \int_0^T \phi_s\, dW_s
  \qquad \text{almost surely}.
\end{equation}
\end{teorema}

The proof lies beyond the scope of the present work, but we can
describe the strategy in detail; it is the same one we already
employed when constructing the stochastic integral in
Section~\ref{subsec:extension_H2}: identify a sufficiently large
subspace of $L^2(\Omega, \F_T, \PP)$ where the representation is
direct, and extend it by density and isometry to the whole space.

Specifically, one considers the subspace
\[
  K \eqdef \biggl\{ M_0 + \int_0^T \phi_s\, dW_s :
                    M_0 \in \R,\ \phi \in \mathcal{H}^2[0, T] \biggr\}
  \subseteq L^2(\Omega, \F_T, \PP),
\]
consisting precisely of the random variables that admit the desired
representation. From here the proof rests on two facts.

The first is that $K$ is a \emph{closed} subspace of
$L^2(\Omega, \F_T, \PP)$. The verification is direct: if
$(Y_n) \subset K$ converges in $L^2$ to some $Y$, the constants
$\E[Y_n]$ converge to $\E[Y]$ by continuity of expectation, and the
associated integrands $\phi^{(n)}$ form a Cauchy sequence in
$\mathcal{H}^2[0, T]$ by the \emph{Itô isometry}
(Theorem~\ref{teo:propiedades_integral_H2}); their limit $\phi$
provides the representation of $Y$, so $Y \in K$.

The second is that the \emph{stochastic exponentials} of the form
\begin{equation*}
  \mathcal{E}(f)
  \eqdef \exp\!\biggl(\int_0^T f(s)\, dW_s
                       - \tfrac{1}{2} \int_0^T f(s)^2\, ds\biggr),
\end{equation*}
with $f : [0, T] \to \R$ deterministic and square-integrable, belong
to $K$ and, moreover, span a dense subspace of
$L^2(\Omega, \F_T, \PP)$. Membership in $K$ is a direct computation
with Itô's formula applied to the SDE
$d\mathcal{E}(f)_t = \mathcal{E}(f)_t\, f(t)\, dW_t$, which in
integral form gives
$\mathcal{E}(f) = 1 + \int_0^T \mathcal{E}(f)_s\, f(s)\, dW_s$.
Density, on the other hand, is the truly delicate part of the
argument: it is proved by showing that no $Z \in L^2(\Omega, \F_T, \PP)$
can be orthogonal to all the exponentials, which reduces, via an
analytic change of parameter, to an injectivity argument for the
Fourier transform of the finite-dimensional Gaussian distributions of
Brownian motion.

Combining both facts, $K$ is a closed subspace containing a dense
subset, and therefore coincides with all of
$L^2(\Omega, \F_T, \PP)$. The uniqueness of the integrand $\phi$
follows, once again, from the Itô isometry. For a detailed
exposition we refer to \cite{etheridge}, Section~4.6.

\medskip

\noindent\textbf{Version for martingales.}
Theorem~\ref{teo:representacion_l2} admits the following general
formulation, which is the one we shall work with in
Chapter~\ref{cap:bs}.

\begin{teorema}[Brownian martingale representation]
\label{teo:representacion_martingalas}
Let $\{M_t\}_{t \in [0, T]}$ be a martingale with respect to the
Brownian filtration $(\F_t)$, with $M_t \in L^2$ for all
$t \in [0, T]$. There exists a unique process
$\phi \in \mathcal{H}^2[0, T]$ such that
\begin{equation}\label{eq:representacion_martingalas}
  M_t = M_0 + \int_0^t \phi_s\, dW_s,
  \qquad t \in [0, T],
  \quad \text{almost surely}.
\end{equation}
\end{teorema}

\begin{proof}
Applying Theorem~\ref{teo:representacion_l2} to the variable $M_T$,
there exists a unique $\phi \in \mathcal{H}^2$ with
\begin{equation*}
  M_T = M_0 + \int_0^T \phi_s\, dW_s
\end{equation*}
almost surely, where we have used $\E[M_T] = M_0$ by the martingale
property. To extend the identity to $t < T$, we take conditional
expectations with respect to $\F_t$ on both sides. The left-hand side
gives $\E[M_T \mid \F_t] = M_t$ by the martingale property of $M$.
The right-hand side gives
$M_0 + \E[\int_0^T \phi_s\, dW_s \mid \F_t] = M_0 + \int_0^t \phi_s\, dW_s$,
since the indefinite integral up to $T$ is a martingale
(Theorem~\ref{teo:martingala_integral_indefinida}) and therefore its
conditional expectation given $\F_t$ is the integral up to $t$.
Identity~\eqref{eq:representacion_martingalas} is proved. Uniqueness
follows from the uniqueness in Theorem~\ref{teo:representacion_l2}.
\end{proof}

\medskip

\noindent\textbf{Financial reading of the integrand.} The financial
interpretation of the integrand $\phi$ in
Theorem~\ref{teo:representacion_martingalas} is entirely analogous to
that of the discrete case. If $\{M_t\}$ is the martingale associated
with the discounted value of a derivative, constructed in
Chapter~\ref{cap:bs} as $M_t = \E_\Q[e^{-rT}\, C_T \mid \F_t]$
(following the same pattern as in
\eqref{eq:def_M_replicante_discreto}), then the process $\phi_t$
appearing in its representation is the \emph{quantity of the
underlying} that the replicating portfolio must hold at time $t$. It
is the generalization to the continuous setting of the explicit
formula~\eqref{eq:phi_explicito} of the binomial model, with the
caveat already noted: in the continuous case $\phi$ exists but is not
constructed explicitly.

\section{Synthesis: from the binomial model to stochastic calculus}
\label{sec:sintesis_cap4}

We close the chapter by establishing the correspondence between the
tools of the binomial model and those of continuous stochastic
calculus.

\medskip

\noindent\textbf{The correspondence between the two settings.} Each
of the four pieces of continuous stochastic calculus has its
counterpart in the discrete case. The \emph{Itô stochastic
integral}, $\int_0^t \phi_s\, dW_s$, is the continuous version of the
discrete stochastic integral $(\phi \cdot \widetilde S)_n$
of~\eqref{eq:integral_discreta}: both express the accumulated gain of
a predictable strategy $\phi$, and the predictability condition,
formulated in discrete time as $\F_{k-1}$-measurability and in
continuous time as measurability with respect to $\F_{t_i}$ (on
simple processes), is what makes the integral preserve the martingale
property (Theorems~\ref{teo:integral_preserva_martingala}
and~\ref{teo:martingala_integral_indefinida}). \emph{Itô's formula}
is the continuous version of the classical chain rule, with the
novelty of the extra term $\tfrac{1}{2}\, b_t^2\, \partial_{xx} f\, dt$
in~\eqref{eq:formula_ito_general}, which has no counterpart in the
discrete case and reflects the nonzero quadratic variation of
Brownian motion.

\emph{Girsanov's theorem} is the continuous version of the explicit
construction of the measure $\Q$ in the binomial
model~\eqref{eq:def_Q_multi}. In both cases, the change of measure is
governed by a density $L = d\Q/d\PP$. The difference is that in
discrete time $L$ is constructed combinatorially, assigning to each
path a probability $q^j(1-q)^{N-j}$, whereas in continuous time $L_T$
is the exponential martingale
$\exp(-\int \theta\, dW - \tfrac{1}{2}\int \theta^2\, ds)$. The
\emph{Brownian martingale representation theorem}
(Theorem~\ref{teo:representacion_martingalas}) is the continuous
version of the binomial representation theorem
(Theorem~\ref{teo:representacion_binomial}). In the discrete case,
the proof was constructive and we have an explicit formula for the
integrand; in the continuous case, by contrast, the proof appeals to
abstract density and completeness arguments and provides no explicit
expression: knowing how to construct it will be, in each specific
problem, an additional question to be settled by means of Itô's
formula.

\medskip

\noindent\textbf{The three-step recipe, continuous version.} The four
pieces above are the ingredients with which we shall obtain the
solution of the pricing problem in the continuous setting. The
strategy, faithful to the scheme of the discrete binomial model
(Subsection~\ref{subsec:tres_pasos_discreto}), will follow the same
three steps:

\begin{enumerate}[label=\textbf{Step \arabic*.}, leftmargin=*]
  \item \emph{Change of measure.} Identify the measure $\Q$ under
        which the discounted price of the underlying is a martingale.
        This identification, carried out in the discrete case by
        explicit construction, will be carried out in the continuous
        case by applying \emph{Girsanov's theorem}.
  \item \emph{Construction of the discounted value process.} Define
        the process $\{M_t\}$ associated with the discounted
        \emph{payoff} of the derivative via
        $M_t = \E_\Q[e^{-rT}\, C_T \mid \F_t]$. By the tower property,
        $\{M_t\}$ is a $\Q$-martingale. This step is identical to
        that of the discrete case.
  \item \emph{Application of the representation theorem.} Apply the
        \emph{Brownian martingale representation theorem} to obtain
        a process $\phi$ such that
        $M_t = M_0 + \int_0^t \phi_s\, d\widetilde{W}_s$, where
        $\widetilde{W}$ is the Brownian motion under $\Q$ produced by
        Girsanov. The integrand $\phi$ encodes the quantity of the
        underlying in the replicating portfolio.
\end{enumerate}

The structure is the same as in the discrete case: only the technical
sophistication of the tools changes. Chapter~\ref{cap:bs} applies
this recipe to geometric Brownian motion as the underlying and
obtains, as its central result, the Black-Scholes formula for the
European \emph{call} option.

% =====================================================================
%   CHAPTER 5 — THE BLACK-SCHOLES MODEL
% =====================================================================
\chapter{The Black-Scholes model}\label{cap:bs}

We reach the ultimate goal of this work: to apply the apparatus of
stochastic calculus to the \emph{Black-Scholes model} and obtain, as
its central result, the classical pricing formula for the European
\emph{call} option. The structure of the chapter follows the same
pattern as Chapter~\ref{cap:discreto}: model assumptions, execution
of the three-step recipe, specialization to the \emph{call}, and
recovery of the classical PDE.

\section{The Black-Scholes model}
\label{sec:modelo_bs}

\subsection{Market assumptions and dynamics of the underlying}
\label{subsec:hipotesis_modelo_bs}

\noindent\textbf{The model assumptions.} The Black-Scholes model
describes an idealized market built on the following structural
assumptions:

\begin{enumerate}[label=\textnormal{(H\arabic*)}, leftmargin=*]
  \item The market consists of two assets: a risky \emph{underlying}
        whose price is a stochastic process $\{S_t\}_{t \in [0, T]}$,
        and a \emph{risk-free bond} with constant interest rate
        $r > 0$.
  \item The underlying follows a geometric Brownian motion (GBM): its
        dynamics are governed by the stochastic differential equation
        \begin{equation}\label{eq:bs_dinamica}
          dS_t = \mu\, S_t\, dt + \sigma\, S_t\, dW_t,
          \qquad S_0 > 0,
        \end{equation}
        where $\mu \in \R$ is the expected return and $\sigma > 0$ is
        the \emph{volatility}, both constant, and $\{W_t\}$ is a
        standard Brownian motion on a filtered probability space
        $(\Omega, \F, (\F_t), \PP)$ with $(\F_t)$ the augmented
        Brownian filtration.
  \item The market is \emph{frictionless}: there are no transaction
        costs, assets are infinitely divisible, short selling is
        allowed, and lending and borrowing take place at the same
        rate $r$.
  \item The market is \emph{arbitrage-free}: there is no strategy
        which, starting from zero capital, produces with strictly
        positive probability a strictly positive terminal value and
        never a negative one.
\end{enumerate}

Assumptions (H1) and (H3)–(H4) are the literal translation to the
continuous setting of the assumptions of the binomial model of
Chapter~\ref{cap:discreto}: the market structure, the idealization
and the no-arbitrage condition remain intact. The novelty is (H2):
the specific choice of the dynamics of the underlying, which in the
discrete case was the multiplicative binomial random walk and in the
continuous case is the GBM.

\medskip

\noindent\textbf{Explicit form of the underlying.} The SDE
\eqref{eq:bs_dinamica} admits a closed-form solution whose expression
we already obtained as Example~2 of
Section~\ref{subsec:formula_ito_general}: the GBM with SDE
\eqref{eq:bs_dinamica} is exactly
\begin{equation}\label{eq:bs_S_explicito}
  S_t = S_0\, \exp\!\left(\Bigl(\mu - \tfrac{1}{2}\sigma^2\Bigr) t
                          + \sigma\, W_t\right),
  \qquad t \in [0, T].
\end{equation}
Identity~\eqref{eq:bs_S_explicito} shows that $S_t$ is strictly
positive almost surely (an economically required property for the
price of an asset), and that its logarithm $\log S_t$ has Gaussian
distribution
$\mathcal{N}\bigl(\log S_0 + (\mu - \tfrac{1}{2}\sigma^2) t,\, \sigma^2 t\bigr)$.
This log-normality property will be the key ingredient in the
explicit computation of the Black-Scholes formula in
Section~\ref{sec:formula_bs}.

\medskip

\noindent\textbf{The risk-free bond.} With constant interest rate
$r$, the price of the risk-free bond evolves deterministically
according to
\begin{equation}\label{eq:bs_bono}
  B_t = e^{rt},
  \qquad t \in [0, T],
\end{equation}
which is the solution of the ordinary differential equation
$dB_t = r\, B_t\, dt$ with $B_0 = 1$, with no stochastic term (the
bond carries no risk).

\subsection{Portfolios in continuous time}
\label{subsec:portfolio_continuo}

\begin{definicion}[Portfolio in continuous time]
\label{def:portfolio_continuo}
A \emph{portfolio} (or \emph{strategy}) on the horizon $[0, T]$ is a
pair of stochastic processes $(\phi_t, \psi_t)_{t \in [0, T]}$ where:
\begin{enumerate}[label=\textnormal{(\roman*)}]
  \item $\phi_t$ represents the quantity of the underlying in the
        portfolio at time $t$;
  \item $\psi_t$ represents the quantity of the risk-free bond in the
        portfolio at time $t$;
  \item both processes are adapted to the Brownian filtration
        $(\F_t)$ and satisfy the technical integrability conditions
        \begin{equation*}
          \phi \in \mathcal{H}^2[0, T], \qquad
          \int_0^T |\psi_s|\, ds < \infty \text{ almost surely}.
        \end{equation*}
\end{enumerate}
The \emph{value} of the portfolio at time $t$ is denoted $V_t$ and is
given by
\begin{equation}\label{eq:valor_portfolio}
  V_t \eqdef \phi_t\, S_t + \psi_t\, B_t.
\end{equation}
\end{definicion}

The integrability conditions on $(\phi, \psi)$ are those guaranteeing
that the integrals that will appear next, the stochastic integral
$\int \phi\, dS$ and the ordinary integral $\int \psi\, dB$, are well
defined in the sense of Chapter~\ref{cap:itocalculus}. The
adaptedness of $\phi$ and $\psi$ to $(\F_t)$ is the translation to
the continuous setting of discrete predictability
(Section~\ref{subsec:predecibles_int_disc}): the portfolio is decided
with the information available before the market moves.

\subsection{The \emph{self-financing} condition}
\label{subsec:self_financing_continuo}

The most delicate structural piece of the continuous model is the
translation to the continuous setting of the \emph{self-financing}
condition introduced in the discrete case in
Section~\ref{subsec:tres_pasos_discreto}. Let us recall its content:
a binomial strategy was \emph{self-financing} if at each node of the
tree
\begin{equation}\label{eq:self_financing_discreto_recordatorio}
  \phi_{k+1}\, S_k + \psi_{k+1}\, B_k = \phi_k\, S_k + \psi_k\, B_k,
\end{equation}
that is, if the amount of capital available at the beginning of
period $k+1$, after reallocating the portfolio, coincides with that
available at the end of period $k$. Rearranging terms,
\eqref{eq:self_financing_discreto_recordatorio} is equivalent to
\begin{equation*}
  S_k\, (\phi_{k+1} - \phi_k) + B_k\, (\psi_{k+1} - \psi_k) = 0,
\end{equation*}
and consequently, summing over the periods up to time $n$,
\begin{equation*}
  V_n - V_0
  = \sum_{k=0}^{n-1} \bigl[\phi_{k+1}\, (S_{k+1} - S_k)
                            + \psi_{k+1}\, (B_{k+1} - B_k)\bigr].
\end{equation*}
The value of the portfolio changes exclusively due to market
movements, not to injections or withdrawals of capital.

\medskip

\begin{definicion}[\emph{Self-financing} strategy]
\label{def:self_financing_continuo}
A strategy $(\phi_t, \psi_t)_{t \in [0, T]}$ is said to be
\emph{self-financing} if its value $V_t = \phi_t\, S_t + \psi_t\, B_t$
satisfies
\begin{equation}\label{eq:self_financing_continuo}
  dV_t = \phi_t\, dS_t + \psi_t\, dB_t,
  \qquad t \in [0, T],
\end{equation}
or, in integral form,
\begin{equation}\label{eq:self_financing_integral}
  V_t = V_0 + \int_0^t \phi_s\, dS_s + \int_0^t \psi_s\, dB_s,
  \qquad t \in [0, T].
\end{equation}
\end{definicion}

Identity~\eqref{eq:self_financing_continuo} is the exact translation
to the continuous setting of the discrete condition: it states that
the change in the value of the portfolio at time $t$ comes
exclusively from changes in the prices of the underlying and of the
bond, weighted respectively by the quantities $\phi_t$ and $\psi_t$
held at each instant.

\medskip

\noindent\textbf{Equivalence with the primitive formulation.}
Definition~\ref{def:self_financing_continuo} is the operational
formulation of the self-financing condition. It is worth verifying
its equivalence with the original formulation, which requires zero
cost in the infinitesimal reallocation of the portfolio.

\begin{proposicion}[Equivalence of the two formulations]
\label{prop:equivalencia_self_financing}
Let $(\phi_t, \psi_t)$ be a strategy with
$V_t = \phi_t\, S_t + \psi_t\, B_t$. The following conditions are
equivalent:
\begin{enumerate}[label=\textnormal{(\roman*)}]
  \item $dV_t = \phi_t\, dS_t + \psi_t\, dB_t$;
  \item $S_t\, d\phi_t + B_t\, d\psi_t + d\langle \phi, S\rangle_t = 0$,
\end{enumerate}
where $\langle \phi, S\rangle_t$ denotes the quadratic covariation of
the processes $\phi$ and $S$.
\end{proposicion}

\begin{proof}
We apply the product rule (a particular case of Itô's formula,
Theorem~\ref{teo:formula_ito_general}) to the product $\phi_t\, S_t$:
\begin{equation*}
  d(\phi_t\, S_t) = \phi_t\, dS_t + S_t\, d\phi_t + d\langle \phi, S\rangle_t.
\end{equation*}
Similarly, for $\psi_t\, B_t$, where $B_t$ has bounded variation, so
that $\langle \psi, B\rangle_t = 0$,
\begin{equation*}
  d(\psi_t\, B_t) = \psi_t\, dB_t + B_t\, d\psi_t.
\end{equation*}
Adding both identities, we obtain the differential of the value:
\begin{equation*}
  dV_t = d(\phi_t S_t + \psi_t B_t)
       = \phi_t\, dS_t + \psi_t\, dB_t
       + \bigl[S_t\, d\phi_t + B_t\, d\psi_t + d\langle \phi, S\rangle_t\bigr].
\end{equation*}
Condition (i) is therefore equivalent to the vanishing of the term in
brackets, which is exactly condition (ii).
\end{proof}
\section{Application of the three-step recipe}
\label{sec:tres_pasos_continuo}

In this section we execute the three-step recipe set out in
Subsection~\ref{subsec:tres_pasos_discreto} for the discrete case,
now in the continuous setting.

\subsection{Statement of the problem}
\label{subsec:planteamiento_continuo}

\noindent\textbf{European derivatives.}

\begin{definicion}[European derivative]
\label{def:derivado_europeo}
A \emph{European derivative} with maturity $T$ on the underlying
$\{S_t\}$ is a contract that pays the holder, at time $t = T$, an
amount $C_T = h(S_T)$, where $h : (0, \infty) \to \R$ is a measurable
function, the \emph{payoff} of the derivative, satisfying the
integrability condition
\begin{equation}\label{eq:integrabilidad_payoff}
  \E_\PP\bigl[h(S_T)^2\bigr] < \infty.
\end{equation}
\end{definicion}

Condition~\eqref{eq:integrabilidad_payoff} is the minimum required
for the processes built from the \emph{payoff} to lie in the $L^2$
spaces demanded by the Brownian martingale representation theorem.
For standard \emph{payoffs}, including the European \emph{call}
$h(S_T) = (S_T - K)^+$ that we address in
Section~\ref{sec:formula_bs}, the condition holds automatically by
the log-normality of $S_T$.

\medskip

\noindent\textbf{The pricing problem.} The central problem is to
determine, at each time $t \in [0, T]$, a fair price $V(t, S_t)$ for
the derivative, that is, a price compatible with the no-arbitrage
assumption (H4). The procedure to follow, as in the discrete scheme,
goes through the notion of \emph{replicating strategy}:

\begin{definicion}[Replicating strategy]
\label{def:estrategia_replicante}
A \emph{replicating strategy} for the European derivative $h(S_T)$ is
a \emph{self-financing} strategy $(\phi_t, \psi_t)_{t \in [0, T]}$
whose terminal value coincides with the \emph{payoff}, that is,
\begin{equation*}
  V_T = \phi_T\, S_T + \psi_T\, B_T = h(S_T)
  \qquad \text{almost surely}.
\end{equation*}
\end{definicion}

By the no-arbitrage argument of the discrete case, if the replicating
strategy exists, its value $V_t$ at each time is necessarily the
price of the derivative at $t$. The problem therefore reduces to
constructing a replicating strategy.

\subsection{Step 1: change of measure via Girsanov}
\label{subsec:paso1_girsanov}

We replace the market measure $\PP$ by an equivalent measure $\Q$
under which the discounted price of the underlying is a martingale.
The identification is carried out via Girsanov's
theorem.\footnote{Here reappears, in its continuous form, the idea
anticipated in Remark~\ref{obs:cambio_medida_general}: passing to an
equivalent measure (the risk-neutral one) to simplify the problem.
What in the binomial model was an explicit reweighting of the
probabilities on the paths of the tree is now the change of measure
established by Girsanov's theorem.}

\medskip

\noindent\textbf{The dynamics of the discounted price.} To set up the
change of measure, it is convenient first to compute the SDE of the
discounted price process
\begin{equation}\label{eq:precio_descontado_def}
  \widetilde S_t \eqdef e^{-rt}\, S_t,
  \qquad t \in [0, T].
\end{equation}
We apply Itô's formula (Theorem~\ref{teo:formula_ito_general}) to
the function $f(t, x) = e^{-rt}\, x$ and the process $X_t = S_t$. The
relevant derivatives are
\begin{equation*}
  \partial_t f = -r\, e^{-rt}\, x, \qquad
  \partial_x f = e^{-rt}, \qquad
  \partial_{xx} f = 0.
\end{equation*}
With $a_t = \mu\, S_t$, $b_t = \sigma\, S_t$ the coefficients of the
SDE \eqref{eq:bs_dinamica}, Itô's formula gives
\begin{align*}
  d\widetilde S_t
  &= \bigl[-r\, e^{-rt}\, S_t + e^{-rt}\, \mu\, S_t + 0\bigr]\, dt
   + e^{-rt}\, \sigma\, S_t\, dW_t \\
  &= (\mu - r)\, e^{-rt}\, S_t\, dt + \sigma\, e^{-rt}\, S_t\, dW_t.
\end{align*}
Since $e^{-rt} S_t = \widetilde S_t$, we obtain
\begin{equation}\label{eq:sde_S_descontado}
  \boxed{\; d\widetilde S_t
        = (\mu - r)\, \widetilde S_t\, dt
        + \sigma\, \widetilde S_t\, dW_t. \;}
\end{equation}
The SDE~\eqref{eq:sde_S_descontado} shows that the discounted price
is a new Itô process with drift coefficient
$(\mu - r)\, \widetilde S_t$ (proportional to the excess return of
the underlying over the risk-free rate) and, therefore,
$\{\widetilde S_t\}$ is not a martingale under $\PP$. The goal of the
change of measure is to remove that drift.

\medskip

\noindent\textbf{The choice of $\theta$.} The SDE
\eqref{eq:sde_S_descontado} fits exactly into the scheme of Girsanov's
theorem reformulated in terms of Itô processes,
eq.~\eqref{eq:girsanov_operativo}: we have a process with dynamics
$dX_t = a_t\, dt + b_t\, dW_t$ and we want to remove the drift $a_t$
by means of the choice
\begin{equation*}
  \theta_t \eqdef \frac{a_t}{b_t}
              = \frac{(\mu - r)\, \widetilde S_t}
                     {\sigma\, \widetilde S_t}
              = \frac{\mu - r}{\sigma}.
\end{equation*}
The quantity
\begin{equation}\label{eq:precio_riesgo}
  \theta \eqdef \frac{\mu - r}{\sigma},
\end{equation}
constant in the Black-Scholes model, goes by the classical name of
\emph{market price of risk}: it is exactly the quantity that
Girsanov's theorem requires in order to cancel the drift of the
discounted price.

\medskip

\noindent\textbf{Verification of Novikov's condition.} Since $\theta$
is constant, Novikov's condition~\eqref{eq:novikov} holds trivially:
\begin{equation*}
  \E_\PP\!\left[\exp\!\biggl(\tfrac{1}{2}\int_0^T \theta^2\, ds\biggr)\right]
  = \exp\!\biggl(\tfrac{1}{2}\theta^2\, T\biggr) < \infty.
\end{equation*}
We are therefore in a position to apply Theorem~\ref{teo:girsanov}
with $\theta_t \equiv \theta$.

\medskip

\noindent\textbf{Application of Girsanov's theorem.} The density
process associated with $\theta$ is, according
to~\eqref{eq:L_t_recordada},
\begin{equation}\label{eq:L_T_BS}
  L_T = \exp\!\biggl(-\theta\, W_T - \tfrac{1}{2}\,\theta^2\, T\biggr).
\end{equation}
We define the measure $\Q$ on $(\Omega, \F_T)$ by
\begin{equation*}
  \Q(A) \eqdef \int_A L_T\, d\PP, \qquad A \in \F_T.
\end{equation*}
By Theorem~\ref{teo:girsanov}, $\Q$ is a probability measure
equivalent to $\PP$, with $d\Q/d\PP = L_T$, and the process
\begin{equation}\label{eq:W_tilde_BS}
  \widetilde W_t \eqdef W_t + \theta\, t,
  \qquad t \in [0, T],
\end{equation}
is a standard Brownian motion under $\Q$.

\medskip

\noindent\textbf{Characterization of $\Q$ as a martingale measure.}
Substituting $dW_t = d\widetilde W_t - \theta\, dt$ into the SDE
\eqref{eq:sde_S_descontado} of the discounted price:
\begin{equation*}
  d\widetilde S_t = (\mu - r)\, \widetilde S_t\, dt
                  + \sigma\, \widetilde S_t\, (d\widetilde W_t - \theta\, dt).
\end{equation*}
Recalling $\theta = (\mu - r)/\sigma$, the $dt$ terms cancel:
\begin{equation}\label{eq:S_descontado_Q_martingala}
  \boxed{\; d\widetilde S_t = \sigma\, \widetilde S_t\, d\widetilde W_t. \;}
\end{equation}
Identity~\eqref{eq:S_descontado_Q_martingala} shows that, under $\Q$,
the discounted price of the underlying solves a driftless SDE: it is
a pure stochastic integral with respect to the Brownian motion
$\widetilde W$. Thus, by
Theorem~\ref{teo:martingala_integral_indefinida}, $\{\widetilde S_t\}$
is a $\Q$-martingale. The measure $\Q$ is accordingly called the
\emph{equivalent martingale measure} or \emph{risk-neutral measure}.

\medskip

\noindent\textbf{The dynamics of the underlying under $\Q$.} It is
also worth making explicit the dynamics of the underlying $S_t$
itself (undiscounted) under the new measure. Substituting
$dW_t = d\widetilde W_t - \theta\, dt$ into~\eqref{eq:bs_dinamica},
\begin{equation}\label{eq:S_dinamica_Q}
  \boxed{\; dS_t = r\, S_t\, dt + \sigma\, S_t\, d\widetilde W_t. \;}
\end{equation}
Under $\Q$, the underlying is still a GBM, but with \emph{trend $r$
instead of $\mu$}. This is the translation to the continuous setting
of the phenomenon we already saw in the binomial model: the change of
measure modifies the dynamics of the underlying so that its expected
return coincides with that of the risk-free bond, whence the name
\emph{risk-neutral}.

\noindent The explicit form of the underlying under $\Q$, analogous
to~\eqref{eq:bs_S_explicito}, is obtained by solving the
SDE~\eqref{eq:S_dinamica_Q} by the same technique as in Example~2 of
Section~\ref{subsec:formula_ito_general}:
\begin{equation}\label{eq:S_explicito_Q}
  S_t = S_0\, \exp\!\left(\Bigl(r - \tfrac{1}{2}\sigma^2\Bigr) t
                          + \sigma\, \widetilde W_t\right).
\end{equation}
Identity~\eqref{eq:S_explicito_Q} will be the key tool for the
explicit computations of Section~\ref{sec:formula_bs}. Under $\Q$,
$\log S_T$ has Gaussian distribution
$\mathcal{N}\bigl(\log S_0 + (r - \tfrac{1}{2}\sigma^2) T,\, \sigma^2 T\bigr)$.

\subsection{Step 2: the $\Q$-martingale associated with the discounted \emph{payoff}}
\label{subsec:paso2_martingala}

Having identified the risk-neutral measure $\Q$, the second step of
the recipe, in analogy with the discrete case
(Subsection~\ref{subsec:tres_pasos_discreto}), consists in
constructing a $\Q$-martingale process associated with the
\emph{payoff}.

\medskip

\noindent\textbf{Definition of the process.} Let $h(S_T)$ be the
\emph{payoff} of the derivative. We define the process
$\{M_t\}_{t \in [0, T]}$ by
\begin{equation}\label{eq:def_M_continuo}
  M_t \eqdef \E_\Q\!\bigl[e^{-rT}\, h(S_T) \,\big|\, \F_t\bigr],
  \qquad t \in [0, T].
\end{equation}
The financial interpretation of $M_t$ is clear: it represents the
best estimate, conditional on the information available at time $t$,
of the discounted value of the terminal \emph{payoff}, computed under
the risk-neutral measure $\Q$.

\medskip

\noindent\textbf{Well-definedness and square integrability.} For the
\emph{payoffs} of standard derivatives, including the European
\emph{call} $(S_T - K)^+$, the condition $h(S_T) \in L^2(\Q)$ holds
thanks to the log-normality of the underlying under $\Q$
(see~\eqref{eq:S_explicito_Q}) and the equivalence between $\PP$ and
$\Q$ (see \cite{karatzas-shreve}, Section 5.8, for the details).
Consequently, $M_t \in L^2(\Q)$ for all $t$, a necessary condition
for applying the Brownian martingale representation theorem in the
next step.

\medskip

\noindent\textbf{Martingale property.}

\begin{proposicion}[$\{M_t\}$ is a $\Q$-martingale]
\label{prop:M_es_Q_martingala}
The process $\{M_t\}$ defined in~\eqref{eq:def_M_continuo} is a
$\Q$-martingale with respect to the filtration $(\F_t)$.
\end{proposicion}

\begin{proof}
The adaptedness of $M_t$ to $\F_t$ is immediate from the very
definition of conditional expectation. Square integrability follows
from $h(S_T) \in L^2(\Q)$ and Jensen's inequality applied to
conditional expectation:
$\E_\Q[M_t^2] \leq \E_\Q\bigl[\E_\Q[(e^{-rT} h(S_T))^2 \mid \F_t]\bigr] = \E_\Q[(e^{-rT} h(S_T))^2] < \infty$.

For the martingale property, let $0 \leq s \leq t \leq T$. Applying
the tower property
(Proposition~\ref{prop:propiedades_esperanza_cond}, property (iv))
with $\F_s \subseteq \F_t$,
\begin{equation*}
  \E_\Q[M_t \mid \F_s]
  = \E_\Q\!\left[\,\E_\Q\bigl[e^{-rT}\, h(S_T) \,\big|\, \F_t\bigr]
                 \,\Big|\, \F_s\right]
  = \E_\Q\!\bigl[e^{-rT}\, h(S_T) \,\big|\, \F_s\bigr]
  = M_s.
\end{equation*}
\end{proof}

\subsection{Step 3: representation, replication and \emph{self-financing}}
\label{subsec:paso3_replicacion}

\noindent\textbf{Application of the representation theorem.} The
$\Q$-martingale $\{M_t\}$ is adapted to the filtration of the
Brownian motion $\widetilde W$ since, by~\eqref{eq:W_tilde_BS}, the
filtration $(\F_t)$ generated by $W$ and augmented with the null sets
is the same as the one generated by $\widetilde W$; moreover, it
satisfies $M_t \in L^2(\Q)$ for all $t$ by
Proposition~\ref{prop:M_es_Q_martingala}. We are therefore in a
position to apply Theorem~\ref{teo:representacion_martingalas} under
the measure $\Q$ and with respect to the Brownian motion
$\widetilde W$: there exists a unique process
$\widetilde\phi \in \mathcal{H}^2[0, T]$ (where $\mathcal{H}^2$ is
now understood with respect to $\Q$) such that
\begin{equation}\label{eq:representacion_M}
  M_t = M_0 + \int_0^t \widetilde\phi_s\, d\widetilde W_s,
  \qquad t \in [0, T].
\end{equation}

\medskip

\noindent\textbf{Identification of the portfolio.} The integrand
$\widetilde\phi_t$ provided by~\eqref{eq:representacion_M} is the
piece that allows us to construct the replicating strategy. The idea
is the following: we want to define $(\phi_t, \psi_t)$ so that the
discounted value of the portfolio, $\widetilde V_t \eqdef e^{-rt}\, V_t$,
coincides with $M_t$ at every time. To this end, observe that if
such an identity holds, then by~\eqref{eq:representacion_M} the
differential of $\widetilde V_t$ would be
$d\widetilde V_t = \widetilde\phi_t\, d\widetilde W_t$. On the other
hand, if the strategy is \emph{self-financing} and
$\widetilde V_t = e^{-rt} V_t$, a direct computation (which we shall
make precise below) gives
\begin{equation*}
  d\widetilde V_t = \phi_t\, d\widetilde S_t
                  = \phi_t\, \sigma\, \widetilde S_t\, d\widetilde W_t,
\end{equation*}
where the last equality uses~\eqref{eq:S_descontado_Q_martingala}.
Comparing both expressions for $d\widetilde V_t$, we identify
$\widetilde\phi_t = \phi_t\, \sigma\, \widetilde S_t$, that is,
\begin{equation}\label{eq:phi_BS}
  \boxed{\;
    \phi_t \eqdef \frac{\widetilde\phi_t}{\sigma\, \widetilde S_t}
              = \frac{\widetilde\phi_t\, e^{rt}}{\sigma\, S_t}. \;}
\end{equation}
The bond holding $\psi_t$ is then defined by forcing the value of the
portfolio to equal the discounted value given by the representation,
\begin{equation}\label{eq:psi_BS}
  \boxed{\;
    \psi_t \eqdef \frac{V_t - \phi_t\, S_t}{B_t},
    \qquad \text{with } V_t \eqdef e^{rt}\, M_t. \;}
\end{equation}
Thus $V_t = \phi_t\, S_t + \psi_t\, B_t$ by construction and, in
particular,
$V_T = e^{rT}\, M_T = e^{rT}\, e^{-rT}\, h(S_T) = h(S_T)$: the
strategy would therefore satisfy the replication condition
(Definition~\ref{def:estrategia_replicante}). Consequently, it
suffices to verify that $(\phi_t, \psi_t)$ is \emph{self-financing}
in the sense of Definition~\ref{def:self_financing_continuo}.

\medskip

\noindent\textbf{Verification of the \emph{self-financing} condition.}

\begin{teorema}[The strategy is \emph{self-financing}]
\label{teo:replicante_self_financing}
The strategy $(\phi_t, \psi_t)$ defined by~\eqref{eq:phi_BS}
and~\eqref{eq:psi_BS} is \emph{self-financing}, that is, it satisfies
\begin{equation*}
  dV_t = \phi_t\, dS_t + \psi_t\, dB_t.
\end{equation*}
\end{teorema}

\begin{proof}
The proof consists in computing both sides of the identity and
checking that they coincide.

\emph{Step 1: computation of $dV_t$.} By construction,
$V_t = e^{rt}\, M_t$. We apply Itô's formula
(Theorem~\ref{teo:formula_ito_general}) to the function
$g(t, m) = e^{rt}\, m$ and the process $X_t = M_t$. The derivatives
are $\partial_t g = r\, e^{rt}\, m$, $\partial_m g = e^{rt}$ and
$\partial_{mm} g = 0$. Since $\{M_t\}$ has
differential~\eqref{eq:representacion_M}, with zero drift and
diffusion coefficient $\widetilde\phi_t$ with respect to
$\widetilde W$, the formula gives
\begin{equation}\label{eq:dV_calculado}
  dV_t = r\, e^{rt}\, M_t\, dt + e^{rt}\, \widetilde\phi_t\, d\widetilde W_t
       = r\, V_t\, dt + e^{rt}\, \widetilde\phi_t\, d\widetilde W_t.
\end{equation}

\emph{Step 2: computation of $\phi_t\, dS_t + \psi_t\, dB_t$.} For
$\phi_t\, dS_t$, we substitute~\eqref{eq:S_dinamica_Q} into the
definition~\eqref{eq:phi_BS}:
\begin{align*}
  \phi_t\, dS_t
  &= \frac{\widetilde\phi_t\, e^{rt}}{\sigma\, S_t}\,
     \bigl(r\, S_t\, dt + \sigma\, S_t\, d\widetilde W_t\bigr) \\
  &= \frac{r\, \widetilde\phi_t\, e^{rt}}{\sigma}\, dt
   + \widetilde\phi_t\, e^{rt}\, d\widetilde W_t.
\end{align*}
For $\psi_t\, dB_t$, using $dB_t = r\, B_t\, dt$ and the
definition~\eqref{eq:psi_BS} of $\psi_t$:
\begin{equation*}
  \psi_t\, dB_t = r\, \psi_t\, B_t\, dt
              = r\, (V_t - \phi_t\, S_t)\, dt.
\end{equation*}
Adding both contributions,
\begin{equation*}
  \phi_t\, dS_t + \psi_t\, dB_t
  = \frac{r\, \widetilde\phi_t\, e^{rt}}{\sigma}\, dt
   + \widetilde\phi_t\, e^{rt}\, d\widetilde W_t
   + r\, V_t\, dt - r\, \phi_t\, S_t\, dt.
\end{equation*}
Substituting $\phi_t\, S_t = \widetilde\phi_t\, e^{rt} / \sigma$
(solving from~\eqref{eq:phi_BS}), the terms
$\frac{r\, \widetilde\phi_t\, e^{rt}}{\sigma}\, dt$ and
$-r\, \phi_t\, S_t\, dt$ cancel, leaving
\begin{equation}\label{eq:dV_self_financing}
  \phi_t\, dS_t + \psi_t\, dB_t
  = r\, V_t\, dt + \widetilde\phi_t\, e^{rt}\, d\widetilde W_t.
\end{equation}

\emph{Step 3: comparison.} The
expressions~\eqref{eq:dV_calculado} and~\eqref{eq:dV_self_financing}
coincide term by term. The strategy $(\phi_t, \psi_t)$ is therefore
\emph{self-financing}.
\end{proof}

\medskip

\noindent\textbf{The central result.} Combining the three steps, we
are in a position to state the main result of the section.

\begin{teorema}[Pricing of European derivatives in the Black-Scholes model]
\label{teo:valoracion_BS}
Let $h(S_T)$ be a European \emph{payoff} satisfying the integrability
condition~\eqref{eq:integrabilidad_payoff}. Then there exists a
unique \emph{self-financing} replicating strategy $(\phi_t, \psi_t)$
for the derivative, with portfolio value at each time $t \in [0, T]$
given by
\begin{equation}\label{eq:formula_valoracion}
  \boxed{\;
    V(t, S_t) = \E_\Q\!\bigl[e^{-r(T - t)}\, h(S_T) \,\big|\, \F_t\bigr]. \;}
\end{equation}
By the no-arbitrage argument, $V(t, S_t)$ is the fair price of the
derivative at time $t$.
\end{teorema}

\begin{proof}
Existence follows from the construction carried out in Steps 1, 2
and 3, together with the \emph{self-financing} verification of
Theorem~\ref{teo:replicante_self_financing}.
Formula~\eqref{eq:formula_valoracion} is obtained from
$V_t = e^{rt}\, M_t = e^{rt}\, \E_\Q[e^{-rT}\, h(S_T) \mid \F_t]
= \E_\Q[e^{-r(T-t)}\, h(S_T) \mid \F_t]$, using the
definition~\eqref{eq:def_M_continuo} of $M_t$ and the
$\F_t$-measurability of the factor $e^{rt}$.

Uniqueness follows from the uniqueness of the integrand
$\widetilde\phi$ in the representation theorem
(Theorem~\ref{teo:representacion_martingalas}): if two
\emph{self-financing} replicating strategies produced the same
terminal value $h(S_T)$, their respective martingales
$\E_\Q[e^{-rT}\, h(S_T) \mid \F_t]$ would coincide, their
representation would be unique, and by
\eqref{eq:phi_BS}–\eqref{eq:psi_BS} the two strategies would
necessarily be equal.
\end{proof}

\medskip

\noindent\textbf{Reading of the result.}
Formula~\eqref{eq:formula_valoracion} is the continuous-time
translation of Theorem~\ref{teo:valoracion_binomial}: the price at
each time is the conditional expectation, under the risk-neutral
measure, of the discounted \emph{payoff}.

\medskip

We have thus solved the general pricing problem. Formula
\eqref{eq:formula_valoracion} expresses, for any European
\emph{payoff} $h$, the price of the derivative as a conditional
expectation under $\Q$. The next section specializes this formula to
the case of the European \emph{call} option, $h(S_T) = (S_T - K)^+$,
and obtains by direct computation the classical Black-Scholes
formula.
\section{The Black-Scholes formula for the European \emph{call} option}
\label{sec:formula_bs}

We now specialize the formula of Theorem~\ref{teo:valoracion_BS} to
the \emph{payoff} of the European \emph{call}, $h(S_T) = (S_T - K)^+$,
thereby obtaining the classical Black-Scholes formula (1973) with its
coefficients $d_1$ and $d_2$.

\medskip

\noindent\textbf{The European \emph{call} option.} Let us briefly
recall the contract. The European \emph{call} option with
\emph{strike} $K > 0$ and maturity $T$ is the European derivative
whose \emph{payoff} at $T$ is
\begin{equation}\label{eq:payoff_call}
  h(S_T) = (S_T - K)^+ \eqdef \max(S_T - K, 0).
\end{equation}
Its holder acquires the right (but not the obligation) to buy the
underlying at the price $K$ at $T$: if at that time $S_T > K$ they
exercise the right and obtain a profit $S_T - K$; if $S_T \leq K$
they do not exercise it and the \emph{payoff} is zero. The
integrability condition~\eqref{eq:integrabilidad_payoff} holds
automatically for this \emph{payoff} since
$0 \leq (S_T - K)^+ \leq S_T$ and the moments of the GBM are finite
of all orders, so that Theorem~\ref{teo:valoracion_BS} applies.

\medskip

\noindent\textbf{Strategy of the computation.} We shall first
compute the price of the \emph{call} at $t = 0$, that is,
$V(0, S_0) = \E_\Q[e^{-rT}\, (S_T - K)^+]$, and at the end obtain the
generalization to the case $t > 0$ by an argument based on the Markov
property of $\{S_t\}$. For $t = 0$, the strategy consists in
decomposing the expectation by means of the exercise indicator
function:
\begin{equation}\label{eq:descomposicion_call}
  \E_\Q\bigl[(S_T - K)^+\bigr]
  = \E_\Q\bigl[(S_T - K)\, \mathbf{1}_{\{S_T > K\}}\bigr]
  = \underbrace{\E_\Q\bigl[S_T\, \mathbf{1}_{\{S_T > K\}}\bigr]}_{(\mathrm{I})}
   - K\, \underbrace{\Q(S_T > K)}_{(\mathrm{II})}.
\end{equation}

\noindent\textbf{Expression of $S_T$ under $\Q$.}
By~\eqref{eq:S_explicito_Q} with $t = T$,
\begin{equation}\label{eq:S_T_Q}
  S_T = S_0\, \exp\!\left(\Bigl(r - \tfrac{1}{2}\sigma^2\Bigr) T
                          + \sigma\, \widetilde W_T\right),
\end{equation}
where $\widetilde W_T \sim \mathcal{N}(0, T)$ under $\Q$. Rewriting
$\widetilde W_T = \sqrt{T}\, Z$ with $Z \sim \mathcal{N}(0, 1)$ under
$\Q$,
\begin{equation*}
  S_T = S_0\, \exp\!\left(\Bigl(r - \tfrac{1}{2}\sigma^2\Bigr) T
                          + \sigma\, \sqrt{T}\, Z\right).
\end{equation*}
The inequality $S_T > K$ is equivalent, taking logarithms, to
\begin{equation*}
  \log\!\left(\frac{S_0}{K}\right)
  + \Bigl(r - \tfrac{1}{2}\sigma^2\Bigr) T
  + \sigma\sqrt{T}\, Z > 0,
\end{equation*}
that is, $Z > -d_2$, where
\begin{equation}\label{eq:def_d2}
  d_2 \eqdef \frac{\log(S_0/K) + \bigl(r - \tfrac{1}{2}\sigma^2\bigr) T}
                  {\sigma\, \sqrt{T}}.
\end{equation}

\noindent\textbf{Computation of the term $(\mathrm{II})$.} By the
symmetry of the Gaussian distribution,
$(\mathrm{II}) = \Q(Z > -d_2) = N(d_2)$, where
$N(x) \eqdef (2\pi)^{-1/2} \int_{-\infty}^x e^{-y^2/2}\, dy$ is the
standard normal distribution function.

\medskip

\noindent\textbf{Computation of the term $(\mathrm{I})$.} We have
\begin{equation*}
  (\mathrm{I}) = \E_\Q\bigl[S_T\, \mathbf{1}_{\{Z > -d_2\}}\bigr]
              = S_0\, e^{(r - \frac{1}{2}\sigma^2) T}\,
                \E_\Q\!\left[e^{\sigma\sqrt{T}\, Z}\,
                              \mathbf{1}_{\{Z > -d_2\}}\right].
\end{equation*}
To evaluate the expectation, we write the Gaussian integral
explicitly, using that $Z \sim \mathcal{N}(0, 1)$, and complete the
square in the exponent:
\begin{align*}
  \E_\Q\!\left[e^{\sigma\sqrt{T}\, Z}\, \mathbf{1}_{\{Z > -d_2\}}\right]
  &= \frac{1}{\sqrt{2\pi}} \int_{-d_2}^{\infty}
       e^{\sigma\sqrt{T}\, z}\, e^{-z^2/2}\, dz \\
  &= \frac{1}{\sqrt{2\pi}} \int_{-d_2}^{\infty}
       \exp\!\left(-\tfrac{1}{2}\bigl(z - \sigma\sqrt{T}\bigr)^2
                   + \tfrac{1}{2}\sigma^2 T\right) dz \\
  &= e^{\frac{1}{2}\sigma^2 T} \cdot
     \frac{1}{\sqrt{2\pi}} \int_{-d_2}^{\infty}
       \exp\!\left(-\tfrac{1}{2}\bigl(z - \sigma\sqrt{T}\bigr)^2\right) dz.
\end{align*}
We make the change of variable $u \eqdef z - \sigma\sqrt{T}$, with
$du = dz$ and new limits $u > -d_2 - \sigma\sqrt{T}$:
\begin{equation*}
  \E_\Q\!\left[e^{\sigma\sqrt{T}\, Z}\, \mathbf{1}_{\{Z > -d_2\}}\right]
  = e^{\frac{1}{2}\sigma^2 T}\,
    \frac{1}{\sqrt{2\pi}} \int_{-d_2 - \sigma\sqrt{T}}^{\infty}
       e^{-u^2/2}\, du
  = e^{\frac{1}{2}\sigma^2 T}\, N(d_2 + \sigma\sqrt{T}),
\end{equation*}
where the last equality again uses the symmetry of the Gaussian.
Defining
\begin{equation}\label{eq:def_d1}
  d_1 \eqdef d_2 + \sigma\sqrt{T}
       = \frac{\log(S_0/K) + \bigl(r + \tfrac{1}{2}\sigma^2\bigr) T}
              {\sigma\, \sqrt{T}},
\end{equation}
and substituting into the expression for $(\mathrm{I})$, the factor
$e^{\frac{1}{2}\sigma^2 T}$ cancels exactly with the corresponding
term in the drift $r - \tfrac{1}{2}\sigma^2$ of the exponent of
$S_T$, leaving
\begin{equation}\label{eq:termino_I}
  (\mathrm{I}) = S_0\, e^{rT}\, N(d_1).
\end{equation}

\medskip

\noindent\textbf{The Black-Scholes formula.} Substituting
$(\mathrm{I})$ and $(\mathrm{II})$ into~\eqref{eq:descomposicion_call}
and multiplying by the discount factor $e^{-rT}$, we obtain the
central formula:

\begin{teorema}[Black-Scholes formula for the European \emph{call}]
\label{teo:formula_BS}
The price at $t = 0$ of the European \emph{call} option with
\emph{strike} $K$ and maturity $T$ in the Black-Scholes model is
\begin{equation}\label{eq:formula_BS}
  \boxed{\;
    V(0, S_0) = S_0\, N(d_1) - K\, e^{-rT}\, N(d_2),
  \;}
\end{equation}
where
\begin{equation*}
  d_1 = \frac{\log(S_0/K) + \bigl(r + \tfrac{1}{2}\sigma^2\bigr) T}
              {\sigma\, \sqrt{T}},
  \qquad
  d_2 = d_1 - \sigma\sqrt{T}
      = \frac{\log(S_0/K) + \bigl(r - \tfrac{1}{2}\sigma^2\bigr) T}
              {\sigma\, \sqrt{T}},
\end{equation*}
and $N(\cdot)$ is the standard normal distribution function.
\end{teorema}

\medskip

\noindent\textbf{Generalization to the case $t > 0$.} Formula
\eqref{eq:formula_BS} has been derived at $t = 0$. The
generalization to the case $t \in (0, T)$ is obtained by taking into
account the following observation: the process $\{S_t\}$ under $\Q$
is Markovian (since it solves the SDE \eqref{eq:S_dinamica_Q}, whose
coefficients depend only on the current state) and, therefore, the
conditional expectation $\E_\Q[(S_T - K)^+ \mid \F_t]$ depends
exclusively on $S_t$, not on the whole past path. By the Markov
property, computing this conditional expectation is equivalent to
computing the unconditional expectation of the problem ``started''
at $t$ with initial value $S_t$ and duration $T - t$. Thus, the
derivation we carried out for $t = 0$, repeated with $S_0$ replaced
by $S_t$ and $T$ replaced by $T - t$, gives
\begin{equation}\label{eq:formula_BS_t}
  \boxed{\;
    V(t, S_t) = S_t\, N(d_1(t, S_t)) - K\, e^{-r(T-t)}\, N(d_2(t, S_t)),
  \;}
\end{equation}
where
\begin{align*}
  d_1(t, S_t) &= \frac{\log(S_t/K) + \bigl(r + \tfrac{1}{2}\sigma^2\bigr)(T-t)}
                       {\sigma\, \sqrt{T-t}}, \\
  d_2(t, S_t) &= d_1(t, S_t) - \sigma\sqrt{T-t}.
\end{align*}
Formula~\eqref{eq:formula_BS_t} expresses the price of the
\emph{call} at any time $t \in [0, T)$ as a function of the current
price of the underlying $S_t$ and the time to maturity $T - t$. It is
the \emph{value function} of the derivative, to which we shall return
in the next section when deriving the Black-Scholes PDE.

\section{The Black-Scholes PDE}
\label{sec:edp_bs}

We close the chapter with the derivation of the classical
Black-Scholes PDE from the martingale calculus.

\medskip

The derivation we have carried out in the preceding sections (via
change of measure, construction of the $\Q$-martingale associated
with the \emph{payoff}, and application of the representation
theorem) is the \emph{martingale route} to the Black-Scholes theory.
There is an alternative classical derivation, historically earlier,
which starts from the idea of \emph{instantaneous hedging} and leads
to a partial differential equation. This section recovers this
second formulation from formula~\eqref{eq:formula_BS_t}, showing the
consistency between the two routes.

\medskip

\noindent\textbf{Application of Itô to the value function.} The
value function of the European \emph{call}, $V(t, S_t)$, given
by~\eqref{eq:formula_BS_t}, is of class $C^{1,2}$ on the domain
$[0, T) \times (0, \infty)$. We apply Itô's formula
(Theorem~\ref{teo:formula_ito_general}) to $V(t, S_t)$, with $S_t$
an Itô process under $\Q$ with dynamics~\eqref{eq:S_dinamica_Q}:
\begin{equation}\label{eq:Ito_V}
  dV(t, S_t)
  = \biggl[\partial_t V + r\, S_t\, \partial_S V
           + \tfrac{1}{2}\, \sigma^2\, S_t^2\, \partial_{SS} V\biggr]\, dt
   + \sigma\, S_t\, \partial_S V\, d\widetilde W_t.
\end{equation}

\medskip

\noindent\textbf{Comparison with the dynamics of the replicating
portfolio.} On the other hand, in
Subsection~\ref{subsec:paso3_replicacion} we proved that the value
of the replicating portfolio satisfies (combining
\eqref{eq:dV_calculado} and the identification \eqref{eq:phi_BS}):
\begin{equation}\label{eq:dinamica_V_portfolio}
  dV_t = r\, V_t\, dt + \sigma\, S_t\, \phi_t\, d\widetilde W_t,
\end{equation}
where $\phi_t = \widetilde\phi_t\, e^{rt} / (\sigma\, S_t)$ is the
quantity of the underlying. Since $V(t, S_t) = V_t$ by construction,
expressions~\eqref{eq:Ito_V} and~\eqref{eq:dinamica_V_portfolio}
must coincide term by term.

\emph{Equating the coefficients of $d\widetilde W_t$:}
$\sigma\, S_t\, \partial_S V = \sigma\, S_t\, \phi_t$, whence
$\phi_t = \partial_S V$.

\emph{Equating the coefficients of $dt$:}
\begin{equation*}
  \partial_t V + r\, S_t\, \partial_S V
  + \tfrac{1}{2}\, \sigma^2\, S_t^2\, \partial_{SS} V
  = r\, V.
\end{equation*}
Since this identity must hold for all $(t, S_t)$ with $S_t$ ranging
over $(0, \infty)$ with positive probability, it translates into a
partial differential equation for the function $V(t, S)$:

\begin{teorema}[Black-Scholes equation]
\label{teo:edp_BS}
The value function $V(t, S)$ of the European \emph{call} in the
Black-Scholes model satisfies the partial differential equation
\begin{equation}\label{eq:EDP_BS}
  \boxed{\;
    \partial_t V + r\, S\, \partial_S V
    + \tfrac{1}{2}\, \sigma^2\, S^2\, \partial_{SS} V = r\, V,
  \;}
\end{equation}
on the domain $(0, T) \times (0, \infty)$, with the terminal
boundary condition
\begin{equation}\label{eq:condicion_terminal_BS}
  V(T, S) = (S - K)^+.
\end{equation}
\end{teorema}

Identity~\eqref{eq:EDP_BS} is the classical \emph{Black-Scholes
equation}, in its original form published in 1973. By what has been
proved, the function defined by~\eqref{eq:formula_BS_t} is a solution
of this boundary value problem; a uniqueness argument shows that it
is in fact the unique sufficiently regular solution (see
\cite{karatzas-shreve}, Section 5.8).

\medskip

\noindent\textbf{The connection via Feynman-Kac.} The equivalence
between the two derivation routes, martingales and PDE, can be
expressed abstractly through the \emph{Feynman-Kac formula}, which
establishes the identity between the solution of the
PDE~\eqref{eq:EDP_BS}–\eqref{eq:condicion_terminal_BS} and the
probabilistic formula~\eqref{eq:formula_valoracion}. Its rigorous
proof lies outside the scope of this work and can be found in
\cite{shreve}, Chapter 6.

% =====================================================================
%   CHAPTER 6 — CONCLUSIONS
% =====================================================================
\chapter{Conclusions}\label{cap:conclusiones}
\section{Summary of the work}
\label{sec:recapitulacion}

This work has focused on the problem of \emph{pricing the European
derivative}, addressed in its discrete version
(Chapter~\ref{cap:discreto}, binomial model) and in its continuous
version (Chapters~\ref{cap:brown} to~\ref{cap:bs}, Black-Scholes
model). In both settings the same three-step structure is observed:
change of measure to a risk-neutral measure, construction of the
$\Q$-martingale associated with the discounted \emph{payoff}, and
application of a representation theorem. The translation between the
two settings is not trivial (combinatorial and constructive in the
discrete case, via Girsanov and abstract arguments in the continuous
one), but the logical structure is identical.

The structural relationship between the discrete binomial model and
the Black-Scholes model admits, beyond this conceptual analogy, a
precise quantitative formulation: the Black-Scholes formula is
\emph{the limit} of the multi-period binomial formula as the number of
steps tends to infinity. If in the binomial tree on $[0, T]$ with $N$
periods of length $\Delta t = T/N$ the up and down parameters are
defined as
\begin{equation*}
  u = e^{\mu\Delta t+\sigma\sqrt{\Delta t}}, \qquad
  d = e^{\mu\Delta t-\sigma\sqrt{\Delta t}}, \qquad
  q = \frac{e^{r \Delta t} - d}{u - d},
\end{equation*}
and one lets $N \to \infty$, the binomial random walk converges in
distribution to the GBM under $\Q$ (by the central limit theorem),
and the price of the \emph{call} obtained by finite recursion on the
tree converges to the Black-Scholes formula \eqref{eq:formula_BS}.
Analogously, the quantity $\phi_k^A$ of the discrete replicating
portfolio, given by the explicit formula~\eqref{eq:phi_explicito},
converges to $N(d_1)$. The rigorous proof of these convergences can
be found in \cite{etheridge}, Chapter 2. The Black-Scholes formula is
therefore not an object independent of the binomial model, but its
limit as the time steps become infinitely small.

\section{Reflection on the approach}
\label{sec:reflexion_enfoque}

The fundamental structural decision of this work has been to approach
Black-Scholes via martingales rather than via PDEs. Both formulations
are equivalent, and their consistency is expressed through
Feynman-Kac (Section~\ref{sec:edp_bs}); however, the martingale route
allows the passage from the discrete to the continuous model to be
presented as a natural extension, preserving the conceptual ideas of
Chapter~\ref{cap:discreto}. This transition, less evident in a
PDE-based presentation, justifies the approach adopted.


It is worth stressing, finally, that the apparatus developed to
pursue this route (the change of measure, Girsanov's theorem and
martingale representation) is of interest well beyond the pricing
problem: these are central tools of modern probability and
stochastic analysis. The financial problem has been, in a sense, the
thread that motivated the introduction of that machinery.

\section{Further directions}
\label{sec:lineas_continuacion}

This work has been restricted to European options on a single
underlying in the Black-Scholes model with constant parameters. Each
of these restrictions can be relaxed in various ways:

\begin{enumerate}[leftmargin=*]
  \item \emph{American options}: incorporating the notion of optimal
        stopping time into the scheme developed here.
  \item \emph{More realistic models}: stochastic volatility (Heston),
        jumps (Merton), stochastic interest rates (Vasicek,
        Hull-White) and combinations thereof.
  \item \emph{Numerical methods}: Monte Carlo simulation and finite
        differences for \emph{payoffs} without a closed-form formula
        (Asian, \emph{barrier}, \emph{basket options}, etc.).
  \item \emph{Calibration and implied volatility}: practical
        estimation of $\sigma$ from market prices of liquid options.
\end{enumerate}

Each of them presupposes the foundations developed in the present
work and constitutes an independent chapter of contemporary
mathematical finance.

% =====================================================================
%   BIBLIOGRAPHY
% =====================================================================
\begin{thebibliography}{99}

\bibitem{bachelier}
L. Bachelier.
\emph{Théorie de la spéculation}.
Annales scientifiques de l'École Normale Supérieure, 3\textsuperscript{e}
série, 17, 21--86, 1900.

\bibitem{black-scholes}
F. Black and M. Scholes.
The pricing of options and corporate liabilities.
\emph{Journal of Political Economy}, 81 (3), 637--654, 1973.

\bibitem{etheridge}
A. Etheridge.
\emph{A Course in Financial Calculus}.
Cambridge University Press, Cambridge, 2002.

\bibitem{karatzas-shreve}
I. Karatzas and S. E. Shreve.
\emph{Brownian Motion and Stochastic Calculus}.
2nd edition, Graduate Texts in Mathematics, vol.~113.
Springer, New York, 1991.

\bibitem{merton}
R. C. Merton.
Theory of rational option pricing.
\emph{Bell Journal of Economics and Management Science}, 4 (1),
141--183, 1973.

\bibitem{oksendal}
B.~\O{}ksendal.
\emph{Stochastic Differential Equations: An Introduction with
Applications}.
6th edition, Universitext. Springer, Berlin, 2003.

\bibitem{shreve}
S. E. Shreve.
\emph{Stochastic Calculus for Finance II: Continuous-Time Models}.
Springer Finance. Springer, New York, 2004.

\bibitem{williams}
D. Williams.
\emph{Probability with Martingales}.
Cambridge Mathematical Textbooks. Cambridge University Press,
Cambridge, 1991.

\end{thebibliography}

\end{document}

